\chapter{Hecke algebras, congruence ideals and the theorem of Diamond, Flach and Guo}\label{ch:congruence}
\chaptermark{Congruence ideals and Diamond-Flach-Guo}

This chapter is where Section 4 of the paper is proved. It joins three threads that earlier chapters have prepared separately. The first is
the commutative algebra of congruences: a surjection $\pi\colon\T\onto\mathcal O$ from a finite free algebra has a congruence ideal $\eta(\pi)$ that measures how
far $\T$ is from splitting off the factor $\mathcal O$, and a lemma of Darmon, Diamond and Taylor compares $\eta(\pi)$ with the Gram determinant of a pairing on a
module (\Cref{ch:modules} proves every ingredient; the lemma is proved here). The second is the arithmetic of modular forms: the Hecke algebra $\T$ of
level $N$, the congruence number $r_E$, the depleted forms $f^\Sigma$ and the twist map $\iota$ (\Cref{ch:hecke,ch:moddeg,ch:modgalrep}). The third is the
theorem of Diamond, Flach and Guo \cite{DFG04}: the length of an imprimitive adjoint Selmer group is the valuation of a naive congruence ideal, and that
ideal changes with the level by the naive Euler factors of the adjoint $L$-function. With the Selmer bookkeeping of \Cref{ch:blochkato} and the Euler
factors of \Cref{ch:lfunctions}, these three threads give the paper's Theorem thm:main, as an equality, and Theorems thm:twist and thm:quant.

The chapter assembles; it does not re-derive what other chapters prove. Its own full proofs are: Lemma lem:ddt417 (\Cref{lem:congruence:ddt417}); the assembly of Theorem thm:main
(\Cref{thm:congruence:main}); Theorem thm:twist (\Cref{thm:congruence:twist}); Theorem thm:quant with its equality clause (\Cref{thm:congruence:quant}); and
the equality $v_\ell(r_E)=v_\ell(\eta^{\mathrm{Hecke}}(g^S))$ in Step 1 of the proof of Theorem thm:quant, which the paper's Remark rem:afterquant (2) calls open
(\Cref{lem:congruence:quant-step1}). The deep theorems (the Taylor-Wiles method, the theorems of Diamond, Flach and Guo, Hida's theorem, the formula of
Kim and Ota) are stated precisely, with an outline of their proofs and the references; for part (c) of the proposition of Diamond, Flach and Guo we sketch the
deduction from parts (a) and (b) along the lines of their published proof, given the integral comparison isomorphism of their integral structures, which we cite (\Cref{thm:congruence:dfg14}).

\section*{Core results at a glance}

Throughout, $f$ is a newform of weight $2$ on $\Gamma_0(M)$ (from Section \ref{sec:congruence:assembly} on, $f=f_E$ with $E$ optimal of conductor $N$), $\T$ the full Hecke algebra ($U_p$ included), $\ell$ an odd
prime, $\pi\colon\T_{\mathcal O}\onto\mathcal O$ a surjection with $\wp=\ker\pi$, $I=\Ann(\wp)$, $\eta(\pi)=\pi(I)$, and $\Frob_p^{-1}$ the geometric Frobenius.

\begin{enumerate}[label=\textbf{\thechapter.\arabic*}]
\item \textbf{Conventions and warnings.} Notation C(a) to C(h); the paper's $I,J$ of Lemma lem:crit are our $\wp,I$; the Hecke ideal $\eta^{\mathrm{Hecke}}$ and the cohomological ideal $\eta^\Sigma_{f,\lambda}$ are
different objects. \emph{(convention)} \Cref{conv:congruence:setup}, \Cref{warn:congruence:letters}, \Cref{warn:congruence:two-ideals}.
\item \textbf{Local components.} $\T_\ell=\prod_{\mf m}\T_{\mf m}$; $\mf m_{f,\ell}=(\ell,T_n-a_n(f))$; $e_{\mf m}S_2\otimes\overline\Q_\ell=\bigoplus_{\mf m_\phi=\mf m}V(\phi)$, and forms of $V_h$ in $\mf m_{f,\ell}$ have $\bar\rho_h^{\mathrm{ss}}\simeq\bar\rho_f^{\mathrm{ss}}$; the pairings
$S_2(\Zl)\times\T_\ell\to\Zl$ and $\hat\delta_!$ are perfect on each component, the latter orthogonally. \emph{(proved here by citation of \Cref{ch:modules,ch:hecke,ch:moddeg})} \Cref{prop:congruence:hecke-local}.
\item \textbf{Level $539$.} At $\ell=3$ the newforms $539$c, $539$d share one maximal ideal and the depleted $(77\mathrm c)^{\{7\}}$, $(11\mathrm a)^{\{7\}}$ another; the congruence modules have orders $288$, $504$, $288$, $504$.
\emph{(computed)} \Cref{ex:congruence:level-539}.
\item \textbf{Congruence ideal, module, cotangent module} (definition): $\eta(\pi)=\pi(I)$, $\mathcal C(\pi)=\T/(\wp+I)$, $\mathrm{Cot}(\pi)=\wp/\wp^2$. \Cref{def:congruence:congruence-module}.
\item \textbf{Properties.} $\mathcal C(\pi)\simeq\mathcal O/\eta$; $\eta\ne0\iff\wp\cap I=0\iff$ the factor of $\T_K$ at $\pi$ is $K$; then $I\simeq\eta$, $\Ann(I)=\wp$, $\eta$ is the denominator of the idempotent; $\eta$ is computed in the
local factor of $\pi$ (the other factors lie in $\wp$); $\mathrm{Cot}(\pi)$ is finite iff the factor is $K$. \emph{(proved here, given \Cref{ch:modules})} \Cref{prop:congruence:congruence-module}.
\item \textbf{Toy algebras.} $\T_r$: $\len\mathrm{Cot}=v(\eta)=r$; $\T_3^\flat$: $2>1$; $\mathcal O[x]/(x^2)$: $\eta=0$. \emph{(proved here)} \Cref{ex:congruence:toy}.
\item \textbf{Hecke congruence ideal} (definition): $\eta^{\mathrm{Hecke}}_f=\eta(\phi_f)$, and $\eta^{\mathrm{Hecke}}(g)$ for an eigenform $g$. It is nonzero whenever the generalised eigenspace is a line (newforms, depleted forms).
\emph{(proved here, given \Cref{ch:hecke})} \Cref{def:congruence:newform-ideal}, \Cref{prop:congruence:newform-ideal}.
\item \textbf{Lemma lem:crit as congruence modules.} $\T/(\wp+I)\simeq\Z/r_E$, $\eta^{\mathrm{Hecke}}_f=r_E\Z$, the congruence moduli are the divisors of $r_E$, also over $\Zl$; and
$S_2(\Z)/(\Z f\oplus L)\simeq\Hom(\mathcal C(\phi_f),\Q/\Z)$. \emph{(stated; proved in \Cref{ch:modgalrep}; the duality proved here)} \Cref{thm:congruence:crit}.
\item \textbf{Line eigenspaces.} For $\pi\colon\T_N\to\Z$ with generalised eigenspace $\C g$ and $\T$-stable complement $g^c$: $\eta(\pi)=r(\pi)\Z$ with $r(\pi)=[S_2(\Z):\Z g\oplus(g^c\cap S_2(\Z))]$ the largest congruence modulus;
the Petersson complement does not work in general. \emph{(proved in \Cref{ch:hecke}; recalled)} \Cref{prop:congruence:crit-general}, \Cref{warn:congruence:complement}.
\item \textbf{Hida.} Outside an explicit finite set, $\lambda$ divides $\langle f,f\rangle/\Omega^+_f\Omega^-_f$ iff $\lambda$ is a congruence prime. \emph{(stated; reference)} \Cref{thm:congruence:hida}, \Cref{rem:congruence:hida-to-dfg}.
\item \textbf{Three ideals.} For optimal $E$, $\ell\nmid2N$, $E[\ell]$ irreducible: $\eta^{\mathrm{Hecke}}_f\Zl=r_E\Zl$, $v_\ell(r_E)=v_\ell(\degphi)=v_\ell(\eta^\emptyset_{f,\ell})$; at \curve{99a1}, $\ell=3$: $1\ne0$. \emph{(proved here, given the cited theorems)}
\Cref{rem:congruence:two-ideals}.
\item \textbf{Wiles's numerical criterion.} For local $\T$ finite free over $\mathcal O$ and $\eta\ne0$: $\len\mathrm{Cot}(\pi)\ge v(\eta)$, with equality iff $\T$ is a complete intersection; $R\onto\T$ with $\len\mathrm{Cot}(\pi_R)\le v(\eta)$ is an
isomorphism. \emph{(proof sketched here; Wiles, Lenstra)} \Cref{thm:congruence:numerical-criterion}.
\item \textbf{Lemma lem:ddt417.} $\T$ finite free over a discrete valuation ring $\mathcal O$, acting self-adjointly on $L$ free with a perfect pairing; $L[\wp]$ of rank $d$ with Gram determinant $\Delta\ne0$: then
$d\,v(\eta)\ge v(\Delta)$, with equality if $L$ is free over $\T$. \emph{(proved here, given \Cref{ch:modules})} \Cref{lem:congruence:ddt417}; Commentary on the paper's statement
\Cref{rem:congruence:ddt417-commentary}; hypotheses \Cref{rem:congruence:ddt417-hypotheses}; test cases \Cref{ex:congruence:ddt417-checks}.
\item \textbf{Depletion} (definition, recalled): $\delta_p$, $M^\Sigma=M\prod_{p\in\Sigma}p^{\delta_p}$, $f^\Sigma$, the naive factor $L^{\nv}_p(A_f,s)$. \Cref{def:congruence:depletion-recall}.
\item \textbf{Integral structure} (definition, after Diamond, Flach and Guo) and its comparison: $M(M,1)_{!,\lambda,\mf m}=H^1(X_0(M),\Zl)_{\mf m}$ for non-Eisenstein $\mf m$. \emph{(proof sketched here)}
\Cref{def:congruence:integral-structure}, \Cref{prop:congruence:integral-structure}.
\item \textbf{The pairing $\hat\delta_!$} (definition) and its properties: on pulled-back classes it is $\frac12\langle x,w_M^*y\rangle$; perfect, alternating, $\T$ self-adjoint, perfect on components. \emph{(proved here, given
\Cref{ch:moddeg}; the normalisation of Diamond, Flach and Guo as the paper describes it)} \Cref{def:congruence:dfg-pairing}, \Cref{prop:congruence:dfg-pairing}.
\item \textbf{$M_{f^\Sigma,\lambda}$ and $\eta^\Sigma_{f,\lambda}$} (definition): the kernel lattice of $\ker\theta_{f^\Sigma}$ and the ideal of values of $\hat\delta_!$ on it. For $\bar\rho_f$ irreducible: saturated of rank $2$, in the
$\mf m$-component, Gram determinant $(\eta^\Sigma_{f,\lambda})^2$; for $\Sigma=\emptyset$ and optimal $E$, $M_{f,\lambda}=\varphi_E^*H^1(E,\Zl)\simeq T_\ell E(-1)$. \emph{(proved here, given \Cref{ch:moddeg} and
\Cref{prop:congruence:integral-structure})} \Cref{def:congruence:eta-sigma}, \Cref{prop:congruence:eta-sigma}.
\item \textbf{Lemma lem:hecke.} For optimal $E$, $\ell$ odd, $\ell\nmid N$, $E[\ell]$ irreducible: $v_\ell(\degphi)=v_\ell(r_E)=v_\lambda(\eta^\emptyset_{f,\lambda})$, with no freeness hypothesis. \emph{(proved here, given
\cite[Theorem 2.2]{ARS12}, \Cref{ch:moddeg} and \Cref{prop:congruence:integral-structure}; the definitions of \cite{DFG04} it uses are checked against the published text)}
\Cref{thm:congruence:lem-hecke}; values \Cref{ex:congruence:eta-values}.
\item \textbf{Remark rem:free.} Multiplicity one gives $H^1(X_0(N),\Zl)_{\mf m}\simeq\T_{\mf m}^2$; no longer needed for Theorems thm:main and thm:quant. \emph{(stated; proved in \Cref{ch:moddeg} given multiplicity one)}
\Cref{rem:congruence:rem-free}.
\item \textbf{The map $\gamma$} (definition): $\gamma=\sum_m\epsilon_m\phi_m$, $\gamma(f)=f^\Sigma$. \Cref{def:congruence:gamma}.
\item \textbf{Theorem thm:dfg14} (Diamond, Flach and Guo, Proposition 1.4). For $\ell\nmid2M$, $\ell\notin\Sigma$: $\gamma^t\gamma=u\prod_pL^{\nv}_p(A_f,1)^{-1}$ on $M_f$; if $\bar\rho_f$ is irreducible,
$\eta^\Sigma_{f,\lambda}=\eta^\emptyset_{f,\lambda}\prod_{p\in\Sigma}L^{\nv}_p(A_f,1)^{-1}$; false for $\ell\in\Sigma$ (\curve{11a1}, $\ell=3$). \emph{(proof sketched here, (c) given (a), (b) and the cited integral comparison isomorphism, as in \cite[Proposition 1.4]{DFG04};
reference)} \Cref{thm:congruence:dfg14}, \Cref{rem:congruence:dfg14-status}.
\item \textbf{The three cases.} $L^{\nv}_p(A_f,1)^{-1}$ is $1$ (additive), $\frac{p-1}p\cdot\frac{(p+1)^2-a_p^2}{p^2}$ (good; Wiles's factor), $1-p^{-2}$ (multiplicative). \emph{(proved here)} \Cref{prop:congruence:three-cases};
interpretation \Cref{rem:congruence:factors-count}; Ihara and the Eisenstein case \Cref{rem:congruence:ihara}; refined structures \Cref{rem:congruence:refined-structures}.
\item \textbf{The exceptional set} (definition): $\lambda\mid Nk!$ or $\bar\rho|_{G_{\Q(\sqrt{\ell^*})}}$ not absolutely irreducible. For $f_E$ and $\ell\nmid2N$: $\ell\notin S_f$ iff $E[\ell]$ irreducible and not ($\ell=3$, image $\Ns$), iff
moreover $h_\ell(E)=0$. \emph{(definition checked against \cite{DFG04}; the proposition proved in \Cref{ch:galrep})} \Cref{def:congruence:exceptional-set}, \Cref{prop:congruence:exceptional-set}.
\item \textbf{$H^1_\Sigma(\Q,W)$} (definition): unramified outside $\Sigma\cup\set\ell$, Bloch-Kato at $\ell$, free on $\Sigma$; equal to the Bloch-Kato conditions outside $\Sigma$ when $\Sigma$ contains the ramified
primes. \Cref{def:congruence:selmer-sigma}.
\item \textbf{Minimal ramification} (definition); the primes dividing $N$ form an admissible $\Sigma$. \Cref{def:congruence:minimally-ramified}.
\item \textbf{Theorem thm:dfg37} (Diamond, Flach and Guo, Theorem 3.7): for $\lambda\notin S_f$ and $\Sigma\not\ni\ell$ with $M_{f,\lambda}$ minimally ramified outside $\Sigma$, $\len H^1_\Sigma(\Q,A_{f,\lambda}\otimes K_\lambda/\mathcal O_\lambda)=v_\lambda(\eta^\Sigma_{f,\lambda})$.
\emph{(stated; reference)} \Cref{thm:congruence:dfg37}.
\item \textbf{Theorem 0.2 of Diamond, Flach and Guo} (quoted): the same for $\Sigma\supseteq\set{p\mid N}$, $\lambda\nmid\Sigma$. \emph{(stated; checked against \cite{DFG04})} \Cref{thm:congruence:dfg02}.
\item \textbf{Theorem thm:dfg27.} For $\ell\notin S_f$: $H^0(\Q,V^D)=H^1_f(\Q,V^D)=0=H^0(\Q,A_f)=H^1_f(\Q,A_f)$. \emph{(stated; reference)} \Cref{thm:congruence:dfg27}.
\item \textbf{Theorem 0.1 of Diamond, Flach and Guo} (quoted): the $\lambda$-part of the Bloch-Kato conjecture for $A_f$ and $A_f(1)$, $\lambda\notin S_f$. \emph{(stated; checked against \cite{DFG04})} \Cref{thm:congruence:dfg01}.
\item \textbf{$\ell=3$.} As published, the theorems cover $\ell=3$ in weight $2$ ($S_f$ excludes only $\lambda\mid2N$ and the induced case); the proof treats $\ell=3$ explicitly (Fontaine-Laffaille theory for $M_{f,\lambda}$, Section 1.8, Proposition 2.16). \emph{(checked
against \cite{DFG04})} \Cref{rem:congruence:ell-three}.
\item \textbf{Deformation rings} (definition) and Mazur's theorem: the ring $R_\Sigma$ exists, and the Pontryagin dual of its cotangent module at a point $\rho$ is the Selmer group $H^1_\Sigma(\Q,\ad M_\rho\otimes K/\mathcal O)$. \emph{(proof sketched here; \cite{DDT})} \Cref{def:congruence:deformation-ring},
\Cref{thm:congruence:deformation-ring}.
\item \textbf{The Hecke side} (definition): $\T_\Sigma$ and $R_\Sigma\onto\T_\Sigma$. \Cref{def:congruence:hecke-side}.
\item \textbf{Taylor-Wiles.} For $\rho_0$ absolutely irreducible over $\Q(\sqrt{\ell^*})$, twist-minimal, Serre weight in $[2,\ell-1]$: $R_\Sigma\simeq\T_\Sigma$ complete intersections, cohomology free, and $\len H^1_\Sigma=v(\eta^\Sigma)$.
\emph{(proof sketched here; references)} \Cref{thm:congruence:tw-axiomatic}; hypotheses \Cref{rem:congruence:where-hypothesis}; why length equals valuation \Cref{rem:congruence:why-length-equals-eta}.
\item \textbf{Theorem thm:main.} For optimal $E$, $\ell$ odd, $\ell\nmid N$, $\ell\notin S_f$:
\begin{gather*}
v_\ell(\degphi)=\len H^1_f(\Q,W)+\sum_{q\mid N}v_\ell(\Tam_q(\ad T))+\sum_{q\mid N}e_\ell(E,q),\\
v_\ell(\degphi)\ge\sum_{q\mid N}e_\ell(E,q)+\sum_{q\mid N}t_\ell(E,q),
\end{gather*}
an equality with no multiplicity-one hypothesis. \emph{(proved here, given the theorems of Diamond, Flach and Guo and \Cref{ch:blochkato,ch:lfunctions,ch:moddeg})} \Cref{thm:congruence:main}; hypotheses
\Cref{rem:congruence:hypothesis-table}.
\item \textbf{Corollary cor:explicit.} $v_\ell(\degphi)=\len H^1_f(\Q,W)+\sum_qe_\ell(E,q)+\sum_{v_q(j)<0}\bigl(t_\ell+\min(t_\ell,v_\ell(q-1),v_\ell(c(u_q)))\bigr)$. \emph{(proved here, given \Cref{thm:congruence:main} and
\Cref{ch:blochkato})} \Cref{cor:congruence:explicit}; examples (\curve{11a1}, \curve{121a1}, \curve{539c1} at $3$, all with $H^1_f=0$) \Cref{ex:congruence:main-examples}.
\item \textbf{Remark rem:ko(1).} The Tamagawa excess $\min(v_\ell(n),v_\ell(q-1),v_\ell(c(u)))$ raises the bound by exactly that amount when $\ell\notin S_f$. \emph{(proved here, given \Cref{ch:blochkato})}
\Cref{rem:congruence:rem-ko-1}.
\item \textbf{Remark rem:ko(2).} Kim and Ota's formula gives $v_\ell(r_E)\ge\sum_{q\mid N^-}t_\ell(E,q)$, contained in (eq:thm), and with \Cref{cor:congruence:explicit} a Selmer formula for the $N^-$-new congruence ideal.
\emph{(proved here, given \cite{KimOta})} \Cref{prop:congruence:kim-ota}, \Cref{rem:congruence:kim-ota-identification}.
\item \textbf{The twist map.} $\iota_S$ preserves integrality, twists eigensystems by $\chi_S$, maps $g^\perp$ into $f^\perp$ when $\iota(g)=f$; levels $N_S=N/D$, $N/D^2$, $N$ for types $\kod{I}_n^*$, $\kod{I}_0^*$, potentially good.
\emph{(proved in \Cref{ch:hecke,ch:modgalrep}; recalled)} \Cref{prop:congruence:twist-inputs}.
\item \textbf{Theorem thm:twist.} If $q^2\mid N$ ($q\in S$) and $\lcm(N_S,D^2)\mid N$, then $r_{E_S}\mid r_E$; if also $\lcm(N,D^2)\mid N_S$, then $r_{E_S}=r_E$. \emph{(proved here, given \Cref{ch:hecke})}
\Cref{thm:congruence:twist}; cases and data ($25{,}477$ and $14{,}880$ pairs below conductor $10{,}000$) \Cref{ex:congruence:twist-cases}.
\item \textbf{Theorem thm:quant.} For optimal $E$, $\ell$ odd, $\ell\nmid N$, $E[\ell]$ irreducible, $S$ the odd primes of type $\kod{I}_0^*$, $\kod{I}_n^*$ ($q=3$ allowed): $v_\ell(\degphi)\ge v_\ell(r_{E_S})+\sum_{q\in S}e_\ell(E,q)$, with
equality if $\ell\notin S_f$. \emph{(proved here, given \Cref{thm:congruence:dfg14}, \Cref{thm:congruence:main} and the cited chapters)} \Cref{thm:congruence:quant}.
\item \textbf{The set-up.} $g^S$ lives at level $N$ with generalised eigenspace $\C g^S$, and $\iota(g^S)=f$. \emph{(proved here, given \Cref{ch:hecke,ch:reduction})} \Cref{lem:congruence:quant-setup}.
\item \textbf{Step 1 is an equality.} For every $E$ of conductor $N$, with $S$ and $g^S$ as in the set-up: $r_E\mid r_\pi$ and $v_\ell(r_E)=v_\ell(\eta^{\mathrm{Hecke}}(g^S))$ for every prime $\ell\notin S$, $\ell=2$ included, with no other hypothesis. \emph{(proved here, given \Cref{ch:hecke})}
\Cref{lem:congruence:quant-step1}; Commentary on Remark rem:afterquant (2) and Open question 6 \Cref{rem:congruence:step1-commentary}.
\item \textbf{Steps 2 to 4.} $v_\ell(\eta^{\mathrm{Hecke}}(g^S))\ge v_\ell(\eta^S_{g,\ell})=v_\ell(\eta^\emptyset_{g,\ell})+\sum_{q\in S}e_\ell(E,q)$ and $v_\ell(\eta^\emptyset_{g,\ell})=v_\ell(r_{E_S})$. \emph{(proved here, given \Cref{thm:congruence:dfg14})}
\Cref{lem:congruence:quant-steps234}.
\item \textbf{All equalities.} For $\ell\notin S_f$: $v_\ell(r_E)=v_\ell(\eta^{\mathrm{Hecke}}(g^S))=v_\ell(\eta^S_{g,\ell})=v_\ell(r_{E_S})+\sum_{q\in S}e_\ell(E,q)$. \emph{(proved here)} \Cref{cor:congruence:quant-equality}; remarks
\Cref{rem:congruence:outer-inequality}, \Cref{rem:congruence:after-quant}.
\item \textbf{Beyond the hypothesis.} For image $\Ns$ at $3$, (eq:identity) with $-h_3(E)$ would follow from $H^1_f(\Q,V^D)=0$ and $v_3(\eta^\Sigma)=\len H^1_\Sigma$; what Open questions 1 and 2 need. \emph{(conditional)}
\Cref{rem:congruence:rem-ns}, \Cref{rem:congruence:open-question-2}, \Cref{rem:congruence:open-question-1}.
\item \textbf{The citations of Diamond, Flach and Guo,} with their status. \Cref{rem:congruence:imported-table}.
\end{enumerate}

\section*{Where the paper uses this chapter}

\begin{itemize}
\item \textbf{Section 1} (``The congruence number is the congruence ideal of $f$ in the full Hecke algebra \dots\ from Hida to Diamond, Flach and Guo''; the hypothesis of Theorem thm:main and
\cite[Theorem 0.1]{DFG04}): \Cref{thm:congruence:crit}, \Cref{thm:congruence:hida}, \Cref{thm:congruence:dfg01}, \Cref{def:congruence:exceptional-set}.
\item \textbf{Lemma lem:crit}, restated as a statement about congruence modules: \Cref{thm:congruence:crit}, \Cref{prop:congruence:crit-general}.
\item \textbf{Section 4.1} ($M^\Sigma$, $f^\Sigma$, the integral structure, $M_{f^\Sigma,\lambda}$, $\hat\delta_!$, $\eta^\Sigma_{f,\lambda}$): \Cref{def:congruence:depletion-recall} to \Cref{prop:congruence:eta-sigma}.
\item \textbf{Theorem thm:dfg14 and the paragraph after it}: \Cref{thm:congruence:dfg14}, \Cref{prop:congruence:three-cases}.
\item \textbf{Proposition prop:pullback and Lemma lem:hecke}: \Cref{prop:congruence:eta-sigma}(iii), \Cref{thm:congruence:lem-hecke}.
\item \textbf{Lemma lem:ddt417}: \Cref{lem:congruence:ddt417}, \Cref{rem:congruence:ddt417-commentary}. \textbf{Remark rem:free}: \Cref{rem:congruence:rem-free}.
\item \textbf{Section 4.2} ($S_f$, Proposition prop:Sf, Theorems thm:dfg37 and thm:dfg27, the minimal-ramification paragraph): \Cref{def:congruence:exceptional-set} to \Cref{rem:congruence:ell-three}.
\item \textbf{Lemmas lem:dfg21, lem:fpr, lem:h0}: used in Step 4 of \Cref{thm:congruence:main}, proved in \Cref{ch:blochkato}.
\item \textbf{Section 4.3} (the proof of Theorem thm:main, Corollary cor:explicit): \Cref{thm:congruence:main}, \Cref{cor:congruence:explicit}.
\item \textbf{Remark rem:ko}: \Cref{rem:congruence:rem-ko-1}, \Cref{prop:congruence:kim-ota}.
\item \textbf{Theorem thm:twist and the paragraph after it}: \Cref{thm:congruence:twist}, \Cref{ex:congruence:twist-cases}.
\item \textbf{Theorem thm:quant, its proof and Remark rem:afterquant}: \Cref{thm:congruence:quant} to \Cref{rem:congruence:after-quant}.
\item \textbf{Remark rem:ns, Open questions 1, 2 and 6}: \Cref{rem:congruence:rem-ns}, \Cref{rem:congruence:open-question-2}, \Cref{rem:congruence:open-question-1}, \Cref{rem:congruence:step1-commentary}.
\item \textbf{Every citation of \cite{DFG04}}: \Cref{rem:congruence:imported-table}.
\end{itemize}

\section{Overview and conventions}\label{sec:congruence:overview}

This section fixes notation and gives a map of Section 4 of the paper. Every sentence of that section is either a citation of an imported
theorem, a result proved in an earlier chapter, or a step proved here; the table below says which. The two warnings concern the two places where
the paper's notation can mislead: the letters $I$ and $J$ of Lemma lem:crit, and the two different congruence ideals of Section 4.

\begin{convention}[standing notation]\label{conv:congruence:setup}
\begin{enumerate}[label=(\alph*)]
\item \emph{Newforms and coefficients.} $f$ is a newform of weight $2$ on $\Gamma_0(M)$ with trivial character. In Sections \ref{sec:congruence:hecke-local} to
\ref{sec:congruence:taylor-wiles} its coefficients may lie in a number field $K$ with a prime $\lambda\mid\ell$; then $\mathcal O=\mathcal O_{K,\lambda}$, with uniformiser $\varpi$ and
valuation $v$, $v(\varpi)=1$. From Section \ref{sec:congruence:assembly} on, $f=f_E$ has rational coefficients, $K=\Q$, $\lambda=\ell$ and $\mathcal O=\Zl$ (the paper's
``$\lambda=\ell$''). The prime $\ell$ is odd and $\ell\nmid M$ unless we say otherwise.
\item \emph{Hecke algebras.} $\T=\T_M=\Z[T_n:n\ge1]\subset\End S_2(\Gamma_0(M))$ is the \emph{full} Hecke algebra, the operators $U_p=T_p$ for $p\mid M$ included
(\Cref{ch:hecke}). $\T_\ell=\T\otimes\Zl$, $\T_{\mathcal O}=\T\otimes\mathcal O$; for a maximal ideal $\mf m$ of $\T_\ell$ (or $\T_{\mathcal O}$), $\T_{\mf m}$ is the localisation, a direct factor
cut out by an idempotent $e_{\mf m}$. $\T$ acts on $H^1(X_0(M),\Zl)$ by pull-back $t^*$ (\Cref{ch:moddeg}).
\item \emph{Congruence data.} For a surjective homomorphism of $\mathcal O$-algebras $\pi\colon\T_{\mathcal O}\onto\mathcal O$ we write $\wp=\ker\pi$, $I=\Ann(\wp)$, $\eta(\pi)=\pi(I)$ (the
\emph{congruence ideal}), $\mathcal C(\pi)=\T_{\mathcal O}/(\wp+I)$ (the \emph{congruence module}) and $\mathrm{Cot}(\pi)=\wp/\wp^2$ (the \emph{cotangent module}). The letters
$\mathcal C$ and $\mathrm{Cot}$ are this chapter's; the literature writes $\Omega$ and $\Phi$, which the book reserves for other objects (Appendix B).
\item \emph{Two congruence ideals.} $\eta^{\mathrm{Hecke}}(\pi)$ is $\eta(\pi)$ of (c), an invariant of the algebra $\T$; $\eta^\Sigma_{f,\lambda}$ is the naive $\Sigma$-finite congruence ideal of
Diamond, Flach and Guo (\Cref{def:congruence:eta-sigma}), an invariant of a lattice with a pairing. For a newform at its level, $\eta^{\mathrm{Hecke}}_f=\eta^{\mathrm{Hecke}}(\phi_f)$.
\item \emph{Lengths.} All counts of finite $\mathcal O$-modules are lengths $\len_{\mathcal O}$; for $\mathcal O=\Zl$, $\#X=\ell^{\len X}$ (Appendix B, B-C10).
\item \emph{Frobenius and duals.} $\Frob_p$ is an arithmetic Frobenius and $\Frob_p^{-1}$ the geometric one; $L$-factors use the geometric Frobenius,
$L_p(V,s)=\det(1-\Frob_p^{-1}p^{-s}\mid V^{I_p})^{-1}$ (B-C2). $M^\vee=\Hom(M,\mathcal O)$ is the linear dual, $X^*$ the Pontryagin dual, $V^D=\Hom(V,\Ql(1))$, and $t^\dagger$ the
adjoint of an operator. With $T=T_\ell E$: $A=\ad T$, $A_f=\ad V_\ell E$, $W=A\otimes\Qlmod$, $T^D\simeq A(1)$, $V^D$, $W^D$ as in \Cref{ch:blochkato}.
\item \emph{Twists.} For a finite set $S$ of odd primes, $D=\prod_{q\in S}q$, $\chi_S=\prod_{q\in S}\chi_{q^*}$, $\iota=\iota_S\colon h\mapsto h\otimes\chi_S=\sum\chi_S(n)a_n(h)q^n$ (the naive twist,
\Cref{ch:hecke}); $E_S=E\otimes\chi_S$ of conductor $N_S$, $g=f_{E_S}$, and $E^{(q)}=E\otimes\chi_{q^*}$. In this chapter $D$ has no other meaning.
\item \emph{Depletion.} $\delta_p=2,1,0$ according as $p\nmid M$, $p\,\|\,M$, $p^2\mid M$; $M^\Sigma=M\prod_{p\in\Sigma}p^{\delta_p}$; $f^\Sigma=\sum_{(n,\Sigma)=1}a_n(f)q^n$.
\end{enumerate}
\end{convention}

\begin{warning}[the letters of Lemma lem:crit]\label{warn:congruence:letters}
In the paper's Lemma lem:crit, $I=\ker\phi_f$ and $J=\Ann_\T(f^\perp)$. In Lemma lem:ddt417 and in this chapter, $\wp=\ker\pi$ and $I=\Ann_\T(\wp)$. So for $\pi=\phi_f$ the paper's
$I$ is our $\wp$, and the paper's $J$ is our $I$ (that $J=\Ann_\T(\ker\phi_f)$ is part of Lemma lem:crit, \Cref{thm:congruence:crit}). We use the paper's letters only when we
quote Lemma lem:crit.
\end{warning}

\begin{warning}[two congruence ideals]\label{warn:congruence:two-ideals}
The Hecke congruence ideal $\eta^{\mathrm{Hecke}}(\pi)$ and the cohomological ideal $\eta^\Sigma_{f,\lambda}$ of Diamond, Flach and Guo are different objects. The first is defined from the
ring $\T$ and the homomorphism $\pi$ alone; the second from the pairing $\hat\delta_!$ on a rank-two lattice in cohomology. Lemma lem:ddt417 compares them
($v(\eta^{\mathrm{Hecke}})\ge v(\eta^{\mathrm{coh}})$), with equality when the cohomology is free over the Hecke algebra; for the newform of an optimal curve at $\ell\nmid2N$ they agree
without any freeness hypothesis (\Cref{rem:congruence:two-ideals}). The paper's Section 4 uses both, and Step 2 of the proof of Theorem thm:quant is exactly the comparison.
\end{warning}

The paper's Section 4, sentence by sentence, against this chapter and the chapters that supply the inputs:
\begin{center}\small
\begin{tabular}{p{0.36\textwidth}p{0.26\textwidth}p{0.28\textwidth}}
\toprule
paper (Section 4) & here & inputs proved in\\
\midrule
4.1: $M^\Sigma$, $f^\Sigma$, the integral structure, $M_{f^\Sigma,\lambda}$, $\hat\delta_!$, $\eta^\Sigma_{f,\lambda}$ & \Cref{def:congruence:depletion-recall} to \Cref{def:congruence:eta-sigma} & \Cref{ch:hecke,ch:moddeg,ch:modules}\\
Theorem thm:dfg14 and the paragraph after it & \Cref{thm:congruence:dfg14}, \Cref{prop:congruence:three-cases} & \Cref{ch:lfunctions,ch:modgalrep}\\
Proposition prop:pullback, Lemma lem:hecke & \Cref{thm:congruence:lem-hecke} & \Cref{ch:moddeg}\\
Lemma lem:ddt417 & \Cref{lem:congruence:ddt417} & \Cref{ch:modules}\\
Remark rem:free & \Cref{rem:congruence:rem-free} & \Cref{ch:moddeg,ch:modules}\\
4.2: $S_f$, Proposition prop:Sf & \Cref{def:congruence:exceptional-set} & \Cref{ch:galrep,ch:blochkato}\\
Theorems thm:dfg37, thm:dfg27 & \Cref{thm:congruence:dfg37}, \Cref{thm:congruence:dfg27} & imported \cite{DFG04}\\
Lemmas lem:dfg21, lem:fpr, lem:h0 & Step 4 of \Cref{thm:congruence:main} & \Cref{ch:blochkato,ch:galcoh,ch:galrep}\\
4.3: proof of Theorem thm:main, Corollary cor:explicit & \Cref{thm:congruence:main} & \Cref{ch:blochkato,ch:lfunctions,ch:moddeg}\\
Remark rem:ko & \Cref{rem:congruence:rem-ko-1}, \Cref{prop:congruence:kim-ota} & \Cref{ch:blochkato,ch:modgalrep}\\
4.4: Theorem thm:twist & \Cref{thm:congruence:twist} & \Cref{ch:hecke,ch:modgalrep}\\
Theorem thm:quant and its proof & \Cref{thm:congruence:quant} & \Cref{ch:hecke,ch:reduction,ch:lfunctions}\\
Remark rem:afterquant & \Cref{lem:congruence:quant-step1}, \Cref{rem:congruence:after-quant} & \Cref{ch:hecke}\\
\bottomrule
\end{tabular}
\end{center}

For a reader with a physics background, two images may help, one sentence each. The congruence module measures how far a lattice fails to split as the
direct sum of its intersections with two complementary eigenspaces, as the unit cell of a crystal need not be a product of cells of two sublattices. The
numerical criterion of Wiles (\Cref{thm:congruence:numerical-criterion}) is an index theorem: the size of a tangent space bounds a defect, with equality
exactly in the most rigid (complete intersection) case.

\section{Hecke algebras over $\Zl$ and their local components}\label{sec:congruence:hecke-local}

This section recalls, with references, the structure of the Hecke algebra over $\Zl$ that every later section uses: the splitting into local
components, the maximal ideal attached to a newform, and which forms live in a given component. The paper uses these facts in Lemma lem:hecke
(``$\mf m$ the maximal ideal of $\T\otimes\Zl$ containing $\ell$ and $\ker\phi_f$''), in Remark rem:free and in Steps 1 and 2 of the proof of Theorem thm:quant.
Nothing here is new; the point is to have the statements in one place, in the form in which they are used.

\begin{proposition}[local components, recollection]\label{prop:congruence:hecke-local}
Let $M\ge1$, $\T=\T_M$, $d=\dim S_2(\Gamma_0(M))$ and $\ell$ a prime.
\begin{enumerate}[label=(\roman*)]
\item $\T$ is a free $\Z$-module of rank $d$; so $\T_\ell$ is a finite free $\Zl$-algebra, it has finitely many maximal ideals, all containing $\ell$, and
$\T_\ell=\prod_{\mf m}\T_{\mf m}$ with $\T_{\mf m}=e_{\mf m}\T_\ell$ local and finite free over $\Zl$.
\item Let $\phi\colon\T\to\overline\Q_\ell$ be a ring homomorphism (a system of eigenvalues); its values lie in the ring of integers $\mathcal O$ of a finite extension of $\Ql$, and
$\mf m_\phi=\ker(\T_\ell\to\mathcal O\to\mathcal O/\varpi)$ is a maximal ideal. If $\phi=\phi_f$ for a newform $f$ with rational coefficients, $\mf m_{f,\ell}=\mf m_{\phi_f}$ is generated by $\ell$ and
the elements $T_n-a_n(f)$.
\item $e_{\mf m}\bigl(S_2(\Gamma_0(M);\Zl)\otimes\overline\Q_\ell\bigr)$ is the sum of the generalised eigenspaces $V(\phi)$ of the systems $\phi$ with $\mf m_\phi=\mf m$, that is, with
$\phi\equiv\phi'$ modulo the maximal ideal of $\overline\Z_\ell$ for one (hence every) $\phi'$ with $\mf m_{\phi'}=\mf m$. If $V(\phi)$ lies in the old space $V_h$ of a newform $h$ and
$\mf m_\phi=\mf m_{f,\ell}$, then $a_p(h)\equiv a_p(f)$ for all $p\nmid M\ell$, so $\bar\rho_h^{\mathrm{ss}}\simeq\bar\rho_f^{\mathrm{ss}}$.
\item The pairing $S_2(\Gamma_0(M);\Zl)\times\T_\ell\to\Zl$, $(h,t)\mapsto a_1(th)$, is perfect and restricts to a perfect pairing $e_{\mf m}S_2(\Gamma_0(M);\Zl)\times\T_{\mf m}\to\Zl$. The pairing
$\hat\delta_!(x,y)=\langle x,w_M^*y\rangle$ on $H^1(X_0(M),\Zl)$ is perfect and alternating, every $t\in\T$ is self-adjoint for it, and $H^1(X_0(M),\Zl)=\bigoplus_{\mf m}e_{\mf m}H^1(X_0(M),\Zl)$
is an orthogonal decomposition on whose summands $\hat\delta_!$ is perfect.
\end{enumerate}
\end{proposition}

\begin{proof}
(i) The freeness of $\T$ is proved in \Cref{ch:hecke} (the perfect pairing with the integral forms); the splitting of a finite free algebra over the complete ring $\Zl$
is the idempotent-lifting theorem of \Cref{ch:modules}, which also shows that the maximal ideals contain $\ell$. (ii) The eigenvalues of $T_n$ are algebraic integers
(\Cref{ch:hecke}), so $\phi(\T)$ is a finitely generated $\Z$-module of integral elements and $\phi(\T_\ell)$ an order in a finite extension; its image in the residue field is
a field, so $\mf m_\phi$ is maximal. If $\phi=\phi_f$ is $\Z$-valued, $\T_\ell/(\ell,T_n-a_n(f):n\ge1)$ is a quotient of $\Z/\ell$ (each $T_n$ becomes a scalar), and it is nonzero
because it maps onto $\F_\ell$ through $\phi_f$; so the ideal generated is $\mf m_{f,\ell}$. (iii) The first sentence is the theorem on generalised eigenspaces over a complete
discrete valuation ring of \Cref{ch:modules}, applied to $L=S_2(\Gamma_0(M);\Zl)$, which is a faithful $\T_\ell$-module. For the second, $\phi(T_p)=a_p(h)$ for $p\nmid M$ on $V_h$
(multiplicity one, \Cref{ch:hecke}), so $\mf m_\phi=\mf m_{f,\ell}$ gives $a_p(h)\equiv a_p(f)$ modulo the maximal ideal for $p\nmid M\ell$, and Brauer-Nesbitt with Chebotarev gives the
isomorphism of semisimplifications (\Cref{ch:modgalrep}). (iv) The first statement is \Cref{ch:hecke} (the pairing is perfect over $\Z$ and stays perfect after $\otimes\Zl$), and
its restriction to components is \Cref{ch:hecke} again (the local components proposition). The second is \Cref{ch:moddeg}: $t^\dagger=w_M^*t^*w_M^*$ for the cup product, so
$\hat\delta_!$ makes $\T$ self-adjoint, and an idempotent of a self-adjoint commutative algebra gives an orthogonal decomposition on whose summands a perfect
pairing stays perfect (\Cref{ch:modules}).
\end{proof}

\begin{paperlink}
The paragraph after the paper's Lemma lem:hecke says that ``the other local factors of $\T\otimes\Zl$ lie in the kernel $\wp$, $\phi_f$ factors through $\T_{\mf m}$, and
$\operatorname{Ann}(\wp)$ decomposes along the idempotents; and the idempotent $e_{\mf m}$ is self-adjoint for $\hat\delta_!$, so $H^1(X_0(N),\Zl)$ splits orthogonally into its local
components and the pairing is perfect on each''. These are \Cref{prop:congruence:hecke-local}(iv) and \Cref{prop:congruence:congruence-module}(iv) below.
\end{paperlink}

\begin{example}[level $539$ at $\ell=3$]\label{ex:congruence:level-539}
This example runs through Sections \ref{sec:congruence:modules} and \ref{sec:congruence:quant}. At level $539=7^2\cdot11$ consider four eigenforms of $\T_{539}$ with integer
coefficients: the newforms $f_{\text{539c}}$ and $f_{\text{539d}}$, and the depleted forms $(f_{\text{77c}})^{\{7\}}$ and $(f_{\text{11a}})^{\{7\}}$, which live at level $539$ because $\delta_7=1$ for
\curve{77c1} (multiplicative at $7$, so $77\cdot7=539$) and $\delta_7=2$ for \curve{11a1} (good at $7$, so $11\cdot49=539$). The curve \curve{539c1} has type $\kod{I}_3^*$ at $7$ and is the twist of
\curve{77c1} by $\chi_{-7}$; \curve{539d1} has type $\kod{I}_0^*$ at $7$ and is the twist of \curve{11a1} by $\chi_{-7}$. So $\iota_{\{7\}}$ sends the two depleted forms to the two newforms.
Computed with PARI/GP (all $n\le130$, beyond Sturm's bound $112$ for $\Gamma_0(539)$): $a_n(f_{\text{539c}})\equiv a_n(f_{\text{539d}})\pmod3$ for every $n$, and
$a_n((f_{\text{77c}})^{\{7\}})\equiv a_n((f_{\text{11a}})^{\{7\}})\pmod3$ for every $n$; but $a_3(f_{\text{539c}})=-2$ while $a_3((f_{\text{77c}})^{\{7\}})=2$. So $\T_{539}\otimes\Z_3$ has a maximal ideal $\mf m_1$
containing the two newforms' systems and another, $\mf m_2$, containing the two depleted systems; $\iota$ exchanges them (it multiplies $T_p$ by $\chi_{-7}(p)$). The Hecke
congruence modules, computed as the index $[S_2(539;\Z):\Z g\oplus(g^c\cap S_2(539;\Z))]$ with $g^c$ the $\T$-stable complement (\Cref{prop:congruence:crit-general}), are
\[
\#\mathcal C(f_{\text{539c}})=288,\quad\#\mathcal C(f_{\text{539d}})=504,\quad\#\mathcal C((f_{\text{77c}})^{\{7\}})=288,\quad\#\mathcal C((f_{\text{11a}})^{\{7\}})=504,
\]
with $288=2^53^2$ and $504=2^33^27$. The modular degrees are $\deg\varphi_{\curve{539c1}}=288$ and $\deg\varphi_{\curve{539d1}}=72$ (\Cref{ch:moddeg}): at $\ell=7$, where $7^2\mid539$, the
congruence number $r(\curve{539d1})=504$ and the degree differ. That the depleted modules equal the congruence numbers of the twisted newforms is
\Cref{lem:congruence:quant-step1} at every $\ell\ne7$. (Computational note: the index must be computed with the integral structure cut out by Sturm's bound $\frac{2}{12}[\SL_2(\Z):\Gamma_0(539)]=112$,
which is what \texttt{mfsturm([539,2])} returns; applied to a space created by \texttt{mfinit}, \texttt{mfsturm} returns $88$ or $50$ according to the space, too few coefficients to test
integrality. Truncating at fewer than $112$ coefficients gives wrong orders that depend on the truncation: with $a_1,\dots,a_{50}$ the four orders come out as $1152$, $2016$, $5760$, $10080$, with
$a_1,\dots,a_{51}$ as $288$, $504$, $1440$, $2520$; with $a_1,\dots,a_{112}$ they are the correct $288$, $504$, $288$, $504$.)
\end{example}

\section{Congruence modules and congruence ideals}\label{sec:congruence:modules}

This section defines the congruence ideal of a surjection from a finite free algebra, records its basic properties (proved in \Cref{ch:modules}),
and reads the paper's Lemma lem:crit in this language: the congruence number $r_E$ is the congruence ideal of $\phi_f$ in the full Hecke algebra, and
the same holds for the depleted form $g^S$ of Theorem thm:quant once the right complement is used. In the paper this is the sentence of Section 1
``The congruence number is the congruence ideal of $f$ in the full Hecke algebra of level $N$ (Lemma lem:crit), and the theory of such ideals, from Hida to
Diamond, Flach and Guo, expresses them through the adjoint Selmer group and the adjoint $L$-value'', and it underlies Lemma lem:hecke and Step 1 of the proof
of Theorem thm:quant. Hida's theorem, which started the subject, is stated at the end as context; the paper does not use it in a proof.

\begin{definition}[congruence ideal, congruence module, cotangent module]\label{def:congruence:congruence-module}
Let $\mathcal O$ be a discrete valuation ring (in practice complete) with fraction field $K$, $\T$ a commutative $\mathcal O$-algebra which is free of finite rank as an $\mathcal O$-module, and
$\pi\colon\T\onto\mathcal O$ a surjective homomorphism of $\mathcal O$-algebras. Put $\wp=\ker\pi$ and $I=\Ann_\T(\wp)$. The \emph{congruence ideal} of $\pi$ is $\eta(\pi)=\pi(I)\subseteq\mathcal O$, the
\emph{congruence module} is $\mathcal C(\pi)=\T/(\wp+I)$, and the \emph{cotangent module} is $\mathrm{Cot}(\pi)=\wp/\wp^2$.
\end{definition}

\begin{proposition}[properties of congruence ideals]\label{prop:congruence:congruence-module}
In the situation of \Cref{def:congruence:congruence-module}:
\begin{enumerate}[label=(\roman*)]
\item $\pi$ induces an isomorphism $\mathcal C(\pi)\simeq\mathcal O/\eta(\pi)$.
\item $\eta(\pi)\ne0$ if and only if $\wp\cap I=0$, if and only if the local factor of the finite $K$-algebra $\T_K=\T\otimes K$ through which $\pi_K$ factors is $K$ itself.
\item If $\eta(\pi)\ne0$, then $\pi$ restricts to an isomorphism $I\simeq\eta(\pi)$, so $I$ is free of rank one over $\mathcal O$; $\Ann_\T(I)=\wp$; and $\eta(\pi)=\set{a\in\mathcal O:ae_\pi\in\T}$, where
$e_\pi\in\T_K$ is the idempotent of the factor $K$ of (ii).
\item If $\T=\prod_{\mf m}\T_{\mf m}$ is the decomposition into local factors (which exists when $\mathcal O$ is complete) and $\pi$ factors through $\T_{\mf m_0}$, then the factors $\T_{\mf m}$, $\mf m\ne\mf m_0$, lie in $\wp$, $I\subseteq\T_{\mf m_0}$,
and $\eta(\pi)=\eta(\pi|_{\T_{\mf m_0}})$: the congruence ideal is computed in the local factor of $\pi$.
\item $\mathrm{Cot}(\pi)$ is a finitely generated $\mathcal O$-module, and it is finite (of finite length) if and only if the factor of $\T_K$ through which $\pi_K$ factors is $K$.
\end{enumerate}
\end{proposition}

\begin{proof}
(i) The structure map $\mathcal O\to\T$ is a section of $\pi$, so $\T=\mathcal O\cdot1\oplus\wp$ as $\mathcal O$-modules, $t=\pi(t)+(t-\pi(t))$. Hence $\pi$ induces $\T/\wp\simeq\mathcal O$, which carries
the image of $I$ onto $\pi(I)$, and $\T/(\wp+I)\simeq\mathcal O/\pi(I)$. (ii) and (iii) are the proposition on congruence ideals of \Cref{ch:modules}; as Lemma lem:ddt417 uses them over
a discrete valuation ring that need not be complete, we give the argument, which uses no completeness and only that $\T$ is torsion-free (not that it is reduced). The finite $K$-algebra
$\T_K$ is a product of local artinian $K$-algebras; write $\T_K=A\times B$ with $A$ the factor through which $\pi_K$ factors and $\mf n$ the maximal ideal of $A$, so that $\wp_K=\wp\otimes K=\mf n\times B$
and $\wp=\T\cap\wp_K$. As $\T$ is torsion-free, $t\in\T$ kills $\wp$ iff it kills $\wp_K$, so $I=\T\cap\Ann_{\T_K}(\wp_K)=\T\cap(\Ann_A(\mf n)\times0)$. If $A=K$, then $\mf n=0$ and $I=\T\cap Ke_\pi$ with
$e_\pi=(1,0)$; since $\pi(ae_\pi)=a$, $\pi$ maps $I$ isomorphically onto $\eta(\pi)=\set{a\in K:ae_\pi\in\T}\subseteq\mathcal O$, which is nonzero because a nonzero multiple of $e_\pi$ lies in $\T$; so
$\wp\cap I=0$, and $\Ann_\T(I)=\T\cap(0\times B)=\wp$. If $A\ne K$, then $\mf n\ne0$ and $\Ann_A(\mf n)\subseteq\mf n$ (a unit of $A$ kills nothing), so $I\subseteq\T\cap(\mf n\times0)\subseteq\wp$ and $\eta(\pi)=0$;
moreover $\Ann_A(\mf n)\ne0$ (it contains $\mf n^{k-1}$ for the least $k$ with $\mf n^k=0$), so $I\ne0$ and $\wp\cap I=I\ne0$. This proves (ii) and (iii). For (iv), note that
the idempotent $1-e_{\mf m_0}$ maps to $0$ under $\pi$ (the image of an idempotent in the domain $\mathcal O$ is $0$ or $1$, and $\pi(e_{\mf m_0})=1$), so $\prod_{\mf m\ne\mf m_0}\T_{\mf m}=(1-e_{\mf m_0})\T\subseteq\wp$; an
element of $I$ kills $1-e_{\mf m_0}$, hence lies in $e_{\mf m_0}\T$, and inside $\T_{\mf m_0}$ it kills $\wp\cap\T_{\mf m_0}$; conversely an element of $\T_{\mf m_0}$ killing $\wp\cap\T_{\mf m_0}$ kills
$\wp=(\wp\cap\T_{\mf m_0})\times\prod_{\mf m\ne\mf m_0}\T_{\mf m}$. So $I=\Ann_{\T_{\mf m_0}}(\wp\cap\T_{\mf m_0})$. (v) $\wp$ is a submodule of the finitely generated module $\T$ over the noetherian ring
$\mathcal O$. Localisation is exact, so $\mathrm{Cot}(\pi)\otimes K=\wp_K/\wp_K^2$, where $\wp_K=\ker\pi_K$. Write $\T_K=A\times B$ with $A$ the local factor through which $\pi_K$ factors, so
$\wp_K=\mf n\times B$ with $\mf n$ the maximal ideal of $A$, and $\wp_K/\wp_K^2=\mf n/\mf n^2$ (as $B^2=B$). By Nakayama's lemma in the artinian local ring $A$, $\mf n/\mf n^2=0$ iff $\mf n=0$ iff $A=K$.
A finitely generated $\mathcal O$-module is finite iff it is killed by $\otimes K$.
\end{proof}

\begin{example}[three toy algebras]\label{ex:congruence:toy}
(a) For $r\ge1$ let $\T_r=\set{(a,b)\in\mathcal O^2:a\equiv b\pmod{\varpi^r}}$ and $\pi(a,b)=a$. Then $\wp=0\times\varpi^r\mathcal O$, $I=\varpi^r\mathcal O\times0$, $\eta=\varpi^r\mathcal O$, $\mathcal C(\pi)\simeq\mathcal O/\varpi^r$,
and $\wp^2=0\times\varpi^{2r}\mathcal O$, so $\mathrm{Cot}(\pi)\simeq\mathcal O/\varpi^r$: $\len\mathrm{Cot}(\pi)=v(\eta)=r$. This is the algebra of two eigenforms congruent modulo $\varpi^r$, and
$\T_r\simeq\mathcal O[x]/(x(x-\varpi^r))$ by $x\mapsto(0,\varpi^r)$ is a complete intersection.
(b) Let $\T_3^\flat=\set{(a,b,c)\in\mathcal O^3:a\equiv b\equiv c\pmod\varpi}$ and $\pi(a,b,c)=a$. Then $\wp=0\times\varpi\mathcal O\times\varpi\mathcal O$, $I=\varpi\mathcal O\times0\times0$, $\eta=\varpi\mathcal O$, but
$\wp^2=0\times\varpi^2\mathcal O\times\varpi^2\mathcal O$ and $\mathrm{Cot}(\pi)\simeq(\mathcal O/\varpi)^2$: $\len\mathrm{Cot}(\pi)=2>1=v(\eta)$. Three pairwise congruent forms; $\T_3^\flat$ is not a complete
intersection (it is not even Gorenstein, \Cref{ch:modules}).
(c) Let $\T=\mathcal O[x]/(x^2)$ and $\pi(x)=0$. Then $\wp=\mathcal Ox=I$, $\eta=0$ and $\mathrm{Cot}(\pi)=\mathcal Ox$ is infinite: the factor of $\T_K$ at $\pi$ is $K[x]/(x^2)$, not $K$. This is
the shape of a Jordan block of a Hecke operator, such as $U_2$ at level $88$ on the forms $f_{\text{11a}}(dz)$, $d\mid8$ (\Cref{ch:hecke}).
These three cases are the whole content of Wiles's numerical criterion (\Cref{thm:congruence:numerical-criterion}) in miniature: $\len\mathrm{Cot}\ge v(\eta)$, with equality for
the complete intersection (a), strict inequality for (b), and both sides degenerate in (c).
\end{example}

\begin{definition}[the Hecke congruence ideal of a newform]\label{def:congruence:newform-ideal}
Let $f$ be a newform of weight $2$ on $\Gamma_0(M)$, $K_f$ its coefficient field, $\lambda$ a prime of $K_f$ above $\ell$, and $\mathcal O$ the completion of the ring of integers of $K_f$ at
$\lambda$. The eigencharacter $\phi_f\colon\T\to\mathcal O_{K_f}$ extends to a surjection $\phi_f\colon\T_{\mathcal O}\onto\mathcal O$ (its image is an $\mathcal O$-submodule containing $1$, hence $\mathcal O$). Its
congruence ideal $\eta^{\mathrm{Hecke}}_f=\eta(\phi_f)\subseteq\mathcal O$ is the \emph{Hecke congruence ideal} of $f$. More generally, for an eigenform $g$ of $\T_M$ with eigencharacter $\pi$, we write
$\eta^{\mathrm{Hecke}}(g)=\eta(\pi)$.
\end{definition}

\begin{proposition}[the Hecke congruence ideal is nonzero]\label{prop:congruence:newform-ideal}
In \Cref{def:congruence:newform-ideal}, $\eta^{\mathrm{Hecke}}_f\ne0$. More generally $\eta(\pi)\ne0$ for every eigencharacter $\pi$ whose generalised eigenspace in $S_2(\Gamma_0(M))\otimes\overline K$ is a line,
in particular for $\pi$ the eigencharacter of a depleted form $f^\Sigma$ at level $M^\Sigma$.
\end{proposition}

\begin{proof}
$S_2(\Gamma_0(M))\otimes\C$ is free of rank one over $\T\otimes\C$ (\Cref{ch:hecke}, the structure of $\T$), so the local factor of $\T\otimes\overline K$ at $\pi$ has the dimension of the
generalised eigenspace of $\pi$. If that is a line, the factor is $\overline K$, so the factor of $\T_K$ through which $\pi$ factors is $K$, and $\eta(\pi)\ne0$ by
\Cref{prop:congruence:congruence-module}(ii). For a newform the generalised eigenspace is the line $\C f$ (multiplicity one and strong multiplicity one, \Cref{ch:hecke}); for
a depleted form at level $M^\Sigma$ it is the line $\C f^\Sigma$ (\Cref{ch:hecke}, depletion).
\end{proof}

\begin{theorem}[Lemma lem:crit as a statement about congruence modules]\label{thm:congruence:crit}
Let $E$ have conductor $N$, $f=f_E$, $\T=\T_N$, $\wp=\ker\phi_f$ and $I=\Ann_\T(\wp)$. Then $I=\Ann_\T(f^\perp)$, the congruence module $\mathcal C(\phi_f)=\T/(\wp+I)$ is cyclic of order $r_E$,
$\eta^{\mathrm{Hecke}}_f=r_E\Z$, and the integers $m$ for which $f\equiv h\pmod m$ with $h\in f^\perp\cap S_2(\Gamma_0(N);\Z)$ are exactly the divisors of $r_E$; the same holds over $\Zl$ with $r_E$
replaced by its $\ell$-part. Moreover, with $L=f^\perp\cap S_2(\Gamma_0(N);\Z)$, the perfect pairing $S_2(\Gamma_0(N);\Z)\times\T\to\Z$ induces
\[
S_2(\Gamma_0(N);\Z)/(\Z f\oplus L)\ \simeq\ \Hom\bigl(\mathcal C(\phi_f),\Q/\Z\bigr):
\]
the congruence module is dual to the failure of the integral forms to split along the decomposition $S_2=\C f\oplus f^\perp$.
\end{theorem}

\begin{proof}
Everything except the last display is the paper's Lemma lem:crit, which \Cref{ch:modgalrep} proves in full (with the decomposition $\T\otimes\Q=\Q\times\T'$ of \Cref{ch:hecke});
we do not repeat it. For the display: the forms killed by $\wp$ are $\Z f$ and those killed by $I$ are $L$ (\Cref{ch:modgalrep}, in the proof of Lemma lem:crit), and for a
perfect pairing $S\times\T\to\Z$ with $\T$ acting compatibly, ideals $\mf a,\mf b$ with $\mf a\cap\mf b=0$, $\T/\mf a$ and $\T/\mf b$ torsion-free of complementary ranks, one has
$S/(S[\mf a]\oplus S[\mf b])\simeq(\T/(\mf a+\mf b))^*$ (\Cref{ch:modules}, the duality between an algebra and a module). Take $\mf a=\wp$, $\mf b=I$; $\wp\cap I=0$ by
\Cref{prop:congruence:newform-ideal} and \Cref{prop:congruence:congruence-module}(ii), $\T/\wp\simeq\Z$ has rank $1$ and $\T/I$ embeds in $\T'$, of rank $\rank\T-1$.
\end{proof}

\begin{paperlink}
Lemma lem:crit is stated in the paper's Section 3 with the letters $I=\ker\phi_f$ and $J$ the annihilator of $f^\perp$ (\Cref{warn:congruence:letters}); the paper's proof says
``Dualising the inclusion $\T\subset\T/I\times\T/J$, whose cokernel is $\T/(I+J)$, shows that $\Z f\oplus L$ has index $r'$''. That sentence is the display of
\Cref{thm:congruence:crit}, read as a statement about orders. Its use in Section 4 is the identification $\eta^{\mathrm{Hecke}}_f\Zl=r_E\Zl$ in Lemma lem:hecke and in Step 1 of the proof of Theorem thm:quant.
\end{paperlink}

\begin{proposition}[Lemma lem:crit for an eigencharacter whose generalised eigenspace is a line]\label{prop:congruence:crit-general}
Let $\pi\colon\T_N\to\Z$ be a ring homomorphism whose generalised eigenspace in $S_2(\Gamma_0(N))\otimes\C$ is a line $\C g$ ($g$ normalised), and let $g^c$ be the $\T$-stable complement of
$\C g$, the sum of the generalised eigenspaces of the other systems of eigenvalues. Put $\Lambda=g^c\cap S_2(\Gamma_0(N);\Z)$. Then $\T/(\ker\pi+\Ann_\T\ker\pi)$ is cyclic of order
$r(\pi)=[S_2(\Gamma_0(N);\Z):\Z g\oplus\Lambda]$, $\eta(\pi)=r(\pi)\Z$, and $g\equiv h\pmod m$ for some $h\in\Lambda$ if and only if $m\mid r(\pi)$. The same holds over $\Zl$ for every prime $\ell$, with
$r(\pi)$ replaced by its $\ell$-part. If $g$ is a newform of level $N$, then $g^c=g^\perp$ and $r(\pi)=r_g$.
\end{proposition}

\begin{proof}
This is the proposition on the congruence module of a rational eigencharacter in \Cref{ch:hecke}, whose proof is that of Lemma lem:crit with $g^c$ in place of $f^\perp$.
\end{proof}

\begin{warning}[the complement must be $\T$-stable]\label{warn:congruence:complement}
For $g$ not a newform at level $N$ (for instance a depleted form), the Petersson orthogonal complement of $\C g$ need not be $\T$-stable, because the operators $U_p$, $p\mid N$, are
not normal (\Cref{ch:hecke} gives an example at level $44$). \Cref{prop:congruence:crit-general} is false with the Petersson complement in general. The paper's Step 1 of the
proof of Theorem thm:quant says ``in the $\T$-stable complement of the $g^S$-line (not the Petersson complement, which is not $\T$-stable when the $U_p$ are involved)'', which is
right; Theorem thm:twist uses $g^\perp$ for a newform $g$ at its own level, where the two complements agree.
\end{warning}

\begin{theorem}[Hida]\label{thm:congruence:hida}
Let $f$ be a normalised newform of weight $k\ge2$ and level $N$, and let $\Omega^\pm_f$ be its periods, normalised by integrality of the modular symbols. Then $\langle f,f\rangle$ divided by
$\Omega^+_f\Omega^-_f$ is an algebraic number, and for primes $\lambda$ outside an explicit finite set (containing the primes dividing $N$ and small primes depending on $k$), $\lambda$ divides this
ratio if and only if $\lambda$ is a congruence prime of $f$, that is, $\lambda$ divides $\eta^{\mathrm{Hecke}}_f$. Through the Rankin-Shimura method, $\langle f,f\rangle$ is a multiple, by explicit factors,
of the naive adjoint $L$-value $L^{\nv}(A_f,1)$.
\end{theorem}

\begin{proof}[Reference]
We do not prove this. The statement combines two results of Hida from 1981: \cite{Hida81} proves that a prime dividing the normalised Petersson norm (outside the excluded set) is a
congruence prime, and a companion paper of the same year proves the converse; we cite only the first, have not checked the precise excluded set, and state the theorem as context. See
\cite{HidaBook} for an exposition; \cite{Ribet83} treats the weight-$2$ case from the side of modular curves. The Petersson norm as a naive $L$-value is proved in the book's normalisation in \Cref{ch:lfunctions}: $\langle f,f\rangle=N\,L^{\nv}(\Sym f,2)/(8\pi^3)$ for a
rational newform of weight $2$.
\end{proof}

\begin{remark}[from Hida to Diamond, Flach and Guo]\label{rem:congruence:hida-to-dfg}
Hida's theorem detects congruence primes; it does not compute the exact power. For an elliptic curve the exact power is the subject of this chapter: the
modular degree is a period ratio (the formula of Zagier and Watkins, \Cref{ch:moddeg}), and \Cref{thm:congruence:main} identifies its $\ell$-adic valuation with the length of a
Selmer group plus local terms. The paper cites \cite[Lemma 2.12]{DFG04} for the formula of the Petersson norm on $\Gamma_1(N)$ as a naive $L$-value, which Diamond,
Flach and Guo deduce from Hida's Rankin-Shimura computation; in the book's normalisation (unnormalised product on $\Gamma_0(N)$) it is the formula recalled in the proof
above, proved in \Cref{ch:lfunctions}. The constants agree. \cite[Lemma 2.12]{DFG04} gives $(f,f)_{\Gamma_1(N)}=(k-1)!\,\delta(N)\,N\,\varphi(N)\,L^{\nv}(A_f,1)/(4^k\pi^{k+1})$, with
$\delta(N)=1$ for $N>2$ and $\delta(N)=2$ for $N\le2$, for the product $(g,h)_\Gamma=(-2i)^{-1}\int_{\Gamma\backslash\HH}g\bar h\,y^{k-2}\,dz\wedge d\bar z=\int_{\Gamma\backslash\HH}g\bar h\,y^k\,d\mu$ of their Section 1.5.1, which is
unnormalised. As $[\bar\Gamma_0(N):\bar\Gamma_1(N)]=\delta(N)\varphi(N)/2$, the product on $\Gamma_0(N)$ is $2/(\delta(N)\varphi(N))$ times theirs, and for $k=2$ and trivial character their formula becomes
$\langle f,f\rangle=N\,L^{\nv}(A_f,1)/(8\pi^3)=N\,L^{\nv}(\Sym f,2)/(8\pi^3)$ at every level.
\end{remark}

\begin{remark}[the three ideals of this chapter, at once]\label{rem:congruence:two-ideals}
Let $E$ be $X_0(N)$-optimal, $f=f_E$, $\ell$ an odd prime with $\ell\nmid N$ and $E[\ell]$ irreducible. Then
\begin{align*}
\eta^{\mathrm{Hecke}}_f\Zl&=r_E\Zl&&(\text{\Cref{thm:congruence:crit}}),\\
v_\ell(r_E)&=v_\ell(\degphi)&&(\text{\cite[Theorem 2.2]{ARS12}}),\\
v_\ell(\degphi)&=v_\lambda(\eta^\emptyset_{f,\lambda})&&(\text{\Cref{thm:congruence:lem-hecke}}),
\end{align*}
so all three ideals agree. Lemma lem:ddt417 (\Cref{lem:congruence:ddt417}) gives $v(\eta^{\mathrm{Hecke}})\ge v(\eta^{\mathrm{coh}})$ in general, with equality when the cohomology is free over the
Hecke algebra; the inequality is strict where freeness fails. Outside the hypotheses it can fail to be an equality: for \curve{99a1} at $\ell=3$, with $3\mid99$, the
Hecke ideal has valuation $v_3(r_E)=v_3(12)=1$ while the ideal of values of $\hat\delta_!$ on $M_{f,\lambda}$ is $\degphi\,\Z_3=4\Z_3=\Z_3$ (\Cref{ch:moddeg}; $r=12$ and $\degphi=4$
recomputed with PARI/GP).
\end{remark}

\section{Wiles's numerical criterion and Lemma lem:ddt417}\label{sec:congruence:ddt417}

This section proves the paper's Lemma lem:ddt417, the piece of commutative algebra that compares the Hecke congruence ideal with the Gram determinant
of a pairing on a module. The paper uses it in Step 2 of the proof of Theorem thm:quant (with $d=2$, for a depleted form) and, in the version of 10 October
2026, in the proof of Lemma lem:hecke. To put it in context we first state Wiles's numerical criterion, of which the lemma is the half that survives in the
presence of a module; the criterion itself is used in this book only in the outline of the Taylor-Wiles method (Section \ref{sec:congruence:taylor-wiles}).

\begin{theorem}[Wiles's numerical criterion; Lenstra]\label{thm:congruence:numerical-criterion}
Let $\mathcal O$ be a complete discrete valuation ring with finite residue field, $\T$ a local $\mathcal O$-algebra, finite and free as an $\mathcal O$-module, and $\pi\colon\T\onto\mathcal O$ a surjection with
$\eta=\eta(\pi)\ne0$. Then:
\begin{enumerate}[label=(\roman*)]
\item $\len_{\mathcal O}\mathrm{Cot}(\pi)\ge v(\eta)$;
\item equality holds if and only if $\T$ is a complete intersection over $\mathcal O$, that is, a quotient of a power series ring $\mathcal O[[x_1,\dots,x_n]]$ by $n$ elements;
\item if $R\onto\T$ is a surjection of complete local noetherian $\mathcal O$-algebras compatible with $\pi$ (so that $\pi$ extends to $\pi_R\colon R\onto\mathcal O$) and
$\len\mathrm{Cot}(\pi_R)\le v(\eta)<\infty$, then $R\to\T$ is an isomorphism and both rings are complete intersections.
\end{enumerate}
\end{theorem}

\begin{proof}[Sketch of proof]
For a complete intersection $\T=\mathcal O[[x_1,\dots,x_n]]/(f_1,\dots,f_n)$ with $\pi(x_i)=0$, both $\len\mathrm{Cot}(\pi)$ and $v(\eta)$ equal $v\bigl(\det(\partial f_i/\partial x_j)(0)\bigr)$: the cotangent module is the
cokernel of the Jacobian matrix at $0$, and Tate's description of the annihilator of $\wp$ in a complete intersection (by the determinant of a matrix expressing the $f_i$ in terms
of the $x_j$) gives $\eta$. The general inequality (i) and the converse in (ii) reduce to this case by writing $\T$ as a quotient of a complete intersection with the same
cotangent module and comparing the two congruence ideals; Lenstra's proof uses the Gorenstein property and a flat deformation. Part (iii) follows from (i) for $\T$ and the
surjection $\mathrm{Cot}(\pi_R)\onto\mathrm{Cot}(\pi)$: $\len\mathrm{Cot}(\pi_R)\ge\len\mathrm{Cot}(\pi)\ge v(\eta)\ge\len\mathrm{Cot}(\pi_R)$ forces equalities, so $\T$ is a complete intersection, and
then a complete intersection is not a proper quotient of a ring with the same cotangent length. We do not give the details; see \cite[Appendix]{Wiles95}, \cite{Lenstra95}, and
\cite[Section 5]{DDT} (Wiles's numerical criterion). The examples of \Cref{ex:congruence:toy} show each case.
\end{proof}

\begin{remark}[why the paper needs a version with a module]\label{rem:congruence:why-a-module}
The ideal that Diamond, Flach and Guo control is not $\eta(\pi)$, an invariant of the ring, but the ideal of values of a pairing on the $\pi$-part of a module $L$ (the
cohomology of a modular curve, \Cref{def:congruence:eta-sigma}). Lemma lem:ddt417 compares the two. It is an equality when $L$ is free over $\T$, which is what a
multiplicity-one theorem gives (the freeness criterion of \Cref{ch:modules} and the Gorenstein theorem of \Cref{ch:moddeg}). The paper says: ``We also need the following piece of
commutative algebra, which is the inequality half of \cite[Lemma 4.17]{DDT}.''
\end{remark}

\begin{lemma}[Lemma lem:ddt417]\label{lem:congruence:ddt417}
Let $\mathcal O$ be a discrete valuation ring, $\T$ a commutative $\mathcal O$-algebra that is free of finite rank as an $\mathcal O$-module, and $L$ a $\T$-module that is free of finite rank over $\mathcal O$,
with a perfect $\mathcal O$-bilinear pairing $\langle\ ,\ \rangle\colon L\times L\to\mathcal O$ for which every $t\in\T$ is self-adjoint. Let $\pi\colon\T\onto\mathcal O$ be a surjective homomorphism of $\mathcal O$-algebras,
$\wp=\ker\pi$, $I=\Ann_\T(\wp)$ and $\eta=\pi(I)$. Suppose that $L[\wp]$ has rank $d$ over $\mathcal O$ and that the Gram determinant $\Delta$ of the pairing on $L[\wp]$ is nonzero. Then
\[
\eta^d\subseteq\Delta\mathcal O,\qquad\text{that is,}\qquad d\cdot v(\eta)\ \ge\ v(\Delta),
\]
and equality holds if $L$ is free over $\T$.
\end{lemma}

\begin{proof}
Here $L[\mf a]$ denotes the submodule of $L$ killed by an ideal $\mf a$; it is free over $\mathcal O$ as a submodule of a free module over a principal ideal domain, so its Gram
determinant is defined up to the square of a unit, and $v(\Delta)$ does not depend on the basis (\Cref{ch:modules}).

\emph{Step 0: the case $\eta=0$, and $\wp\cap I=0$.} If $\eta=0$ the inequality says $\infty\ge v(\Delta)$ and there is nothing to prove. Assume $\eta\ne0$ and fix $t\in I$ with $\pi(t)$ a
generator of $\eta$. Then $\wp\cap I=0$ by \Cref{prop:congruence:congruence-module}(ii) (the argument there uses only that $\T$ is torsion-free, not that it is reduced).

\emph{Step 1: $Q=L/(L[\wp]+L[I])$ is killed by $\eta$.} Let $x\in L$. Since $\wp I=0$, every element of $\wp$ kills $tx$, so $tx\in L[\wp]$. For $s\in I$, $s(\pi(t)-t)\in I\wp=0$, so $(\pi(t)-t)x\in L[I]$.
Hence $\pi(t)x=tx+(\pi(t)-t)x\in L[\wp]+L[I]$, and $\pi(t)$, a generator of $\eta$, kills $Q$.

\emph{Step 2: the sum is direct.} If $y\in L[\wp]\cap L[I]$, then $y$ is killed by $\wp+I$, which contains $\pi(t)=t+(\pi(t)-t)$ ($t\in I$, $\pi(t)-t\in\wp$). So $\pi(t)y=0$, and $y=0$ since $L$ is
torsion-free and $\pi(t)\ne0$. Thus $Q=L/(L[\wp]\oplus L[I])$.

\emph{Step 3: $L[\wp]$ is saturated and $L[\wp]^\perp=L[I]$.} $L[\wp]$ is saturated: if $ax\in L[\wp]$ with $0\ne a\in\mathcal O$, then $a(sx)=s(ax)=0$ for $s\in\wp$, so $sx=0$. Let $y\in L[I]$ and $x\in L[\wp]$. Since $\pi(t)-t\in\wp$ kills $x$ and $t\in I$ kills $y$,
\[
\pi(t)\langle x,y\rangle=\langle tx,y\rangle+\langle(\pi(t)-t)x,y\rangle=\langle x,ty\rangle+0=0,
\]
using self-adjointness of $t$, and $\langle x,y\rangle=0$. Thus
$L[I]\subseteq L[\wp]^\perp$. Conversely let $y\perp L[\wp]$. For $s\in I$ and $z\in L$, $sz\in L[\wp]$ (as in Step 1), so $\langle sy,z\rangle=\langle y,sz\rangle=0$ for all $z$; the pairing being perfect,
$sy=0$. So $y\in L[I]$. (This step is the only one that uses self-adjointness.)

\emph{Step 4: $Q$ is the cokernel of the Gram map.} The map $L\to L[\wp]^\vee=\Hom_{\mathcal O}(L[\wp],\mathcal O)$, $y\mapsto\langle-,y\rangle$, is surjective because $L[\wp]$ is saturated and the
pairing is perfect, and its kernel is $L[\wp]^\perp=L[I]$ (\Cref{ch:modules}, the Gram cokernel theorem). So $L/L[I]\simeq L[\wp]^\vee$, and under this isomorphism the image of $L[\wp]$
is the image of the Gram map $L[\wp]\to L[\wp]^\vee$. Hence $Q\simeq\coker(L[\wp]\to L[\wp]^\vee)$, a module of length $v(\Delta)$ (the cokernel of a square matrix over a discrete
valuation ring has length the valuation of its determinant).

\emph{Step 5: the inequality.} $L[\wp]^\vee\simeq\mathcal O^d$ surjects onto $Q$, so $Q$ is generated by $d$ elements; by Step 1 it is killed by $\eta$. So $Q$ is a quotient of $(\mathcal O/\eta)^d$
and $v(\Delta)=\len Q\le d\,v(\eta)$.

\emph{Step 6: equality for free $L$.} Let $L\simeq\T^r$ as $\T$-modules. Then $L[\wp]\simeq(\T[\wp])^r=I^r$ and $L[I]\simeq\Ann_\T(I)^r=\wp^r$ (\Cref{prop:congruence:congruence-module}(iii)), and
$I\simeq\eta$ is free of rank one over $\mathcal O$, so $d=r$. Then $Q\simeq(\T/(\wp+I))^r\simeq(\mathcal O/\eta)^d$ by \Cref{prop:congruence:congruence-module}(i), of length $d\,v(\eta)$.
\end{proof}

\begin{remark}[Commentary: the paper's statement of Lemma lem:ddt417]\label{rem:congruence:ddt417-commentary}
The paper states the lemma for $\mathcal O$ ``a discrete valuation ring with finite residue field'' and concludes ``$\#\mathcal O/\Delta\le\#(\mathcal O/\eta)^d$''. The book's \Cref{lem:congruence:ddt417} is stated with
lengths over any discrete valuation ring; for a finite residue field of order $\kappa$ the two forms agree, since $\#X=\kappa^{\len X}$. Three steps that the paper's proof compresses are
expanded above. (a) The paper's ``First $\wp\cap I=0$: \dots\ $\T$ being torsion-free'' is Step 0; the paper rightly adds that this does not follow from reducedness, since the
full Hecke algebra can have nilpotents (\Cref{ex:congruence:toy}(c)). (b) The paper's ``Since $L$ is torsion-free and $\wp\cap I=0$, the sum $L[\wp]+L[I]$ is direct'' is Step 2; the operative
reason is that the nonzero scalar $\pi(t)$ lies in $\wp+I$. (c) In the paper's proof of $L[\wp]^\perp=L[I]$, the converse inclusion shows that the single element $t$ kills $y$; this suffices
because $I=\mathcal O t$ (\Cref{prop:congruence:congruence-module}(iii)), and Step 3 above uses every $s\in I$ instead. In the paper's direct inclusion the identity
``$\pi(t)\langle x,y\rangle=\langle tx,y\rangle=\langle x,ty\rangle=0$'' uses $(\pi(t)-t)x=0$ for $x\in L[\wp]$, which is how Step 3 reads it.
\end{remark}

\begin{remark}[on the hypotheses]\label{rem:congruence:ddt417-hypotheses}
(a) Without $\eta\ne0$ the intersection $\wp\cap I$ can be nonzero: for $\T=\mathcal O[x]/(x^2)$ acting on $L=\T$, $\wp=I=\mathcal Ox$. The lemma survives because then $\eta=0$. (b) If $\eta=0$ and $L$
is free over $\T$, then $\Delta=0$ (\Cref{ch:modules}): $I\subseteq\wp$ gives $I^2=0$ and the pairing vanishes on $L[\wp]=I^r$. So in the free case the hypothesis $\Delta\ne0$ forces
$\eta\ne0$. (c) Self-adjointness is used only in Step 3. (d) For the Hecke algebra of a modular curve the module $L$ is $H^1(X_0(N),\Zl)_{\mf m}$, the algebra $\T_{\mf m}$ and the
pairing $\hat\delta_!$ restricted to the $\mf m$-component, which is perfect there (\Cref{prop:congruence:hecke-local}(iv)); by \Cref{prop:congruence:congruence-module}(iv) the
congruence ideal of $\pi$ on $\T_\ell$ is the one computed on $\T_{\mf m}$.
\end{remark}

\begin{example}[the lemma on test cases]\label{ex:congruence:ddt417-checks}
(a) $L=\T_r$ (\Cref{ex:congruence:toy}(a)) with $\langle s,u\rangle=\lambda_0(su)$, $\lambda_0(a,b)=(a-b)/\varpi^r$, a perfect pairing for which $\T_r$ is self-adjoint: $L[\wp]=I=\mathcal O(\varpi^r,0)$,
$d=1$, $\Delta=\lambda_0((\varpi^r,0)^2)=\varpi^r$; equality, $L$ being free. (b) $L=\mathcal O^2$ with $\T_r$ acting diagonally and the standard pairing: $L[\wp]=\mathcal Oe_1$, $\Delta=1$, so $0=v(\Delta)<r=v(\eta)$;
strict, and $L$ is not free over $\T_r$. (c) $L=\T_r^2$ with $\langle x,y\rangle=\lambda_0(x_1y_2-x_2y_1)$, alternating: $d=2$, $\Delta=\varpi^{2r}=\eta^2$. This is the shape of Step 2 of the proof
of Theorem thm:quant. (These computations are in \Cref{ch:modules}.) (d) The modular case outside the hypotheses: for \curve{99a1} and $\ell=3$, $L=H^1(X_0(99),\Z_3)_{\mf m}$, $d=2$,
$v_3(\eta^{\mathrm{Hecke}})=v_3(r_E)=1$ and the ideal of values on $L[\wp]$ is $4\Z_3=\Z_3$, so $v_3(\Delta)=0<2$: $L$ is not free over $\T_{\mf m}$ (\Cref{rem:congruence:two-ideals}).
\end{example}

\section[The lattices of Diamond, Flach and Guo; Lemma lem:hecke]{The lattices and ideals of Diamond, Flach and Guo; Lemma lem:hecke}\label{sec:congruence:dfg-lattices}

This section sets up the objects of the paper's Section 4.1: the depleted form $f^\Sigma$ at level $M^\Sigma$, the integral structure on cohomology, the lattice
$M_{f^\Sigma,\lambda}$, the pairing $\hat\delta_!$ and the naive $\Sigma$-finite congruence ideal $\eta^\Sigma_{f,\lambda}$, in the form in which Diamond, Flach and Guo define them for weight $2$
and trivial character \cite[Sections 1.5.3, 1.6.2, 1.7.3]{DFG04}, and compares them with the cohomology of $X_0(M^\Sigma)$, as the paper does. It then states Lemma lem:hecke
as an equality, which \Cref{ch:moddeg} proves, and Remark rem:free. These are the inputs of Step 1 of the proof of Theorem thm:main and of Steps 2 to 4 of the proof of
Theorem thm:quant.

\begin{definition}[depletion and naive factors, recalled]\label{def:congruence:depletion-recall}
Let $f$ be a newform of weight $2$ on $\Gamma_0(M)$ and $\Sigma$ a finite set of primes. For a prime $p$ put $\delta_p=2$, $1$, $0$ according as $p\nmid M$, $p\,\|\,M$, $p^2\mid M$; for $f=f_E$ these
are good, multiplicative and additive reduction, and $\delta_p=\dim V_\ell E^{I_p}$ for $p\ne\ell$ (\Cref{ch:modgalrep}). Put $M^\Sigma=M\prod_{p\in\Sigma}p^{\delta_p}$ and
$f^\Sigma=\sum_{(n,\Sigma)=1}a_n(f)q^n$. The \emph{naive} local factor of the adjoint $L$-function is $L^{\nv}_p(A_f,s)=L_p(A_f,s)$ if $\delta_p>0$ and $1$ if $\delta_p=0$ (\Cref{ch:lfunctions}).
\end{definition}

By \Cref{ch:hecke} (depletion), $f^\Sigma$ is a normalised eigenform of the full Hecke algebra $\T_{M^\Sigma}$, with $U_pf^\Sigma=0$ for $p\in\Sigma$, and its generalised eigenspace in
$S_2(\Gamma_0(M^\Sigma))$ is the line $\C f^\Sigma$; $M^\Sigma$ is the least level at which $f^\Sigma$ is an eigenform of all $T_n$. The paper's Section 1.1 says it follows ``the conventions of
\cite[\S1.7.2]{DFG04}'' for $\delta_p$ and the naive factors; there $\delta_p$ is the dimension of the inertia invariants of the lattice $M_{f,\lambda}$ of \Cref{def:congruence:eta-sigma}, which is the
same number because $M_{f,\lambda}$ is dual to $T_\ell E$ up to homothety.

\begin{definition}[the integral structure]\label{def:congruence:integral-structure}
Let $M\ge1$ and $\ell\nmid 2M$. Diamond, Flach and Guo \cite[Sections 1.2, 1.4]{DFG04} define, for the open compact subgroup $U_0(M)\subset\GL_2(\widehat\Z)$ and the trivial character, a
$\Zl$-lattice $M(M,1)_{!,\lambda}$ in the first cohomology of the modular curve of level $M$: it is the image of the $U_0(M)$-invariants of the compactly supported cohomology
of the (adelic, open) modular curve $Y_{N'}$ of full level $N'$, for an auxiliary $N'\ge3$ with $U_{N'}\subseteq U_0(M)$, in the torsion-free quotient of the $U_0(M)$-invariants of its
cohomology, with the Hecke action by double cosets. It does not depend on $N'$.
\end{definition}

\begin{proposition}[comparison with $H^1(X_0(M),\Zl)$]\label{prop:congruence:integral-structure}
Let $\ell\nmid2M$. Pull-back along $Y_{N'}\to Y_0(M)$ maps $H^1(X_0(M),\Zl)$ (the interior cohomology of $Y_0(M)$) $\T$-equivariantly onto a sublattice of finite index of $M(M,1)_{!,\lambda}$. For
every maximal ideal $\mf m$ of $\T_\ell$ that is not Eisenstein (\Cref{ch:moddeg}), the localisations agree:
\[
M(M,1)_{!,\lambda,\mf m}=H^1(X_0(M),\Zl)_{\mf m}.
\]
In particular this holds for $\mf m=\mf m_{f,\lambda}$ when $\bar\rho_{f,\lambda}$ is irreducible, which is the only case the paper uses.
\end{proposition}

\begin{proof}[Sketch of proof]
The interior cohomology (the image of $H^1_c$ in $H^1$) of $Y_0(M)$ with $\Zl$-coefficients is $H^1(X_0(M),\Zl)$ (\Cref{ch:moddeg}); pull-back to the finite cover $Y_{N'}$ lands in the
$U_0(M)$-invariants and is compatible with the Hecke correspondences, and it is injective with finite cokernel after tensoring with $\Ql$ (the covering has finite degree, and
push-forward composed with pull-back is multiplication by the degree). By the inflation-restriction sequence for the covering group, the cokernel is a subquotient of a
group of homomorphisms from the covering group to $\Qlmod$, that is, it comes from abelian unramified subcovers of congruence covers of $X_0(M)$. On such classes the
Hecke operators $T_p$, $p\nmid MN'$, act through $1+p$ (more precisely through an Eisenstein ideal), as on the Shimura covering $X_1(M)\to X_0(M)$
\cite{LingOesterle,Ribet88}; so the cokernel is supported on Eisenstein maximal ideals and vanishes after localisation at a non-Eisenstein $\mf m$. When $\ell$ does not divide
the order of the covering group the cokernel is $0$ (transfer). Diamond, Flach and Guo use the same device (kernel and cokernel supported at Eisenstein ideals) in
\cite[Lemma 1.8, Corollary 1.9]{DFG04}, and arrange a covering of degree prime to $\ell$ with an auxiliary prime \cite[Section 1.8]{DFG04}. The Eisenstein property of the cokernel
is proved in the cited sources for the Shimura covering; for general congruence covers we have given only the sketch above.
\end{proof}

\begin{paperlink}
The paper (Section 4.1) says that the integral structure ``contains the pull-back of $H^1(X_0(M^\Sigma),\Zl)$ with finite index, and we use it only after localisation at a maximal ideal
$\mf m$ of the Hecke algebra whose residual representation is irreducible, where we take it, as in the version of 10 October 2026, to agree with $H^1(X_0(M^\Sigma),\Zl)_{\mf m}$''. That
is \Cref{prop:congruence:integral-structure}; the word ``take'' marks it, correctly, as the one input of Section 4.1 that the paper sketches rather than proves.
\end{paperlink}

\begin{definition}[the pairing $\hat\delta_!$]\label{def:congruence:dfg-pairing}
Diamond, Flach and Guo define $\hat\delta_!$ on $M(M,1)_!$ as the cup product into $H^2_c$ of the full-level curve, restricted to invariants and divided by the index $[U_0(M):U_{N'}]$, composed
with the operator $[UwU]_\omega$ attached to $w=\sm0{-1}M0$ \cite[Sections 1.5.2, 1.5.3, (13), (14)]{DFG04}; for weight $2$ and trivial character this operator is the Atkin-Lehner
involution $w_M$, with no scalar factor. On $H^1(X_0(M),\Zl)$ the book writes $\hat\delta_!(x,y)=\langle x,w_M^*y\rangle$, the cup product composed with $w_M^*$ (Appendix B, B-C13(f);
\Cref{ch:moddeg}).
\end{definition}

\begin{proposition}[properties of $\hat\delta_!$]\label{prop:congruence:dfg-pairing}
Let $\ell$ be odd.
\begin{enumerate}[label=(\roman*)]
\item On classes pulled back from $X_0(M)$, the pairing of Diamond, Flach and Guo is $c\langle x,w_M^*y\rangle$ with $c=\deg(Y_{N'}\to Y_0(M))/[U_0(M):U_{N'}]=\frac12$, a unit in $\Zl$.
\item $\langle x,w_M^*y\rangle$ is perfect and alternating on $H^1(X_0(M),\Zl)$, every $t\in\T$, the $U_p$ included, is self-adjoint for it, and its restriction to each $\mf m$-component is perfect.
\end{enumerate}
So on a non-Eisenstein component the two pairings agree up to a unit, and the ideals of values and Gram determinants of their restrictions to any sublattice have the same
valuations.
\end{proposition}

\begin{proof}
(i) The degree of a finite map of curves multiplies the fundamental class in $H^2$ by the degree, so the cup product of pulled-back classes is $\deg(Y_{N'}\to Y_0(M))$ times the
cup product on $Y_0(M)$. The group $U_0(M)/U_{N'}$ acts on $Y_{N'}$ over $Y_0(M)$ through its quotient by the image of $\pm1$, because $-1\in\GL_2(\Q)$ acts trivially on the upper and
lower half planes; so the degree is half the index, and $c=\frac12$. The division by the index is that of \cite[Section 1.5.2]{DFG04} (``normalizing by dividing by $[U:U_N]$''). In \cite[Section 1.5.3]{DFG04}, $w$ is the
element with entries $\sm0{-1}M0$ of $\prod_{p\mid M}\GL_2(\Z_p)$, and $[UwU]_\omega$ ``coincides with $\psi(-1_\infty)N^{k-2}[Uw^{-1}U]_\omega$'' (their $N$ is our $M$); for $k=2$ and $\psi=1$
there is no scalar, and $[UwU]_\omega$ is the involution induced by the global matrix $\sm0{-1}M0$, which lies in $w\prod_{p\nmid M}\GL_2(\Z_p)$, that is, by $w_M$. (ii) This is the corollary on $\hat\delta_!$ in \Cref{ch:moddeg}: the transpose of the correspondence $t$ is $w_Mtw_M$,
so $(t^*)^\dagger=w_M^*t^*w_M^*$ for the cup product, whence $\langle tx,w_M^*y\rangle=\langle x,w_M^*ty\rangle$; $w_M^*$ is an involution adjoint to itself, and the cup product is alternating and perfect
on a smooth projective curve (Poincar\'e duality, \Cref{ch:curves}), so $\langle x,w_M^*y\rangle$ is perfect, and alternating because $\langle x,w_M^*x\rangle=-\langle w_M^*x,x\rangle=-\langle w_M^*x,w_M^{*2}x\rangle$ $=-\langle x,w_M^*x\rangle$, as $w_M^*$
preserves the cup product, so $2\langle x,w_M^*x\rangle=0$. The statement on components is \Cref{prop:congruence:hecke-local}(iv).
\end{proof}

\begin{definition}[the kernel lattice and the naive $\Sigma$-finite congruence ideal]\label{def:congruence:eta-sigma}
Let $f$ be a newform of weight $2$ on $\Gamma_0(M)$ with rational coefficients, $\ell\nmid2M$, and $\Sigma$ a finite set of primes. For an eigenform $h$ of $\T_{M'}$ with eigencharacter $\theta_h$ put
\[
M_{h,\lambda}=M(M',1)_{!,\lambda}[\ker\theta_h],
\]
the submodule killed by the kernel of $\theta_h$ on the full Hecke algebra \cite[Sections 1.1.2, 1.6.2]{DFG04}. With $h=f^\Sigma$ at level $M^\Sigma$, the restriction of $\hat\delta_!$ to the rank-two lattice
$M_{f^\Sigma,\lambda}$ is alternating; its ideal of values, the image of $\wedge^2M_{f^\Sigma,\lambda}$ under $\hat\delta_!$, is an ideal of $\Zl$, written $\eta^\Sigma_{f,\lambda}$: the \emph{naive $\Sigma$-finite
congruence ideal} of $f$ \cite[Section 1.7.3]{DFG04}. For $\Sigma=\emptyset$, $f^\emptyset=f$ and $M_{f,\lambda}$ is the lattice of $f$ at level $M$.
\end{definition}

\begin{proposition}[properties of $M_{f^\Sigma,\lambda}$ and $\eta^\Sigma_{f,\lambda}$]\label{prop:congruence:eta-sigma}
In \Cref{def:congruence:eta-sigma} assume that $\bar\rho_f$ is irreducible, and let $\mf m$ be the maximal ideal of $\T_{M^\Sigma}\otimes\Zl$ attached to $f^\Sigma$.
\begin{enumerate}[label=(\roman*)]
\item $M_{f^\Sigma,\lambda}=H^1(X_0(M^\Sigma),\Zl)[\ker\theta_{f^\Sigma}]$; it is saturated, of rank $2$, stable under $G_\Q$, and contained in the $\mf m$-component.
\item The Gram determinant of $\hat\delta_!$ on $M_{f^\Sigma,\lambda}$ is $(\eta^\Sigma_{f,\lambda})^2$ up to a unit: if $a$ generates the ideal of values, the Gram matrix in a suitable basis is $\sm0a{-a}0$.
\item For $\Sigma=\emptyset$ and $f=f_E$ with $E$ optimal, $M_{f,\lambda}=\varphi_E^*H^1(E,\Zl)$, which is $\Hom(T_\ell E,\Zl)\simeq T_\ell E(-1)$ as a Galois module; hence $\ad M_{f,\lambda}\simeq\ad T_\ell E$.
\end{enumerate}
\end{proposition}

\begin{proof}
(i) $\theta=\theta_{f^\Sigma}$ factors through $\T_{\mf m}$, so $\theta(e_{\mf m})=1$ and $1-e_{\mf m}\in\ker\theta$; an element $x$ killed by $\ker\theta$ satisfies $x=e_{\mf m}x$. So $M_{f^\Sigma,\lambda}$ lies in the
$\mf m$-component, where \Cref{prop:congruence:integral-structure} identifies the integral structure with $H^1(X_0(M^\Sigma),\Zl)_{\mf m}$ ($\mf m$ is not Eisenstein: its residual
representation is $\bar\rho_f$, \Cref{ch:modgalrep}). A submodule killed by an ideal is saturated (Step 3 of the proof of \Cref{lem:congruence:ddt417}). By the Eichler-Shimura
isomorphism $H^1(X_0(M^\Sigma),\C)\simeq S_2\oplus\overline{S_2}$, which is Hecke-equivariant, the $\theta$-eigenspace is $\C f^\Sigma\oplus\C\overline{f^\Sigma}$, because the eigenspace of $\theta$ in $S_2$ lies in its
generalised eigenspace, the line $\C f^\Sigma$; so the rank is $2$. The Hecke correspondences are defined over $\Q$, so the \'etale realisation of $M_{f^\Sigma,\lambda}$ is $G_\Q$-stable
(\Cref{ch:moddeg}). (ii) An alternating form on $\Zl^2$ with basis $e_1,e_2$ has Gram matrix $\sm0a{-a}0$ with $a=\hat\delta_!(e_1,e_2)$, and its values $\hat\delta_!(xe_1+ye_2,ze_1+we_2)=(xw-yz)a$ fill $a\Zl$; the
determinant is $a^2$ (\Cref{ch:modules}). (iii) is the theorem ``$M_{f,\lambda}$ is the pull-back'' of \Cref{ch:moddeg} (the paper's Proposition prop:pullback), with the Weil pairing
and the comparison of \'etale and Betti cohomology.
\end{proof}

\begin{paperlink}
The paper's Section 4.2 wrote, in the version of 10 October 2026, that $M_{f,\lambda}$ and $T_\ell E$ are homothetic; the revised paper says that $\varphi_E^*$ identifies $M_{f,\lambda}$ with
$H^1(E,\Zl)\simeq T_\ell E(-1)$, the dual of $T_\ell E$, ``so its adjoint is $\ad T$''. \Cref{prop:congruence:eta-sigma}(iii) is that sentence; the adjoint is insensitive to dualising and twisting
(\Cref{ch:blochkato}), so nothing in the assembly depends on the difference.
\end{paperlink}

\begin{theorem}[Lemma lem:hecke, with Proposition prop:pullback]\label{thm:congruence:lem-hecke}
Let $E$ be $X_0(N)$-optimal, $f=f_E$, $\ell$ an odd prime with $\ell\nmid N$ and $E[\ell]$ irreducible. Then
\[
v_\ell(\degphi)=v_\ell(r_E)=v_\lambda(\eta^\emptyset_{f,\lambda}).
\]
No freeness or multiplicity-one hypothesis is needed. The first equality, and the fact that the ideal of values of $\langle x,w_N^*y\rangle$ on $\varphi_E^*H^1(E,\Zl)$ is
$\degphi\,\Zl$, are proved (given the cited theorem of Agashe, Ribet and Stein); the identification of that ideal with the ideal $\eta^\emptyset_{f,\lambda}$ of \cite{DFG04} uses the normalisation of $[UwU]_\omega$ (\Cref{def:congruence:dfg-pairing}, \Cref{prop:congruence:dfg-pairing}), the kernel lattice of
\cite[Sections 1.1.2, 1.6.2]{DFG04} (\Cref{def:congruence:eta-sigma}) and the definition of $\eta^\Sigma_f$ in \cite[Section 1.7.3]{DFG04}, all checked against the published text
(\Cref{rem:congruence:imported-table}), and the comparison of integral structures, which is sketched (\Cref{prop:congruence:integral-structure}).
\end{theorem}

\begin{proof}
The first equality is the theorem of Agashe, Ribet and Stein, \cite[Theorem 2.2]{ARS12}, since $\ell^2\nmid N$ (\Cref{ch:moddeg}). For the second: by the theorem ``the degree as a
congruence ideal'' of \Cref{ch:moddeg} (the paper's Proposition prop:pullback), for every $a,b\in H^1(E,\Zl)$,
$\langle\varphi_E^*a,w_N^*\varphi_E^*b\rangle_{X_0(N)}=\varepsilon_N(f)\,\degphi\,\langle a,b\rangle_E$, and the cup product on $H^1(E,\Zl)$ is unimodular, so the ideal of values of $\langle x,w_N^*y\rangle$ on
$\varphi_E^*H^1(E,\Zl)=M_{f,\lambda}$ is $\degphi\,\Zl$. By \Cref{prop:congruence:dfg-pairing}(i) and \Cref{prop:congruence:integral-structure}, the pairing of Diamond, Flach and Guo on their
lattice $M_{f,\lambda}$ (which is the same lattice, by \Cref{prop:congruence:eta-sigma}(i), the maximal ideal of $(f,\ell)$ being non-Eisenstein) differs from it by the unit $\frac12$. So
$\eta^\emptyset_{f,\lambda}=\degphi\,\Zl$.
\end{proof}

\begin{paperlink}
The revised paper's proof of Lemma lem:hecke is these two sentences. The version of 10 October 2026 proved only $v_\ell(r_E)\ge v_\lambda(\eta^\emptyset_{f,\lambda})$, by \Cref{lem:congruence:ddt417} applied
to $L=H^1(X_0(N),\Zl)_{\mf m}$ and $\pi=\phi_f$, with $d=2$ and $\Delta=(\eta^\emptyset_{f,\lambda})^2$, and equality when $L$ is free over $\T_{\mf m}$; the paragraph after the lemma in the revised paper
records this argument and the two points that make it correct: the congruence ideal is computed in $\T_{\mf m}$, the other local factors lying in the \emph{kernel} $\wp$ (not in its
annihilator, \Cref{prop:congruence:congruence-module}(iv)); and $\hat\delta_!$ is perfect on the $\mf m$-component (\Cref{prop:congruence:hecke-local}(iv)).
\end{paperlink}

\begin{remark}[Remark rem:free]\label{rem:congruence:rem-free}
For $\ell\nmid2N$ and $\mf m$ non-Eisenstein, $J_0(N)[\mf m]$ is two-dimensional over $\T/\mf m$ (Mazur, Ribet; \cite[Theorem 3.5]{RibetStein} for $J_1(N)$, transported to $J_0(N)$ through the
Shimura subgroup, which is Eisenstein); then $\T_{\mf m}$ is Gorenstein and $H^1(X_0(N),\Zl)_{\mf m}\simeq\T_{\mf m}^2$ \cite{Tilouine97}. Both steps are in \Cref{ch:moddeg} (multiplicity one stated,
Gorenstein and freeness proved given it), and the freeness criterion is in \Cref{ch:modules}. After \Cref{thm:congruence:lem-hecke} this is \emph{not} needed for Theorem thm:main, nor for any
step of Theorem thm:quant (\Cref{thm:congruence:quant}), nor for the equality clauses of either. It remains the conceptual reason why the Hecke and the cohomological
congruence ideals agree at $\ell\nmid2N$ (\Cref{lem:congruence:ddt417} with equality), and with \Cref{lem:congruence:ddt417} it gives a second proof of $v_\ell(r_E)=v_\ell(\degphi)$ for
$\ell\nmid2N$, $E[\ell]$ irreducible. Freeness fails in general when $\ell\mid N$: \curve{99a1} at $\ell=3$ (\Cref{ex:congruence:ddt417-checks}(d)).
\end{remark}

\begin{example}[values of $\eta^\emptyset$]\label{ex:congruence:eta-values}
By \Cref{thm:congruence:lem-hecke}, $v_\ell(\eta^\emptyset_{f,\ell})=v_\ell(\degphi)$ whenever $\ell\nmid2N$ and $E[\ell]$ is irreducible. At $\ell=3$ (all curves below have $E[3]$ irreducible and
$3\nmid N$; degrees from PARI/GP and Cremona's tables): \curve{11a1}, $\degphi=1$, $v_3=0$; \curve{77c1}, $6$, $v_3=1$; \curve{121a1}, $6$, $v_3=1$; \curve{121d1}, $24$, $v_3=1$; \curve{539c1},
$288$, $v_3=2$; \curve{539d1}, $72$, $v_3=2$. For the depleted forms of \Cref{ex:congruence:level-539}, \Cref{thm:congruence:dfg14} below gives
$v_3(\eta^{\{7\}}_{\text{77c},3})=v_3(\eta^\emptyset_{\text{77c},3})+v_3(1-7^{-2})=1+1=2$ and $v_3(\eta^{\{7\}}_{\text{11a},3})=0+v_3\bigl(\frac67\cdot\frac{64-4}{49}\bigr)=0+2=2$, equal to $v_3(288)$ and $v_3(504)$:
this is Theorem thm:quant read backwards (\Cref{thm:congruence:quant}).
\end{example}

\section{The variation of the congruence ideal with $\Sigma$}\label{sec:congruence:dfg14}

This section treats the paper's Theorem thm:dfg14, \cite[Proposition 1.4(c)]{DFG04}: adding a prime $p$ to $\Sigma$ multiplies the naive congruence ideal by the naive Euler
factor $L^{\nv}_p(A_f,1)^{-1}$. The paper uses it twice: in the proof of Theorem thm:main, to pass from level $N$ to $\Sigma=\set{p\mid N}$, and in Step 3 of the proof of Theorem
thm:quant, to pass from the newform $g=f_{E_S}$ at level $N_S$ to its depletion at level $N$. We state it, explain the three cases of the paragraph after it, and show
how part (c) follows from parts (a) and (b), following the published proof: homothety of lattices, and primitivity of $f^\Sigma$ read through the integral comparison isomorphism.

\begin{definition}[the map $\gamma$]\label{def:congruence:gamma}
Let $f$ be a newform of weight $2$ on $\Gamma_0(M)$ and $\Sigma$ a finite set of primes. For $m\mid M^\Sigma/M=\prod_{p\in\Sigma}p^{\delta_p}$ let $\epsilon_m\colon H^1(X_0(M))\to H^1(X_0(M^\Sigma))$ be the
pull-back along the degeneracy map $z\mapsto mz$ (on forms, $h\mapsto h(mz)$; \cite[Section 1.7.3]{DFG04} normalise $\epsilon_m$ by the factor $m^{1-k}=m^{-1}$, a unit when $\ell\nmid m$), and let $\phi_m$ be the endomorphism of the
cohomology at level $M$ which is multiplicative in coprime $m$ with $\phi_1=1$, $\phi_p=-T_p$ and $\phi_{p^2}=p$. Put
\[
\gamma=\sum_{m\mid M^\Sigma/M}\epsilon_m\phi_m,
\]
and let $\gamma^t$ be its adjoint for the pairings $\hat\delta_!$ at the two levels. On $q$-expansions $\gamma(f)=\prod_{p\in\Sigma,\,\delta_p=2}(1-a_pV_p+pV_p^2)\prod_{p\in\Sigma,\,\delta_p=1}(1-a_pV_p)\,f=f^\Sigma$,
with $V_p\colon h(z)\mapsto h(pz)$ (\Cref{ch:hecke}, depletion).
\end{definition}

\begin{theorem}[Diamond, Flach and Guo, Proposition 1.4; the paper's Theorem thm:dfg14]\label{thm:congruence:dfg14}
Let $f$ be a newform of weight $2$ on $\Gamma_0(M)$ with rational coefficients, $\ell$ a prime with $\ell\nmid2M$, and $\Sigma$ a finite set of primes with $\ell\notin\Sigma$.
\begin{enumerate}[label=(\alph*)]
\item $\gamma$ restricts to an isomorphism $M_f\otimes\Q\to M_{f^\Sigma}\otimes\Q$ of the rational eigenspaces, compatible with $f\mapsto f^\Sigma$.
\item On $M_f$, $\gamma^t\circ\gamma=u\cdot\prod_{p\in\Sigma}L^{\nv}_p(A_f,1)^{-1}$ with $u=\prod_{p\in\Sigma,\,\delta_p=1}(-a_p)\prod_{p\in\Sigma,\,\delta_p=2}p$, an $\ell$-adic unit; in particular $\hat\delta_!$ is non-degenerate on
$M_{f^\Sigma}$.
\item If $\bar\rho_f$ is irreducible, then $\gamma(M_{f,\lambda})=M_{f^\Sigma,\lambda}$ and
\[
\eta^\Sigma_{f,\lambda}=\eta^\emptyset_{f,\lambda}\prod_{p\in\Sigma}L^{\nv}_p(A_f,1)^{-1}.
\]
\end{enumerate}
\end{theorem}

\begin{proof}[Sketch of proof]
(a) The maps $\epsilon_m$ are injective on the $f$-isotypic part with images the $f(mz)$-parts of the old space (multiplicity one, \Cref{ch:hecke}), and $\gamma$ maps the $f$-part onto the
line of the depletion and its conjugate, by the $q$-expansion formula; it commutes with $G_\Q$ because the degeneracy maps are defined over $\Q$.

(b) This is a double-coset computation (Diamond, Flach and Guo refer to \cite[Chapter 2]{Wiles95} and to page 121 of the 1995 printing of \cite{DDT} in \emph{Current Developments in Mathematics}); we sketch it in the case of one good prime $p$. The
pairings $\hat\delta_!(\epsilon_ix,\epsilon_jy)$ for $i,j\in\set{1,p,p^2}$ are the cup products at level $M$ of $x$ with $y$ transformed by the composites $\epsilon_i^t\epsilon_j$, which are
Hecke operators at level $M$ ($\epsilon_1^t\epsilon_1$ is the degree $(p+1)p$ of the covering, $\epsilon_1^t\epsilon_p$ is $T_p$, and so on); on $M_f$ each acts by a scalar polynomial in $a_p$ and $p$, and
expanding $\gamma^t\gamma=\sum_{i,j}\phi_i^t\epsilon_i^t\epsilon_j\phi_j$ gives a scalar, which equals $p\,(p-1)((p+1)^2-a_p^2)/p^3$ up to the normalisations of the $\epsilon_m$, that is $u$ times Wiles's factor.
The multiplicative case is the $2\times2$ analogue. The non-degeneracy
follows since every factor is nonzero (\Cref{prop:congruence:three-cases}). We do not give the details; see \cite[Section 1.7.3]{DFG04}.

(c) The lattices $\gamma(M_{f,\lambda})$ and $M_{f^\Sigma,\lambda}$ are $G_\Q$-stable lattices in the same two-dimensional representation, by (a), and $\gamma(M_{f,\lambda})\subseteq M_{f^\Sigma,\lambda}$ because
$\gamma$ preserves integral cohomology ($p\ne\ell$ for $p\in\Sigma$, so the normalising factors of the $\epsilon_m$ are units) and maps the $f$-eigenspace into the $f^\Sigma$-eigenspace by (a), and
$M_{f^\Sigma,\lambda}$ is the intersection of that eigenspace with the integral cohomology. As $\bar\rho_f$ is irreducible, every
$G_\Q$-stable lattice in it is homothetic to a given one (\Cref{ch:modgalrep}, lattices), so $\gamma(M_{f,\lambda})=\ell^nM_{f^\Sigma,\lambda}$ for some $n\ge0$. Diamond, Flach and Guo
get $n=0$ from the de Rham side, not from a statement about the degeneracy maps modulo $\ell$; their proof reads: ``If $\bar{\mathcal M}_{f,\lambda}$ is irreducible, then the image of
$\mathcal M_{f,\lambda}$ must be of the form $\lambda^n\mathcal M_{f^\Sigma,\lambda}$ for some $n\geqslant0$; since $\gamma_{\mathrm{dR}}(f)=f^\Sigma$, we see that $n=0$.'' The input is the integral comparison
isomorphism that is part of their integral structures at $\lambda$ \cite[Sections 1.1.2, 1.2.5]{DFG04}, which we cite and do not prove (it is the theory of Fontaine and Laffaille \cite{FontaineLaffaille}
applied to the modular curve, which has good reduction at $\ell\nmid M^\Sigma$): $M_{f,\lambda}$ and $M_{f^\Sigma,\lambda}$ are the images under the fully faithful Fontaine-Laffaille functor of
filtered modules $D_f$ and $D_{f^\Sigma}$ whose step $\mathrm{Fil}^1$ is $\Zl f$, respectively $\Zl f^\Sigma$ (for a normalised eigenform the first filtration step of the de Rham realisation is spanned by
the form, \cite[(15)]{DFG04}), and $\gamma$ is a morphism of these structures. If $\gamma(M_{f,\lambda})=\ell^nM_{f^\Sigma,\lambda}$, full faithfulness gives an isomorphism $D_f\to\ell^nD_{f^\Sigma}$ of filtered
modules, which maps $\mathrm{Fil}^1D_f=\Zl f$ onto $\ell^n\Zl f^\Sigma$; as it maps $f$ to $\gamma(f)=f^\Sigma$, $n=0$. (Saturation statements of the type of Ihara's lemma for the degeneracy maps
modulo $\ell$ would give another route; the published proof does not use them, and neither does the book.) Now for $x,y\in M_{f,\lambda}$,
\[
\hat\delta_!(\gamma x,\gamma y)=\hat\delta_!(x,\gamma^t\gamma y)=u\prod_{p\in\Sigma}L^{\nv}_p(A_f,1)^{-1}\,\hat\delta_!(x,y),
\]
so the ideal of values on $M_{f^\Sigma,\lambda}=\gamma(M_{f,\lambda})$ is $\eta^\emptyset_{f,\lambda}$ times $\prod_pL^{\nv}_p(A_f,1)^{-1}$, $u$ being a unit ($a_p=\pm1$ at $\delta_p=1$, and $p\ne\ell$).
\end{proof}

\begin{remark}[status of the statement]\label{rem:congruence:dfg14-status}
The paper cites only part (c), as ``\cite[Proposition 1.4(c)]{DFG04}'', under the standing assumptions of its Section 4.1: $f$ with rational coefficients, $\lambda=\ell\nmid2M$, $\Sigma$ a finite set of
primes. It does not say $\ell\notin\Sigma$, and the hypothesis is needed: for $\ell\in\Sigma$ the statement of (c) is false. Take $f=f_{\text{11a}}$, $\ell=3$ (so $\ell\nmid2M=22$) and
$\Sigma=\set3$. Here $\delta_3=2$ and $a_3=-1$, so by \Cref{prop:congruence:three-cases}(ii)
\[
L^{\nv}_3(A_f,1)^{-1}=\frac23\cdot\frac{16-1}{9}=\frac{2\cdot15}{27}=\frac{10}9,\qquad v_3\Bigl(\frac{10}9\Bigr)=-2
\]
(recomputed with PARI/GP, also from the Euler factor $(1+2X-6X^2-27X^3)^{-1}$ of the symmetric square at $3$, evaluated at $X=3^{-2}$). As $\eta^\emptyset_{f,3}=\degphi\,\Z_3=\Z_3$
(\Cref{ex:congruence:eta-values}), the right side of (c) would be $3^{-2}\Z_3$, which is not contained in $\Z_3$; but however $\eta^{\{3\}}_{f,3}$ is read at the level $M^\Sigma=99$, which $3$ divides (for instance as the
ideal of values of $\langle x,w_{99}^*y\rangle$ on $H^1(X_0(99),\Z_3)[\ker\theta_{f^\Sigma}]$), it is an ideal of values of an integral pairing on an integral lattice, so it lies in $\Z_3$. Every use in the paper has $\ell\notin\Sigma$ ($\Sigma=\set{p\mid N}$ with $\ell\nmid N$, and $\Sigma=S$ with $\ell\nmid N$), so we state the theorem with $\ell\notin\Sigma$. In \cite{DFG04} the hypothesis is a standing assumption: in their Section 1.7.3, $\Sigma$ is a finite
subset of $S=S_{\mathbf f}(K)\setminus\set\lambda$, the finite primes of $K$ other than $\lambda$. Part (b) is as printed: \cite[Proposition 1.4(b)]{DFG04} reads
$\gamma^t\circ\gamma=\phi_{N^\Sigma/N_f}\prod_{p\in\Sigma}L^{\nv}_p(A_f,1)^{-1}$ on $M_f$, with $\phi_p=-T_p$ and $\phi_{p^2}=\psi(p)p^{k-1}$, so for weight $2$ and trivial character $\phi_{N^\Sigma/N_f}$ acts on $M_f$ by $u$;
the proof of their Theorem 2.15 records that this element, $\prod_{\delta_p=1}(-a_p)\prod_{\delta_p=2}\psi(p)p^{k-1}$, is a unit at $\lambda$. The book uses only (c).
\end{remark}

\begin{proposition}[the three cases]\label{prop:congruence:three-cases}
Let $f$ be a newform of weight $2$ on $\Gamma_0(M)$ with rational coefficients and $p$ a prime. The factor $L^{\nv}_p(A_f,1)^{-1}$ of \Cref{thm:congruence:dfg14} is:
\begin{enumerate}[label=(\roman*)]
\item $1$ if $\delta_p=0$ ($p^2\mid M$, additive for $f=f_E$): adding $p$ to $\Sigma$ changes neither the level nor the ideal;
\item if $\delta_p=2$ ($p\nmid M$, good), with $\alpha,\beta$ the roots of $X^2-a_pX+p$,
\[
\Bigl(1-\frac\alpha\beta p^{-1}\Bigr)\Bigl(1-p^{-1}\Bigr)\Bigl(1-\frac\beta\alpha p^{-1}\Bigr)=\frac{p-1}p\cdot\frac{(p+1)^2-a_p^2}{p^2},
\]
Wiles's level-raising factor, of $\ell$-adic valuation $v_\ell(p-1)+v_\ell((p+1)^2-a_p^2)$ for $\ell\ne p$;
\item $1-p^{-2}$ if $\delta_p=1$ ($p\,\|\,M$, multiplicative), of valuation $v_\ell(p^2-1)$.
\end{enumerate}
So only the primes with $\delta_p\ge1$ contribute.
\end{proposition}

\begin{proof}
(i) is the definition of the naive factor. (ii) The geometric Frobenius acts on $A_f=\ad V$ with eigenvalues $\alpha/\beta$, $1$, $\beta/\alpha$ (Appendix B, B-C3(d); this set does not depend on the choice of
Frobenius), so $L_p(A_f,1)^{-1}$ is the displayed product. Expanding the outer factors, $(1-\frac\alpha\beta p^{-1})(1-\frac\beta\alpha p^{-1})=1-\frac{\alpha^2+\beta^2}{\alpha\beta}p^{-1}+p^{-2}=1-\frac{a_p^2-2p}{p^2}+\frac1{p^2}=\frac{(p+1)^2-a_p^2}{p^2}$,
using $\alpha+\beta=a_p$, $\alpha\beta=p$. (iii) At a multiplicative prime $(\ad V)^{I_p}$ is a line on which the geometric Frobenius acts by $p^{-1}$ (\Cref{ch:lfunctions}), so $L_p(A_f,s)^{-1}=1-p^{-1-s}$, which is
$1-p^{-2}$ at $s=1$; equivalently $L^{\nv}_p(A_f,s)^{-1}=1-a_p^2p^{-1-s}$ with $a_p^2=1$. The valuations follow, $p$ being prime to $\ell$.
\end{proof}

\begin{remark}[what the factors count]\label{rem:congruence:factors-count}
At a good prime $p\ne\ell$, $v_\ell((p+1)^2-a_p^2)>0$ is Ribet's level-raising condition $a_p\equiv\pm(p+1)\pmod\ell$ for a newform of level $Mp$ congruent to $f$ (\Cref{ch:modgalrep}), and
$v_\ell(p-1)>0$ allows congruences with forms of level $Mp^2$ whose local component at $p$ is a ramified principal series of $\ell$-power order. The valuation of the factor is the
number of new congruences that $f^\Sigma$ acquires at the higher level. At a multiplicative prime the factor $1-p^{-2}$ is the analogue for raising $p\,\|\,M$ to $p^2$. The paper's
paragraph after Theorem thm:dfg14 (``adding a good prime $p$ multiplies the ideal by \dots\ Wiles's level-raising factor; adding a multiplicative prime multiplies it by $1-p^{-2}$'') is
\Cref{prop:congruence:three-cases}.
\end{remark}

\begin{remark}[Ihara's lemma and the Eisenstein case]\label{rem:congruence:ihara}
In the proof of part (c), as in \cite{DFG04}, irreducibility of $\bar\rho_f$ enters through the homothety of lattices: without it $\gamma(M_{f,\lambda})$ need not be a multiple of
$M_{f^\Sigma,\lambda}$, and the primitivity argument no longer fixes the index. (In the book it also enters through the comparison of integral structures at a non-Eisenstein $\mf m$,
\Cref{prop:congruence:integral-structure}.) Ihara's lemma, the injectivity of degeneracy maps modulo $\ell$ on non-Eisenstein components, is not used. In the Eisenstein case the factor relating $\eta^\Sigma$ and $\eta^\emptyset$ can differ from the Euler factor by a non-unit; the paper's Conjecture conj:eis and
Section sec:eis concern that case, which lies outside every theorem of the paper.
\end{remark}

\begin{remark}[refined integral structures]\label{rem:congruence:refined-structures}
Diamond, Flach and Guo also use refined integral structures \cite[Section 1.8]{DFG04}, which change the level at exceptional primes $p\equiv-1\pmod\ell$ where the local
representation is induced from the unramified quadratic extension, and add an auxiliary prime $r$ at which the Euler factor is a unit, so that the relevant cover has degree
prime to $\ell$ and the pairing is perfect integrally. With them the naive factors are replaced by the true ones; this is how the true Euler factors enter their proof of
the Tamagawa number conjecture. The paper never uses the refined ideals: it combines the naive ideal (\Cref{thm:congruence:dfg14}) with Lemma lem:dfg21, whose local terms carry the
true factors, and the difference between the two is exactly $\sum_qe_\ell(E,q)$ (Step 5 of \Cref{thm:congruence:main}).
\end{remark}

\section[The exceptional set and the theorems of DFG]{The exceptional set and the theorems of Diamond, Flach and Guo}\label{sec:congruence:dfg-theorems}

This section states the theorems of Diamond, Flach and Guo that the paper imports in its Section 4.2: Theorem thm:dfg37 (the length of the $\Sigma$-imprimitive
adjoint Selmer group is the valuation of the naive congruence ideal), Theorem thm:dfg27 (vanishing of global invariants and Selmer groups), and their main theorem on the
Tamagawa number conjecture, together with the exceptional set $S_f$ and the notion of minimal ramification that enter the hypotheses. It also records what has been checked
about the case $\ell=3$, which \Cref{ch:blochkato} left open. The proofs are by the Taylor-Wiles method, outlined in Section \ref{sec:congruence:taylor-wiles}.

\begin{definition}[the exceptional set]\label{def:congruence:exceptional-set}
Let $f$ be a newform of weight $k\ge2$ and level $N$ with coefficients in $K$. Diamond, Flach and Guo define $S_f$ as the set of finite primes $\lambda$ of $K$ such that either $\lambda\mid Nk!$, or the
residual representation $M_{f,\lambda}/\lambda M_{f,\lambda}$ is not absolutely irreducible when restricted to $G_F$, where $F=\Q(\sqrt{(-1)^{(\ell-1)/2}\ell})$ and $\lambda\mid\ell$ \cite[Section 0.2]{DFG04}; the
definition is repeated as the display \cite[(31)]{DFG04} at the start of their Section 2 (p.~695), which is the number the paper cites. For $k=2$ and $K=\Q$: $\ell\in S_f$ iff $\ell\mid2N$ or $\bar\rho|_{\Gal(\Qbar/\Q(\sqrt{\ell^*}))}$ is not absolutely irreducible.
\end{definition}

\begin{proposition}[the paper's Proposition prop:Sf]\label{prop:congruence:exceptional-set}
Let $E/\Q$ have conductor $N$, $f=f_E$, and $\ell$ an odd prime with $\ell\nmid N$. Then $\ell\notin S_f$ if and only if $E[\ell]$ is irreducible and, when $\ell=3$, the mod $3$ image is not conjugate to
$N_s(\F_3)$ (Sutherland's label $\Ns$); equivalently, $E[\ell]$ is irreducible and $h_\ell(E)=0$. In particular, for $\ell\ge5$ with $\ell\nmid N$, $\ell\notin S_f$ as soon as $E[\ell]$ is irreducible.
\end{proposition}

\begin{proof}
This is the paper's Proposition prop:Sf, proved in \Cref{ch:galrep} (with the fundamental characters of \Cref{ch:modgalrep} and the subgroups of $\GL_2(\F_3)$ of \Cref{ch:groups});
the equivalence with $h_\ell(E)=0$ is Lemma lem:h0 (\Cref{ch:galrep,ch:blochkato}). In outline: an irreducible $\bar\rho$ is absolutely irreducible (complex conjugation has eigenvalues $\pm1$), and its
restriction to $G_F$ fails to be absolutely irreducible iff $\bar\rho\simeq\bar\rho\otimes\chi_{\ell^*}$; restricted to inertia at $\ell$ this is impossible for $\ell\ge5$ (compare the fundamental
characters in the ordinary and supersingular cases), and for $\ell=3$ it happens exactly for the image $N_s(\F_3)$.
\end{proof}

\begin{definition}[the $\Sigma$-imprimitive Selmer group]\label{def:congruence:selmer-sigma}
Let $\Sigma$ be a finite set of primes not containing $\ell$ and $W=\ad T\otimes\Qlmod$. Following \cite[Section 2.1]{DFG04}, $H^1_\Sigma(\Q,W)$ is the subgroup of $H^1(\Q,W)$ of classes that are
unramified at every $p\notin\Sigma\cup\set{\ell}$ (that is, lie in $H^1_{\ur}(\Q_p,W)$), satisfy the Bloch-Kato condition $H^1_f(\Q_\ell,W)$ at $\ell$, and are unrestricted at $p\in\Sigma$ and at $\infty$ (\Cref{ch:blochkato}).
\end{definition}

If $\Sigma$ contains every prime $p\ne\ell$ at which $W$ is ramified, then at $p\notin\Sigma\cup\set{\ell}$ the module $W$ is unramified, $W^{I_p}=W$ is divisible, and $H^1_{\ur}(\Q_p,W)=H^1_f(\Q_p,W)$
(\Cref{ch:blochkato}); so $H^1_\Sigma(\Q,W)$ is also the group with the Bloch-Kato conditions outside $\Sigma$, and $H^1_f(\Q,W)\subseteq H^1_\Sigma(\Q,W)$. Diamond, Flach and Guo write the unramified condition
in general because $H^1_{\ur}(\Q_p,W)$ need not be divisible at a ramified prime (\Cref{ex:congruence:ex-hur}). At $\ell$ the condition is the Bloch-Kato one; for the finite modules $W_n=W[\lambda^n]$ they define it through Fontaine-Laffaille theory,
as the extensions that lie in the essential image of the Fontaine-Laffaille functor, and \cite[Proposition 2.2]{DFG04} shows that the direct limit over $n$ is the Bloch-Kato condition.

\begin{definition}[minimal ramification]\label{def:congruence:minimally-ramified}
Let $\rho_0\colon G_\Q\to\GL_2(k)$ be a residual representation and $p\ne\ell$. A lift $\rho$ of $\rho_0$ (to $\mathcal O$, or to an artinian $\mathcal O$-algebra $A$) is \emph{minimally ramified at $p$} if
$\rho(I_p)\simeq\rho_0(I_p)$ (reduction is an isomorphism on the images of inertia) when $\#\rho_0(I_p)\ne\ell$, and if the quotient $M/M^{I_p}$ of the module $M$ of $\rho$ by its inertia invariants is free of rank one over $A$
(for $\rho$ over $\mathcal O$: $\dim(M\otimes K)^{I_p}=1$) when $\#\rho_0(I_p)=\ell$ \cite[Section 3.1, p.~715, and p.~717]{DFG04}. (The case distinction of \cite{DFG04} is $\#\rho_0(I_p)=\ell$ or not, not
whether $\ell$ divides $\#\rho_0(I_p)$.) $M_{f,\lambda}$ is \emph{minimally ramified outside $\Sigma$} if it is so at every prime $p\notin\Sigma\cup\set{\ell}$.
\end{definition}

If $\Sigma$ contains every prime $p\ne\ell$ at which $\rho_f$ is ramified, then $M_{f,\lambda}$ is minimally ramified outside $\Sigma$: at $p\notin\Sigma\cup\set{\ell}$, $\rho_f(I_p)=1=\rho_0(I_p)$. For $f=f_E$ and
$\Sigma=\set{p\mid N}$, $\rho_{E,\ell}$ is unramified outside $N\ell$ (\Cref{ch:galrep}), so the paper's choice qualifies. Diamond, Flach and Guo also use smaller sets $\Sigma$ internally (the minimal
deformation problem at the twist-minimal level); the paper never needs them.

\begin{theorem}[Diamond, Flach and Guo, Theorem 3.7; the paper's Theorem thm:dfg37]\label{thm:congruence:dfg37}
Let $f$ be a newform of weight $k\ge2$ and level $N$ with coefficients in $K$, $\lambda\notin S_f$, and $\Sigma$ a finite set of primes not containing $\ell$ such that $M_{f,\lambda}$ is minimally
ramified outside $\Sigma$. Then $H^1_\Sigma(\Q,A_{f,\lambda}\otimes K_\lambda/\mathcal O_\lambda)$ is finite, of length $v_\lambda(\eta^\Sigma_{f,\lambda})$. For $f=f_E$, $\ell\notin S_f$ and $\Sigma=\set{p\mid N}$:
\[
\len H^1_\Sigma(\Q,W)=v_\ell(\eta^\Sigma_{f,\ell}),\qquad W=\ad T_\ell E\otimes\Qlmod.
\]
\end{theorem}

\begin{proof}[Reference]
We do not prove this; it is \cite[Theorem 3.7]{DFG04} (p.~723; checked: the hypotheses are $\lambda\notin S_f$ and $M_{f,\lambda}$ minimally ramified outside $\Sigma$, $\ell\notin\Sigma$, and
$\eta^\Sigma_f$ is the ideal ``defined before Proposition 1.4''); see \cite[Section 3]{DFG04} and the outline in Section \ref{sec:congruence:taylor-wiles}. For $f=f_E$ the module $A_{f,\lambda}=\ad M_{f,\lambda}$ is $\ad T_\ell E$ by
\Cref{prop:congruence:eta-sigma}(iii), and $\Sigma=\set{p\mid N}$ satisfies the hypothesis by the remark after \Cref{def:congruence:minimally-ramified}.
\end{proof}

\begin{theorem}[Diamond, Flach and Guo, Theorem 0.2]\label{thm:congruence:dfg02}
The introduction of \cite{DFG04} states the theorem in the form the paper uses: ``Let $f$ be a newform of weight $k\ge2$ and level $N$ with coefficients in $K$. Suppose $\Sigma$ is a finite
set of rational primes containing those dividing $N$. Suppose that $\lambda$ is a prime of $K$ which is not in $S_f$ and does not divide any prime in $\Sigma$. Then the $\mathcal O_{K,\lambda}$-module
$H^1_\Sigma(G_\Q,A_{f,\lambda}/\mathcal A_{f,\lambda})$ has length $v_\lambda(\eta^\Sigma_f)$.'' It is labelled there ``Theorem 0.2 (= Theorem 3.7)''. Here $A_{f,\lambda}/\mathcal A_{f,\lambda}$ is the divisible module
$A_{f,\lambda}\otimes K_\lambda/\mathcal O_\lambda$, that is $W$ for $f=f_E$.
\end{theorem}

\begin{proof}[Reference]
We do not prove this; it is the case of \Cref{thm:congruence:dfg37} in which $\Sigma$ contains the primes dividing $N$, which suffices for the paper (\Cref{ch:blochkato} states it as a pointer).
\end{proof}

\begin{theorem}[Diamond, Flach and Guo, Theorem 2.7; the paper's Theorem thm:dfg27]\label{thm:congruence:dfg27}
If $\ell\notin S_f$, then $H^0(\Q,V^D)=H^1_f(\Q,V^D)=0$, and the same holds for $A_f=\ad V$: $H^0(\Q,A_f)=H^1_f(\Q,A_f)=0$. Here $V^D=A_f(1)$.
\end{theorem}

\begin{proof}[Reference]
We do not prove this; it is \cite[Theorem 2.7]{DFG04} (p.~699, checked), stated there for $A_f$ and $B_f=A_f(1)$ together with $\operatorname{ord}_{s=0}L(A_f,s)=\operatorname{ord}_{s=0}L(B_f,s)=0$.
\Cref{ch:blochkato} states it and sketches the derivation, which is the published proof: the $H^0$ vanish by absolute irreducibility, $\dim H^1_f(\Q,A_f)=\dim H^1_f(\Q,B_f)$ by the Euler characteristic
formula of Fontaine and Perrin-Riou, and $H^1_f(\Q,A_f)=0$ because the divisible part of $H^1_f(\Q,W)$ lies in the finite group $H^1_\Sigma(\Q,W)$ of \Cref{thm:congruence:dfg37}.
Only $H^1_f(\Q,V^D)=0$ is needed in the paper; $H^0(\Q,V^D)=0$ holds for every $E$ (\Cref{ch:blochkato}).
\end{proof}

\begin{theorem}[Diamond, Flach and Guo, Theorem 0.1]\label{thm:congruence:dfg01}
``Let $f$ be a newform of weight $k\ge2$ and level $N$ with coefficients in $K$. If $\lambda$ is not in $S_f$, then the $\lambda$-part of the Bloch-Kato conjecture holds for $A_f$ and $A_f(1)$.''
\cite[Theorem 0.1]{DFG04}, which is labelled ``(= Theorem 2.15)'' there.
\end{theorem}

\begin{proof}[Reference]
We do not prove this; see \cite[Sections 2.3, 2.4]{DFG04}. Its proof (that of their Theorem 2.15, pp.~709 to 711) takes for $\Sigma$ a nonempty set of primes $p\ne\ell$ containing those
where $M_\lambda$ is ramified, and combines \Cref{thm:congruence:dfg14} (with the unit $u$), the comparison of $H^1_\Sigma$ with $H^1_f$ (their Lemma 2.1, the paper's Lemma lem:dfg21), the local formula of
Fontaine and Perrin-Riou (the paper's Lemma lem:fpr) with \cite[(57)]{DFG04}, the triviality of the Tamagawa factor at $\ell$ \cite[Proposition 2.16]{DFG04} and \Cref{thm:congruence:dfg37}, with the
Petersson norm (\cite[Lemma 2.12]{DFG04}) and Deligne's period \cite[Theorem 2.10]{DFG04} in place of the modular degree:
it is the same assembly as the proof of Theorem thm:main (\Cref{thm:congruence:main}). The statement of the conjecture is in \Cref{ch:blochkato}.
\end{proof}

\begin{remark}[does Diamond, Flach and Guo's theorem cover $\ell=3$ in weight $2$?]\label{rem:congruence:ell-three}
\Cref{ch:blochkato} noted that the adjoint has Hodge-Tate weights $\set{-1,0,1}$, so that Fontaine-Laffaille theory cannot be applied to $\ad T$ directly when $\ell=3$, and left open
how Diamond, Flach and Guo treat that case. What has been checked against the published text is this. The abstract of \cite{DFG04} says that $\lambda$ ``is any prime of $K$ not dividing
$Nk!$'' such that the residual representation is absolutely irreducible over $\Q(\sqrt{(-1)^{(\ell-1)/2}\ell})$; the definition of $S_f$ in the introduction excludes only $\lambda\mid Nk!$ and the
non-absolutely-irreducible case; and Theorems 0.1 and 0.2, quoted above, are stated for every $\lambda\notin S_f$. For $k=2$, $Nk!=2N$, so as \emph{statements} these theorems include $\ell=3$ whenever
$3\nmid N$ and $\bar\rho|_{\Gal(\Qbar/\Q(\sqrt{-3}))}$ is absolutely irreducible; this is what the paper means by ``It covers $\ell=3$ in weight $2$: $S_f$ excludes only the primes dividing $2N$ and the induced case''.
How the proof treats $\ell=3$ has also been checked against the published text.
\begin{enumerate}[label=(\arabic*)]
\item Fontaine-Laffaille theory is applied to $M_{f,\lambda}$, not to the adjoint. A crystalline representation is \emph{short} if $\mathrm{Fil}^0D=D$, $\mathrm{Fil}^\ell D=0$ and no nonzero quotient
$V'$ has $V'\otimes\Ql(\ell-1)$ unramified \cite[Section 1.1.2]{DFG04}, and $M_{f,\lambda\text{-crys}}$ is an object of the relevant category ``if $\ell\nmid N$ and $\ell>k$'' (p.~713), which holds for
$k=2$, $\ell=3$. Section 1.7.1 of \cite{DFG04} records the obstruction of \Cref{ch:blochkato}: the adjoint gives such an object only after twisting ``if $\ell-1>2(k-1)$'', which fails for $k=2$, $\ell=3$.
\item The local condition at $\ell$ for $W=\Hom(V_1,V_2)\otimes\Qlmod$ with $V_1,V_2$ short is defined at finite level by extensions in the essential image of the Fontaine-Laffaille functor and
compared with the Bloch-Kato condition in \cite[Proposition 2.2]{DFG04}; the dimension count of \cite[Corollary 2.3]{DFG04} needs only $\ell>2$.
\item The refined integral structures assume that $\bar\rho$ is irreducible and, ``if $3\in\lambda$'', absolutely irreducible on $G_{\Q(\mu_3)}$ \cite[Section 1.8]{DFG04}; this is the condition $3\notin S_f$.
\item The Taylor-Wiles axioms ask for absolute irreducibility over the quadratic subfield of $\Q(\mu_\ell)$ and Serre weight $2\le k\le\ell-1$ \cite[Section 3.1]{DFG04}, which holds for $k=2$, $\ell=3$.
\item The case $\ell=k+1$ needs one extra input, in the triviality of the Tamagawa factor at $\ell$ for $A_f(1)$ \cite[Proposition 2.16]{DFG04}: when $a_\ell\equiv0\pmod\lambda$ their Lemma 2.17 suffices;
in the ordinary case ($a_\ell\not\equiv0\pmod\lambda$ and $\ell=k+1$) they appeal to Perrin-Riou's results on her conjecture $C_{EP}(V)$ for ordinary representations, with Colmez's proof of the
reciprocity law $\mathrm{Rec}(V)$. This enters only Theorem 0.1, not Theorems 2.7 and 3.7, which are what Theorem thm:main uses.
\end{enumerate}
In the same direction, \cite[Lemma 2.5]{DFG04} (an irreducible residual representation is absolutely irreducible over $F$) assumes $\lambda\nmid N(2k-1)(2k-3)k!$, which for $k=2$ excludes
$\ell=3$: this is the room for the image $\Ns$. The book therefore treats Theorem thm:main at $\ell=3$ (with $3\nmid N$ and image not $\Ns$) as a theorem.
\end{remark}

\section{The Taylor-Wiles method in outline}\label{sec:congruence:taylor-wiles}

This section explains, without proofs, why the length of a Selmer group should equal the valuation of a congruence ideal, which is the content of
\Cref{thm:congruence:dfg37}. The reader's natural question, ``why should a Selmer group know about congruences?'', has a short answer: the Selmer group is the dual of
the cotangent space of a deformation ring, the congruence ideal is the congruence ideal of a Hecke algebra, the two rings are isomorphic, and Wiles's numerical
criterion converts the isomorphism into an equality of lengths. The paper does not use the method beyond \Cref{thm:congruence:dfg37}, but its hypotheses are exactly those of
the method, and Section 1.5 of the paper contrasts it with the ``level lowering and twisting'' of Theorems thm:qual and thm:quant. The standard reference is \cite{DDT};
\cite{Wiles95,TaylorWiles,Diamond96} are the original papers.

\begin{definition}[deformation rings]\label{def:congruence:deformation-ring}
Let $k$ be a finite field of characteristic $\ell$, $\mathcal O$ the ring of integers of a finite extension of $\Ql$ with residue field $k$, and $\rho_0\colon G_\Q\to\GL_2(k)$ continuous, odd and absolutely
irreducible. For a finite set $\Sigma$ of primes not containing $\ell$, a \emph{deformation of type $\Sigma$} of $\rho_0$ to a complete local noetherian $\mathcal O$-algebra $A$ with residue field $k$ is a strict
equivalence class of lifts $\rho\colon G_\Q\to\GL_2(A)$ that are unramified outside $\Sigma\cup\set{\ell}\cup\set{\text{primes where }\rho_0\text{ ramifies}}$, minimally ramified (\Cref{def:congruence:minimally-ramified}) at the
ramified primes outside $\Sigma\cup\set\ell$, unrestricted at $\Sigma$, of fixed determinant $\chi_\ell$, and of Fontaine-Laffaille (``short crystalline'') type at $\ell$.
\end{definition}

\begin{theorem}[Mazur; representability and the tangent space]\label{thm:congruence:deformation-ring}
In the situation of \Cref{def:congruence:deformation-ring} the functor of deformations of type $\Sigma$ is represented by a complete local noetherian $\mathcal O$-algebra $R_\Sigma$. If $\pi\colon R_\Sigma\onto\mathcal O$
corresponds to a lift $\rho$ with module $M_\rho$, there is a natural isomorphism
\[
\Hom_{\mathcal O}\bigl(\mathrm{Cot}(\pi),K/\mathcal O\bigr)\simeq H^1_\Sigma\bigl(\Q,\ad M_\rho\otimes K/\mathcal O\bigr),
\]
the Selmer group with the local conditions of type $\Sigma$.
\end{theorem}

\begin{proof}[Sketch of proof]
Representability follows from Schlessinger's criteria, using absolute irreducibility (so that the deformations have only scalar automorphisms) and the finiteness of
$H^1(G_{\Q,S},\ad\rho_0)$. Infinitesimal deformations of $\rho$ to the ring $(\mathcal O/\varpi^n)[\varepsilon]$, $\varepsilon^2=0$, correspond to cohomology classes in $H^1(\Q,\ad M_\rho/\varpi^n)$, a lift $(1+\varepsilon c)\rho$ being determined by a
cocycle $c$; the conditions of type $\Sigma$ become the local conditions defining $H^1_\Sigma$, and passing to the limit over $n$ gives the isomorphism. We do not give the details; see
\cite[Sections 2.2 to 2.8]{DDT}.
\end{proof}

\begin{definition}[the Hecke side]\label{def:congruence:hecke-side}
With $\rho_0=\bar\rho_{f,\lambda}$ and $\Sigma$ as above, let $\T_\Sigma$ be the localisation at the maximal ideal of $\rho_0$ of the Hecke algebra acting on the cohomology $L_\Sigma$ of the modular curve
of level $N(\rho_0)\prod_{p\in\Sigma}p^{\delta_p}$ (with the minimal conductor $N(\rho_0)$ at the twist-minimal level), and let $\pi\colon\T_\Sigma\onto\mathcal O$ be the eigencharacter of the depletion of $f$.
\end{definition}

The Galois representation with coefficients in $\T_\Sigma$ built from the cohomology (Carayol) is a deformation of type $\Sigma$, so it induces $R_\Sigma\to\T_\Sigma$, which is surjective because $\T_\Sigma$ is
generated by traces of Frobenius elements. The point $\pi$ of $\T_\Sigma$ pulls back to a point of $R_\Sigma$.

\begin{theorem}[Taylor-Wiles, in the axiomatic form of Diamond, Flach and Guo]\label{thm:congruence:tw-axiomatic}
Let $\rho_0$ be odd, absolutely irreducible on $G_F$ with $F=\Q(\sqrt{(-1)^{(\ell-1)/2}\ell})$, of minimal conductor among its twists, with Serre weight $k$, $2\le k\le\ell-1$. Then for every $\Sigma$ the map
$R_\Sigma\to\T_\Sigma$ is an isomorphism of complete intersections, $L_\Sigma$ is free over $\T_\Sigma$, and for every lift $\rho$ of type $\Sigma$ with values in $\mathcal O$ (all of which are modular),
\[
\len_{\mathcal O}H^1_\Sigma(\Q,\ad M_\rho\otimes K/\mathcal O)=v\bigl(\eta(\T_\Sigma,\pi)\bigr)=v\bigl(\eta^\Sigma\bigr),
\]
with $\eta^\Sigma$ the ideal of values of the pairing on $L_\Sigma[\ker\pi]$.
\end{theorem}

\begin{proof}[Sketch of proof]
For $\Sigma=\emptyset$: choose, for each $n$, a set $Q_n$ of auxiliary primes $q\equiv1\pmod{\ell^n}$ at which $\rho_0$ is unramified with $\rho_0(\Frob_q)$ having distinct eigenvalues, of a fixed
cardinality $r=\dim H^1_\emptyset(\Q,\ad\rho_0)$, such that the dual Selmer group with the relaxed conditions at $Q_n$ vanishes. Allowing ramification at $Q_n$ makes $R_{Q_n}$ a quotient of a
power series ring in $r$ variables and $\T_{Q_n}$ free over the group ring of an $\ell$-group $\Delta_{Q_n}$ of rank $r$; patching over $n$ produces a power series ring mapping onto a module free
over a power series ring in the same number of variables, whence $R_\emptyset\simeq\T_\emptyset$ is a complete intersection and the cohomology is free. For general $\Sigma$, both sides change by
the same local factors as $\Sigma$ grows: on the Selmer side by the comparison of local conditions (the paper's Lemma lem:dfg21 and \Cref{ch:galcoh}), on the Hecke side by
\Cref{thm:congruence:dfg14}; with the numerical criterion (\Cref{thm:congruence:numerical-criterion}(iii)) this gives the equality at every $\Sigma$. Diamond, Flach and Guo formulate this
as an axiomatic theorem \cite[Section 3]{DFG04}; see \cite{DDT} (Chapters 3 to 5) for the case of semistable elliptic curves, and \cite{Diamond96} for the freeness and the role of the
cohomology. We do not give the details.
\end{proof}

\begin{remark}[where each part of the hypothesis $\ell\notin S_f$ is used]\label{rem:congruence:where-hypothesis}
(a) $\ell\nmid N$: the representation is crystalline at $\ell$ with Hodge-Tate weights $0,1$, Fontaine-Laffaille theory applies to $M_{f,\lambda}$, and the levels are prime to $\ell$.
(b) $\ell\nmid k!$, that is $\ell\ne2$ for $k=2$: the weights $0,1$ must lie in an interval of length at most $\ell-2$.
(c) Absolute irreducibility over $\Q(\sqrt{\ell^*})$: the auxiliary primes $q\in Q_n$ are found by Chebotarev in $\Gal(\overline\Q/\Q(\mu_{\ell^n}))$, and one needs elements $g$ there whose image $\rho_0(g)$ has
distinct eigenvalues and which detect any given class of the dual Selmer group. When $\rho_0$ is induced from $\Q(\sqrt{\ell^*})\subseteq\Q(\mu_\ell)$, its restriction to $G_{\Q(\mu_\ell)}$ is reducible, and the argument
fails. For elliptic curves with $\ell\nmid2N$ this bites only at $\ell=3$ with image $\Ns$ (\Cref{prop:congruence:exceptional-set}).
The paper's Section 1.5 and Section sec:anomaly turn on (c): Theorems thm:qual and thm:quant avoid it, and the curves with image $\Ns$ are where the identity needs a correction
(Section \ref{sec:congruence:beyond}).
\end{remark}

\begin{remark}[why the length equals the valuation]\label{rem:congruence:why-length-equals-eta}
In three lines, for $\ell\notin S_f$:
(1) $H^1_\Sigma(\Q,W)^*\simeq\mathrm{Cot}(\pi_{R_\Sigma})$ (\Cref{thm:congruence:deformation-ring});
(2) $R_\Sigma\simeq\T_\Sigma$ is a complete intersection, so $\len\mathrm{Cot}(\pi)=v(\eta(\T_\Sigma,\pi))$ (\Cref{thm:congruence:tw-axiomatic}, \Cref{thm:congruence:numerical-criterion}(ii));
(3) $L_\Sigma$ is free over $\T_\Sigma$, so $\eta(\T_\Sigma,\pi)=\eta^\Sigma$ (\Cref{lem:congruence:ddt417}, equality case, with $d=2$).
So the Selmer group ``knows'' the congruences because its dual is the cotangent space of a ring that is isomorphic to the Hecke algebra, and in a complete intersection the
cotangent space has exactly the size of the congruence module. The toy algebras of \Cref{ex:congruence:toy} show what would go wrong otherwise.
\end{remark}

\section{The assembly: proof of Theorem thm:main}\label{sec:congruence:assembly}

This section proves the paper's main theorem. The proof is an assembly: every input is a theorem proved or stated elsewhere, and the work here is to check that the
hypotheses match and that the local terms combine as claimed. The paper's Section 4.3 gives the proof in one paragraph; we give it in six numbered steps, each naming its
input and the hypothesis it uses. The result is an equality, with no multiplicity-one hypothesis, because Lemma lem:hecke is an equality (\Cref{thm:congruence:lem-hecke}).
From now on $f=f_E$ has rational coefficients and $\lambda=\ell$.

\begin{theorem}[the paper's Theorem thm:main]\label{thm:congruence:main}
Let $E/\Q$ be $X_0(N)$-optimal, $f=f_E$, $T=T_\ell E$, $W=\ad T\otimes\Qlmod$, and let $\ell$ be an odd prime with $\ell\nmid N$ and $\ell\notin S_f$ (equivalently, by \Cref{prop:congruence:exceptional-set}, $E[\ell]$
irreducible and, if $\ell=3$, mod $3$ image not $\Ns$). Then
\begin{equation}\label{eq:congruence:identity}
v_\ell(\degphi)=\len_{\Zl}H^1_f(\Q,W)+\sum_{q\mid N}v_\ell\bigl(\Tam_q(\ad T)\bigr)+\sum_{q\mid N}e_\ell(E,q),
\end{equation}
which is the paper's (eq:identity), and consequently
\begin{equation}\label{eq:congruence:thm}
v_\ell(\degphi)\ge\sum_{q\mid N}e_\ell(E,q)+\sum_{q\mid N}t_\ell(E,q),
\end{equation}
which is the paper's (eq:thm).
\end{theorem}

\begin{proof}
Let $\Sigma$ be the set of primes dividing $N$; it is nonempty (no elliptic curve over $\Q$ has good reduction everywhere, \Cref{ch:reduction}) and $\ell\notin\Sigma$.

\emph{Step 1 (the degree is the congruence ideal; uses optimality, $\ell\nmid2N$, $E[\ell]$ irreducible).} By \Cref{thm:congruence:lem-hecke},
\[
v_\ell(\degphi)=v_\ell(\eta^\emptyset_{f,\ell}).
\]

\emph{Step 2 (vary $\Sigma$; uses $\bar\rho$ irreducible, $\ell\nmid2N$, $\ell\notin\Sigma$).} By \Cref{thm:congruence:dfg14}(c),
\[
v_\ell(\eta^\emptyset_{f,\ell})=v_\ell(\eta^\Sigma_{f,\ell})-\sum_{q\in\Sigma}v_\ell\bigl(L^{\nv}_q(A_f,1)^{-1}\bigr).
\]
By \Cref{prop:congruence:three-cases} only the multiplicative primes contribute, each with $v_\ell(q^2-1)$.

\emph{Step 3 (Taylor-Wiles; uses $\ell\notin S_f$).} $M_{f,\ell}$ is minimally ramified outside $\Sigma$ (the remark after \Cref{def:congruence:minimally-ramified}) and $\ad M_{f,\ell}\simeq\ad T$
(\Cref{prop:congruence:eta-sigma}(iii)), so \Cref{thm:congruence:dfg37} gives
\[
v_\ell(\eta^\Sigma_{f,\ell})=\len H^1_\Sigma(\Q,W).
\]

\emph{Step 4 (the Selmer half; uses $\ell\notin S_f$, $\ell$ odd).} By the Selmer identity of \Cref{ch:blochkato} (which packages the paper's Lemma lem:dfg21 in the form proved there, Lemma
lem:fpr at every $q\in\Sigma$, Lemma lem:h0 in the form $h_\ell(E)=0$ for $\ell\notin S_f$, and Theorem thm:dfg27 for the hypothesis $H^1_f(\Q,V^D)=0$, the hypothesis $H^0(\Q,V^D)=0$ being automatic),
\[
\len H^1_\Sigma(\Q,W)=\len H^1_f(\Q,W)+\sum_{q\in\Sigma}\Bigl(v_\ell\bigl(L_q(A_f,1)^{-1}\bigr)+v_\ell\bigl(\Tam_q(\ad T)\bigr)\Bigr).
\]
The term at $\infty$ vanishes because $H^1(\R,T^D)$ is killed by $2$ and $\ell$ is odd, and the term $H^0(\Q,W^D)^*$ of Lemma lem:dfg21 vanishes because $h_\ell(E)=0$.

\emph{Step 5 (cancellation of the naive factors; Definition def:terms).} Combining Steps 1 to 4,
\[
v_\ell(\degphi)=\len H^1_f(\Q,W)+\sum_{q\mid N}v_\ell\bigl(\Tam_q(\ad T)\bigr)+\sum_{q\mid N}v_\ell\Bigl(\frac{L_q(A_f,1)^{-1}}{L^{\nv}_q(A_f,1)^{-1}}\Bigr).
\]
At a multiplicative $q$ the true and naive factors coincide ($\delta_q=1>0$) and the term is $0$; at an additive $q$ the naive factor is $1$ and the term is $v_\ell(L_q(A_f,1)^{-1})$. In both cases the
term is $e_\ell(E,q)$ by Definition def:terms (\Cref{ch:lfunctions}). This is \eqref{eq:congruence:identity}.

\emph{Step 6 (the inequality).} $\len H^1_f(\Q,W)\ge0$. At every prime $q\mid N$ with $v_q(j_E)<0$ (multiplicative, or additive of type $\kod{I}_n^*$, $q=2$ included),
$v_\ell(\Tam_q(\ad T))=t_\ell(E,q)+\min(t_\ell(E,q),v_\ell(q-1),v_\ell(c(u_q)))\ge t_\ell(E,q)$, by Lemma lem:tamfactor and its extension to potentially multiplicative additive primes
(\Cref{ch:blochkato}; at an additive prime one applies Lemma lem:tamfactor to the local quadratic twist that is multiplicative, $\ad T$ being unchanged, with $n=-v_q(j_E)$, not the
index of the Kodaira symbol, which differs at $q=2$). At the other bad primes $t_\ell(E,q)=0$ and $\Tam_q(\ad T)=\Zl$ (Lemma lem:tamadd, \Cref{ch:blochkato}). This gives \eqref{eq:congruence:thm}.
\end{proof}

\begin{corollary}[the paper's Corollary cor:explicit]\label{cor:congruence:explicit}
Under the hypotheses of \Cref{thm:congruence:main},
\begin{multline*}
v_\ell(\degphi)=\len H^1_f(\Q,W)+\sum_{q\mid N}e_\ell(E,q)\\
+\sum_{q:\,v_q(j_E)<0}\Bigl(t_\ell(E,q)+\min\bigl(t_\ell(E,q),v_\ell(q-1),v_\ell(c(u_q))\bigr)\Bigr),
\end{multline*}
with $u_q$ the unit part of the Tate parameter at $q$ and $c(u_q)$ the class of its reduction (Appendix B, B-C14). So $\len H^1_f(\Q,W)$ is computable from $\degphi$ and local data.
\end{corollary}

\begin{proof}
Substitute the values of the Tamagawa factors of \Cref{ch:blochkato} (Step 6 above) into \eqref{eq:congruence:identity}.
\end{proof}

\begin{paperlink}
The paper's proof of Theorem thm:main (Section 4.3), sentence by sentence. ``By Lemma lem:hecke, $v_\ell(\degphi)=v_\lambda(\eta^\emptyset_{f,\lambda})$'': Step 1. ``By Theorem thm:dfg14, \dots'': Step 2.
``By Theorem thm:dfg37, $v_\lambda(\eta^\Sigma_{f,\lambda})=\len H^1_\Sigma(\Q,W)$'': Step 3. ``By Theorem thm:dfg27 the hypotheses of Lemma lem:dfg21 hold, and $H^0(\Q,W^D)=0$ by Lemma lem:h0'', with
the display and ``by Lemma lem:fpr each local term is \dots'': Step 4. The display of $v_\ell(\degphi)$ and ``the last sum is $\sum_{p\mid N}e_\ell(E,p)$ by Definition def:terms'': Step 5. ``For
(eq:thm), the Selmer length is nonnegative, $v_\ell(\Tam_p(\ad T))\ge t_\ell(E,p)$ at every potentially multiplicative prime $p$, multiplicative or additive, $p=2$ included, by Lemma
lem:tamfactor'': Step 6. The version of 10 October 2026 stated (eq:identity) as an inequality with equality under freeness, and argued at additive potentially multiplicative primes
through the global twist $E_{\{p\}}$ of Lemma lem:twist, which exists only for odd $p$; the revised paper, like Step 6, argues locally.
\end{paperlink}

\begin{remark}[the hypotheses, and what fails without each]\label{rem:congruence:hypothesis-table}
\begin{center}\small
\begin{tabular}{p{0.24\textwidth}p{0.2\textwidth}p{0.46\textwidth}}
\toprule
hypothesis & used in & without it\\
\midrule
$E$ optimal & Step 1 & for another curve of the class the parametrisation has degree $\deg\psi\cdot\degphi$; harmless for $v_\ell$ when $E[\ell]$ is irreducible (\Cref{ch:moddeg})\\
$\ell\nmid N$ & Steps 1, 3, 4 & Agashe-Ribet-Stein fails at $\ell^2\mid N$ (\curve{99a1}); Fontaine-Laffaille and the integral structures need $\ell\nmid N$; Lemma lem:fpr needs $q\ne\ell$\\
$\ell$ odd & Steps 1, 4 & the factor $\frac12$ in $\hat\delta_!$; the term at $\infty$; self-duality of $\ad T$\\
$E[\ell]$ irreducible & Steps 1, 2 & non-Eisenstein $\mf m$; homothety of lattices, also in \Cref{thm:congruence:dfg14}(c)\\
irreducible over $\Q(\sqrt{\ell^*})$ & Steps 3, 4 & Taylor-Wiles (\Cref{rem:congruence:where-hypothesis}); $h_\ell(E)=0$ and $H^1_f(\Q,V^D)=0$\\
multiplicity one & nowhere & \Cref{rem:congruence:rem-free}\\
\bottomrule
\end{tabular}
\end{center}
\end{remark}

\begin{example}[three curves at $\ell=3$]\label{ex:congruence:main-examples}
Each curve below has $3\nmid N$ and $E[3]$ irreducible with image not $\Ns$ (Cremona's \texttt{galrep} file), so $3\notin S_f$ and \Cref{thm:congruence:main} applies; the Selmer lengths are
therefore \emph{computed}. Data from PARI/GP.
(a) \curve{11a1}: type $\kod{I}_5$ at $11$, $\degphi=1$. Here $e_3(E,11)=0$ (multiplicative), $t_3(E,11)=v_3(5)=0$, so $v_3(\Tam_{11}(\ad T))=0$, and $0=\len H^1_f(\Q,W)$.
(b) \curve{121a1}: type $\kod{II}$ at $11$ (potentially good), $\degphi=6$. PARI's symmetric square Euler factor at $11$ is $(1+11X)^{-1}$, so $L_{11}(A_f,1)^{-1}=1+\frac{11}{121}=\frac{12}{11}$ and
$e_3(E,11)=v_3(12)=1$; $\Tam_{11}(\ad T)=\Z_3$ (Lemma lem:tamadd). So $1=v_3(6)=\len H^1_f(\Q,W)+0+1$ and $H^1_f(\Q,\ad T_3E\otimes\Q_3/\Z_3)=0$.
(c) \curve{539c1}: type $\kod{I}_3^*$ at $7$ ($v_7(j)=-3$) and $\kod{I}_2$ at $11$, $\degphi=288$, $v_3=2$. At $7$: $e_3=v_3(1-7^{-2})=v_3(48)=1$, $t_3=v_3(3)=1$, and the Tate parameter of the
multiplicative twist \curve{77c1} is $q_E=5\cdot7^3+O(7^4)$, so $\bar u=5$, which is not a cube in $\F_7^\times$ (the cubes are $\pm1$): $\min(1,v_3(6),v_3(c(u)))=0$ and $v_3(\Tam_7(\ad T))=1$. At $11$: $t_3=v_3(2)=0$,
$e_3=0$. So $2=\len H^1_f(\Q,W)+1+1$, and the Selmer group vanishes again. Had $\bar u$ been a cube, the identity would have forced $v_3(\degphi)\ge3$.
Further examples with nonzero Selmer groups are in \Cref{ch:blochkato}.
\end{example}

\section[Remark rem:ko: Tamagawa excess; Kim and Ota]{Remark rem:ko: the Tamagawa excess and the formula of Kim and Ota}\label{sec:congruence:rem-ko}

This section treats the paper's Remark rem:ko, which compares \Cref{thm:congruence:main} with two finer statements: the excess in the Tamagawa factor at a multiplicative prime
(part (1), tested in the paper's Section sec:tam), and the formula of Kim and Ota for the change of the congruence ideal when primes are moved into the new part of the
level (part (2), quoted in the paper's Section 1.5).

\begin{remark}[Remark rem:ko(1)]\label{rem:congruence:rem-ko-1}
Let $q\ne\ell$ be a prime with $v_q(j_E)<0$, $n=-v_q(j_E)$, and $u$ the unit part of the Tate parameter. The Tamagawa factor is
$v_\ell(\Tam_q(\ad T))=t_\ell(E,q)+\min\bigl(v_\ell(n),v_\ell(q-1),v_\ell(c(u))\bigr)$, with $t_\ell(E,q)=v_\ell(n)$ (\Cref{ch:blochkato}). The excess is positive exactly when $\ell\mid n$, $q\equiv1\pmod\ell$ and $\bar u$ is an
$\ell$-th power in $\F_q^\times$, and then it is at least $1$ and can be larger. By \eqref{eq:congruence:identity}, for $\ell\notin S_f$ the lower bound for $v_\ell(\degphi)$ rises by exactly that excess:
\[
v_\ell(\degphi)\ \ge\ \sum_{q\mid N}e_\ell(E,q)+\sum_{q:\,v_q(j_E)<0}\Bigl(t_\ell(E,q)+\min\bigl(t_\ell(E,q),v_\ell(q-1),v_\ell(c(u_q))\bigr)\Bigr),
\]
at multiplicative primes and at additive potentially multiplicative primes alike. The paper's ``gives $v_\ell(\Tam_q)\ge t_\ell(E,q)+1$'' is the case of a positive excess; the abstract's
``raises the bound by at least one'' is the same statement. In \Cref{ex:congruence:main-examples}(c) the excess at $7$ vanishes because $\bar u=5$ is not a cube.
\end{remark}

\begin{proposition}[Remark rem:ko(2): Kim and Ota against Theorem thm:main]\label{prop:congruence:kim-ota}
Let $E$ be optimal with conductor $N=N^+N^-$, $N^-$ square-free, $\ell\ge5$, $\ell\nmid N$, $\ell\notin S_f$, $\bar\rho$ ramified at every $q\mid N^-$ with $q\equiv\pm1\pmod\ell$, so that the formula of Kim and Ota
holds (\Cref{ch:modgalrep}):
\[
\ord_\lambda\eta_f(N)=\ord_\lambda\eta_f(N^+,N^-)+\sum_{q\mid N^-}t_\ell(E,q).
\]
\begin{enumerate}[label=(\roman*)]
\item $v_\ell(r_E)\ge\sum_{q\mid N^-}t_\ell(E,q)$, and this is contained in \eqref{eq:congruence:thm}, which needs neither $\ell\ge5$ nor the condition at $q\equiv\pm1$.
\item If $\eta_f(N)$ is either the Hecke congruence ideal $\eta^{\mathrm{Hecke}}_f\Zl$ of \Cref{def:congruence:newform-ideal} or the ideal $\eta^\emptyset_{f,\lambda}$ of \Cref{def:congruence:eta-sigma}, then
\begin{multline*}
\ord_\lambda\eta_f(N^+,N^-)=\len H^1_f(\Q,W)+\sum_{q\mid N}e_\ell(E,q)+\sum_{q\mid N^-}\min\bigl(t_\ell,v_\ell(q-1),v_\ell(c(u_q))\bigr)\\
+\sum_{q\mid N^+,\,v_q(j_E)<0}\bigl(t_\ell+\min(t_\ell,v_\ell(q-1),v_\ell(c(u_q)))\bigr),
\end{multline*}
with $t_\ell=t_\ell(E,q)$: a Selmer-theoretic reading of the congruence ideal of the $N^-$-new quotient.
\end{enumerate}
\end{proposition}

\begin{proof}
(i) is in \Cref{ch:modgalrep}; it is contained in \eqref{eq:congruence:thm} because the primes of $N^-$ are multiplicative and all other terms of \eqref{eq:congruence:thm} are nonnegative. (ii) In either case
$\ord_\lambda\eta_f(N)=v_\ell(\degphi)$: by \Cref{thm:congruence:crit} and Agashe-Ribet-Stein in the first, by \Cref{thm:congruence:lem-hecke} in the second. Subtract the formula of Kim and Ota from
\Cref{cor:congruence:explicit}; the primes of $N^-$ have $v_q(j_E)<0$, and their terms $t_\ell(E,q)$ cancel.
\end{proof}

\begin{remark}[on the identification in (ii)]\label{rem:congruence:kim-ota-identification}
We have not checked in \cite{KimOta} whether their $\eta_f(N)$ is defined as the congruence ideal of the full Hecke algebra (as the paper's Remark rem:ko(2) says: ``concerns the same
congruence ideal'') or as a cohomological ideal; under the hypotheses of the proposition the two have the same valuation (\Cref{rem:congruence:two-ideals}), so (ii) holds in either reading.
The theorem numbers \cite[Theorem 1.1, Definition 2.12]{KimOta} are cited as the paper cites them.
\end{remark}

\section{Twist monotonicity: Theorem thm:twist}\label{sec:congruence:twist}

This section proves the paper's Theorem thm:twist: under a quadratic twist that does not raise the level, the congruence number can only drop. The proof transports
congruences along the twist map $\iota$ and uses nothing beyond \Cref{ch:hecke}; in particular it needs no hypothesis on $\ell$. It is the model for Step 1 of the proof of
Theorem thm:quant in the next section, and its data paragraph is where the paper first sees the Euler terms that Theorem thm:quant explains.

\begin{proposition}[the twist map, recollection]\label{prop:congruence:twist-inputs}
Let $S$ be a finite set of odd primes, $D=\prod_{q\in S}q$, $\chi_S$ and $\iota=\iota_S$ as in \Cref{conv:congruence:setup}(g).
\begin{enumerate}[label=(\roman*)]
\item $\iota$ maps $S_2(\Gamma_0(M))$ into $S_2(\Gamma_0(\lcm(M,D^2)))$ and integral forms to integral forms; $\iota(h)=\iota(h^S)$, and $\iota\circ\iota(h)=h^S$.
\item If $h$ lies in the generalised eigenspace of the operators $T_p$, $p\nmid MD$, for a system $(\lambda_p)$, then $\iota(h)$ lies in the generalised eigenspace for the system $(\chi_S(p)\lambda_p)$.
\item Let $f$ be a newform of level $N$. If $\iota(h)$ has level dividing $N$ and $h$ lies in a sum of generalised eigenspaces whose twisted systems $(\chi_S(p)\lambda_p)$ differ from the system of $f$
at infinitely many $p$, then $\iota(h)\in f^\perp$; in particular, if $g$ is a newform with $\iota(g)=f$, then $\iota$ maps the forms of level dividing $N$ in
$g^\perp$ into $f^\perp$.
\item (Levels of twists.) Let $E$ have conductor $N$, $S$ a set of odd primes with $q^2\mid N$ for $q\in S$, and suppose $E$ has, at every $q\in S$, type $\kod{I}_n^*$ ($n\ge1$), resp.\ $\kod{I}_0^*$, resp.\
$\kod{II},\kod{III},\kod{IV}$ or a dual with $q\ge5$. Then $N_S=N/D$, resp.\ $N/D^2$, resp.\ $N$, and in every case $\iota(f_{E_S})=f_E$.
\end{enumerate}
\end{proposition}

\begin{proof}
(i) and (ii) are \Cref{ch:hecke} (the twist operator); (iii) is \Cref{ch:hecke} (transport along $\iota$), which rests on strong multiplicity one: a generalised eigenform of the $T_p$, $p\nmid N$,
for a system different from that of $f$ at infinitely many $p$ is orthogonal to every $f(dz)$. (iv) is \Cref{ch:modgalrep} (levels of $E_S$), with the twist lemma of \Cref{ch:reduction} for the
types, and the case of the paper's Lemma lem:twist(iii) for $\iota(f_{E_S})=f_E$ (\Cref{ch:hecke}, the case of the paper).
\end{proof}

\begin{theorem}[the paper's Theorem thm:twist]\label{thm:congruence:twist}
Let $E/\Q$ have conductor $N$, let $S$ be a set of odd primes, $D=\prod_{q\in S}q$, $\chi_S=\prod_{q\in S}\chi_{q^*}$, and $E_S=E\otimes\chi_S$ of conductor $N_S$. If $q^2\mid N$ for every $q\in S$ and $\lcm(N_S,D^2)\mid N$,
then $r_{E_S}\mid r_E$. If moreover $\lcm(N,D^2)\mid N_S$, then $r_{E_S}=r_E$.
\end{theorem}

\begin{proof}
Let $f=f_E$, $g=f_{E_S}$, $\iota=\iota_S$.

\emph{(1) $\iota(g)=f$.} By \Cref{prop:congruence:twist-inputs}(i), $\iota(g)\in S_2(\Gamma_0(\lcm(N_S,D^2)))\subseteq S_2(\Gamma_0(N))$. Both $f$ and $\iota(g)$ have multiplicative coefficients ($\chi_S$ is multiplicative).
For a prime $p\nmid D$ the character $\chi_S$ is unramified at $p$ (also at $p=2$, as $D^*=\prod q^*\equiv1\pmod4$), so $V_\ell E\simeq V_\ell E_S\otimes\chi_S$ gives $a_p(f)=\chi_S(p)a_p(g)$ at primes of good
reduction, and the same holds at multiplicative primes (the unramified twist changes split to nonsplit exactly when $\chi_S(p)=-1$) and at additive primes (both coefficients vanish). The
prime-power coefficients at $p\nmid D$ satisfy the same recursion on both sides, $p$ dividing $N$ if and only if it divides $N_S$. For $q\in S$, $a_{q^k}(f)=0$ because $q^2\mid N$ (\Cref{ch:hecke}), and
$a_{q^k}(\iota(g))=\chi_S(q^k)a_{q^k}(g)=0$. So $\iota(g)=f$.

\emph{(2) $\iota$ maps $g^\perp$ into $f^\perp$.} At level $N_S$, $g$ is a newform, so $g^\perp$ is $\T$-stable and is the sum of the generalised eigenspaces of the systems of eigenvalues different from that of
$g$ (\Cref{ch:hecke}); each such system differs from that of $g$ at infinitely many primes (strong multiplicity one). Twisting multiplies the eigenvalues by $\chi_S(p)$, so the twisted systems differ from
that of $f=\iota(g)$ at infinitely many primes, and \Cref{prop:congruence:twist-inputs}(iii) gives $\iota(g^\perp)\subseteq f^\perp$, the images having level dividing $N$.

\emph{(3) Transport.} By definition of $r_{E_S}$ there is $h\in g^\perp\cap S_2(\Gamma_0(N_S);\Z)$ with $g\equiv h\pmod{r_{E_S}}$ coefficientwise. Applying $\iota$, which preserves integrality and congruences
(it multiplies each coefficient by $0$ or $\pm1$), $f=\iota(g)\equiv\iota(h)\pmod{r_{E_S}}$ with $\iota(h)\in f^\perp\cap S_2(\Gamma_0(N);\Z)$. By Lemma lem:crit (\Cref{thm:congruence:crit}) the moduli of such congruences
are exactly the divisors of $r_E$, so $r_{E_S}\mid r_E$. (The definition of $r_E$ as the largest modulus gives only $r_{E_S}\le r_E$; the divisibility needs more than the definition: the
moduli are closed under $\lcm$, which is the cyclicity step of Lemma lem:crit, its elementary half.)

\emph{(4) Symmetry.} $E=(E_S)\otimes\chi_S$, and if $\lcm(N,D^2)\mid N_S$ then $q^2\mid N_S$ for $q\in S$; so the hypotheses hold with the roles of $E$ and $E_S$ exchanged, and $r_E\mid r_{E_S}$.
\end{proof}

\begin{example}[the three cases, and the data]\label{ex:congruence:twist-cases}
By \Cref{prop:congruence:twist-inputs}(iv) the hypotheses of \Cref{thm:congruence:twist} hold in three cases (the paper's list after the theorem).
(a) Type $\kod{I}_n^*$ at every $q\in S$: $N_S=N/D$, $r_{E_S}\mid r_E$. For instance \curve{539c1} ($\kod{I}_3^*$ at $7$) and \curve{77c1}: $r=288$ and $6$; \curve{980i1} ($\kod{I}_5^*$ at $7$) and \curve{140b1}
($\kod{I}_5$ at $7$): $\degphi=2880$ and $60$; \curve{121d1} ($\kod{I}_1^*$ at $11$) and \curve{11a1}: $r=24$ and $1$.
(b) Type $\kod{I}_0^*$ at every $q\in S$: $N_S=N/D^2$, $E_S$ good at $S$. For instance \curve{539d1} and \curve{11a1}: $r=504$ and $1$.
(c) Potentially good types $\kod{II},\kod{III},\kod{IV}$ or duals at $q\ge5$: $N_S=N$ and $r_{E_S}=r_E$. For instance \curve{121a1} (type $\kod{II}$ at $11$) and its twist by $\chi_{-11}$, which is
\curve{121c2} (type $\kod{IV}^*$; the twist exchanges $\kod{II}$ and $\kod{IV}^*$): $r=6=6$. (The curve \curve{121c1}, of type $\kod{IV}$, is isogenous to \curve{121c2} and is the twist of \curve{121a2},
of type $\kod{II}^*$; $r$ depends only on the newform, so it is $6$ for all four curves.)
(Congruence numbers computed with PARI/GP as the index of \Cref{prop:congruence:crit-general}, degrees with \texttt{ellmoddegree}.)
At $\ell=3$ the $3$-adic valuations of the modular degrees in (a) and (b) are $2$ against $1$, $2$ against $1$, $1$ against $0$, and $2$ against $0$: the difference is $e_3(E,q)$, namely
$v_3(7^2-1)=1$, $v_3(48)=1$, $v_3(11^2-1)=1$, and, for the $\kod{I}_0^*$ prime $7$ of \curve{539d1}, $v_3\bigl(\frac67\cdot\frac{64-4}{49}\bigr)=2$. This is Theorem thm:quant.
\emph{The data, recomputed.} For every $X_0(N)$-optimal curve without complex multiplication of conductor below $10{,}000$ (Cremona's \texttt{degphi} file), let $S$ be the set of odd primes of type
$\kod{I}_0^*$ or $\kod{I}_n^*$ (PARI codes $-1$ and $\le-5$), and consider the pairs $(E,\ell)$ with $S\ne\emptyset$, $\ell\in\set{3,5,7,11,13}$, $\ell\nmid N$ and $E[\ell]$ irreducible (no label $\ell$B or $\ell$Cs in Cremona's
\texttt{galrep} file). There are $25{,}477$ such pairs ($15{,}607$ with $3\in S$; $8$ with $\ell=3$ and image $\Ns$), and in every one
$v_\ell(\degphi)=v_\ell(\deg\varphi_{E_S'})+\sum_{q\in S}e_\ell(E,q)$, with $E_S'$ the optimal curve of the class of $E_S$, located by its conductor and its $a_p$ for $p<80$, and $e_\ell(E,q)=v_\ell(q^2-1)$
(type $\kod{I}_n^*$) or $v_\ell(q-1)+v_\ell((q+1)^2-a_q(E_S)^2)$ (type $\kod{I}_0^*$). For case (c), with $S$ the set of primes $q\ge5$ of type $\kod{II},\kod{III},\kod{IV}$ or a dual and the same conditions on $\ell$,
there are $14{,}880$ pairs ($42$ with image $\Ns$ at $\ell=3$), and $v_\ell(\degphi)=v_\ell(\deg\varphi_{E_S'})$ in every one.
\end{example}

\begin{paperlink}
The paragraph after Theorem thm:twist in the revised paper reports, on the working set, the pairs (curve, prime of type $\kod{I}_n^*$) and (curve, prime of potentially good additive type) with
$v_\ell(\degphi)\ge v_\ell(\deg\varphi_{E_q})$ and ``equality in every pair of case (c)'', and then: ``In cases (a) and (b) Theorem thm:quant predicts a difference equal to the Euler term, which is
positive in case (a) at $\ell=3$ for every $q\ne3$, as $3\mid q^2-1$''; and ``Below conductor $10{,}000$ the difference equals the Euler term in all $25{,}477$ pairs''. The count $25{,}477$ agrees with
\Cref{ex:congruence:twist-cases}; the working-set counts are the paper's and were not recomputed here. The version of 10 October 2026 said ``strict inequality only in case (b)'', which
case (a) contradicts (\curve{539c1}, \curve{980i1}, \curve{121d1}).
\end{paperlink}

\section{Theorem thm:quant}\label{sec:congruence:quant}

This section proves the paper's Theorem thm:quant: at the odd primes of type $\kod{I}_0^*$ and $\kod{I}_n^*$ the omitted Euler factor divides the modular degree, under the sole
hypothesis that $E[\ell]$ is irreducible, and with equality (after removing the congruence number of the twist) when $\ell\notin S_f$. The proof raises the level from $N_S$ to $N$ by
depleting the newform of the twist, and moves congruences between levels with $\iota$. We prove, in addition, that Step 1 of the paper's proof is an equality at every prime $\ell\notin S$
(\Cref{lem:congruence:quant-step1}), which the paper's Remark rem:afterquant (2) leaves open.

\begin{theorem}[the paper's Theorem thm:quant]\label{thm:congruence:quant}
Let $E$ be $X_0(N)$-optimal and let $\ell$ be an odd prime with $\ell\nmid N$ and $E[\ell]$ irreducible. Let $S$ be the set of odd primes $q$ at which $E$ has type $\kod{I}_0^*$ or $\kod{I}_n^*$ ($n\ge1$), $\chi_S=\prod_{q\in S}\chi_{q^*}$,
$E_S=E\otimes\chi_S$, and $E_S'$ the optimal curve of the isogeny class of $E_S$. Then
\begin{equation}\label{eq:congruence:quant}
v_\ell(\degphi)\ \ge\ v_\ell(r_{E_S})+\sum_{q\in S}e_\ell(E,q)\ \ge\ \sum_{q\in S}e_\ell(E,q).
\end{equation}
If moreover $\ell\notin S_f$ (automatic for $\ell\ge5$, \Cref{prop:congruence:exceptional-set}), the first inequality is an equality and $v_\ell(r_{E_S})=v_\ell(\deg\varphi_{E_S'})$:
\[
v_\ell(\degphi)=v_\ell(\deg\varphi_{E_S'})+\sum_{q\in S}e_\ell(E,q).
\]
\end{theorem}

\begin{lemma}[the set-up]\label{lem:congruence:quant-setup}
In \Cref{thm:congruence:quant}, let $f=f_E$, $g=f_{E_S}$ (of level $N_S$), $D=\prod_{q\in S}q$ and $\iota=\iota_S$. Then $\ell\notin S$, and:
\begin{enumerate}[label=(\roman*)]
\item $E_S$ has good reduction at the $\kod{I}_0^*$ primes of $S$ and multiplicative reduction at the $\kod{I}_n^*$ primes; so $\delta_q(g)=2$, resp.\ $1$, and $N_S^S=N_S\prod_{q\in S}q^{\delta_q(g)}=N$;
\item $D^2\mid N$, and at $q\notin S$ the curves $E$ and $E_S$ have the same reduction type, the same Euler factors (true and naive) and the same $e_\ell(E,q)$;
\item $g^S\in S_2(\Gamma_0(N))$ is a normalised eigenform of $\T=\T_N$ with $U_qg^S=0$ for $q\in S$, its generalised eigenspace is the line $\C g^S$, and $\iota(g^S)=\iota(g)=f$.
\end{enumerate}
\end{lemma}

\begin{proof}
$\ell\notin S$ because $S$ consists of primes dividing $N$; in particular $\ell\ne3$ if $3\in S$. (i) is the twist lemma of \Cref{ch:reduction} (the paper's Lemma lem:twist, valid at every odd $q$ for
these exchanges: $\kod{I}_0^*\leftrightarrow\kod{I}_0$, $\kod{I}_n^*\leftrightarrow\kod{I}_n$), with the levels of \Cref{prop:congruence:twist-inputs}(iv). (ii) $v_q(N)=2$ for $q\in S$ (types $\kod{I}_0^*$, $\kod{I}_n^*$ at odd $q$ have
conductor exponent $2$, \Cref{ch:reduction}). At $q\notin S$ the character $\chi_S$ is unramified (at $q=2$ too, as $\prod_{q\in S}q^*\equiv1\pmod4$), so the twist does not change the type, $a_q(E_S)^2=a_q(E)^2$,
and the local factors of $\ad V$, which are twist-invariant, agree (\Cref{ch:lfunctions}). (iii) By (i) and \Cref{ch:hecke} (depletion) $g^S$ is an eigenform of $\T_N$ with generalised eigenspace $\C g^S$;
$\iota$ kills the coefficients at multiples of $D$, so $\iota(g^S)=\iota(g)$, which is $f$ by \Cref{prop:congruence:twist-inputs}(iv).
\end{proof}

\begin{lemma}[Step 1, as an equality]\label{lem:congruence:quant-step1}
Let $E/\Q$ be any elliptic curve of conductor $N$ (optimal or not), $S$ the set of odd primes at which $E$ has type $\kod{I}_0^*$ or $\kod{I}_n^*$, and $f$, $g$, $D$, $\iota$ as in
\Cref{lem:congruence:quant-setup}, whose parts (i) and (iii) use neither a prime $\ell$ nor optimality. Let $\pi\colon\T_N\to\Z$ be the eigencharacter of $g^S$ and
$\eta^{\mathrm{Hecke}}(g^S)=\pi(\Ann_{\T_N}\ker\pi)=r_\pi\Z$ (\Cref{prop:congruence:crit-general}). Then $r_E\mid r_\pi$, and for every prime $\ell\notin S$,
\[
v_\ell(r_E)=v_\ell\bigl(\eta^{\mathrm{Hecke}}(g^S)\bigr).
\]
No hypothesis on $\ell$ other than $\ell\notin S$ is needed: neither $\ell$ odd, nor $\ell\nmid N$, nor $E[\ell]$ irreducible, nor $\ell\notin S_f$.
\end{lemma}

\begin{proof}
\emph{$r_E\mid r_\pi$.} This is \Cref{ch:hecke} (``$r_f$ divides the congruence module of $g^S$'', and its corollary for the paper's situation): take $h\in f^\perp\cap S_2(\Gamma_0(N);\Z)$ with $f\equiv h\pmod{r_E}$; then
$g^S=\iota(f)\equiv\iota(h)\pmod{r_E}$, because $\iota\circ\iota(g)=g^S$ and $\iota(f)=\iota(\iota(g))$; and $\iota(f^\perp)$ lies in the $\T$-stable complement of $\C g^S$ (the systems in $f^\perp$ differ from that of $f$ at infinitely many
primes, so their twists differ from that of $g$, by strong multiplicity one). \Cref{prop:congruence:crit-general} gives $r_E\mid r_\pi$. So $v_\ell(r_E)\le v_\ell(r_\pi)$ for every $\ell$.

\emph{$v_\ell(r_E)\ge v_\ell(r_\pi)$ for $\ell\notin S$ (the paper's Step 1).} Let $t=v_\ell(r_\pi)$ and $\mf m$ the maximal ideal of $\T_\ell$ attached to $\pi$. By \Cref{prop:congruence:crit-general} over $\Zl$ there is
$h\in g^c\cap S_2(\Gamma_0(N);\Zl)$, $g^c$ the $\T$-stable complement of $\C g^S$, with $g^S\equiv h\pmod{\ell^t}$. Since $\pi$ factors through $\T_{\mf m}$, $e_{\mf m}g^S=g^S$; the idempotent $e_{\mf m}\in\T_\ell$ preserves
$S_2(\Gamma_0(N);\Zl)$ and $g^c$, so $h'=e_{\mf m}h\in g^c\cap S_2(\Gamma_0(N);\Zl)$ and $g^S-h'=e_{\mf m}(g^S-h)\in\ell^tS_2(\Gamma_0(N);\Zl)$. By \Cref{prop:congruence:hecke-local}(iii), $h'$ lies in the sum of the generalised
eigenspaces $V(\phi)$ with $\phi\ne\pi$ and $\phi\equiv\pi$ modulo the maximal ideal of $\overline\Z_\ell$. For such $\phi$ and $q\in S$, $\phi(U_q)\equiv\pi(U_q)=0$; every eigenvalue of $U_q$ on $S_2(\Gamma_0(N))$ is $0$ or a unit at every
prime not above $q$ (\Cref{ch:hecke}: it is $0$, $\pm1$ or a root of $X^2-a_q(h)X+q$ for a newform $h$), and $\ell\ne q$, so $\phi(U_q)=0$. Each $V(\phi)$ lies in the old space $V_h$ of a newform $h$ of level dividing $N$.
If $h=g$, then $\phi$ agrees with $\pi$ on the $T_p$, $p\nmid N$, and on the $U_q$, $q\in S$; the operators $U_p$ for $p\mid N_S$, $p\notin S$, act on $V_g$ by the scalar $a_p(g)$ (no oldforms at those primes,
$v_p(N)=v_p(N_S)$); so $\phi=\pi$, which is excluded. Hence $h\ne g$. The twist $\iota(V(\phi))$ lies in $S_2(\Gamma_0(N))$, because the level of a twist by $\chi_S$ divides $\lcm(N,D^2)=N$, and in the generalised
eigenspace of the $T_p$, $p\nmid N$, for the system $\chi_S(p)a_p(h)$; this differs from the system $\chi_S(p)a_p(g)$ of $f$ at infinitely many $p$ (otherwise $h=g$ by strong multiplicity one), so
$\iota(V(\phi))\subseteq f^\perp$ (\Cref{prop:congruence:twist-inputs}(iii)). Therefore $\iota(h')\in f^\perp\cap S_2(\Gamma_0(N);\Zl)$ and
\[
f=\iota(g^S)\equiv\iota(h')\pmod{\ell^t}.
\]
By Lemma lem:crit over $\Zl$ (\Cref{thm:congruence:crit}), $t\le v_\ell(r_E)$.
\end{proof}

\begin{remark}[Commentary: Remark rem:afterquant (2) and Open question 6]\label{rem:congruence:step1-commentary}
The paper proves, in its Step 1, the inequality $v_\ell(r_E)\ge v_\ell(\eta^{\mathrm{Hecke}}(g^S))$ under the hypotheses of Theorem thm:quant, and says in Remark rem:afterquant (2): ``When $\ell\notin S_f$
every inequality of the proof is an equality, in particular those of Steps 1 and 2; no multiplicity one at level $N$ or $N_S$ is used. Whether Step 1 is an equality without that hypothesis is
open; the Hecke congruence modules at level $539$ (of orders $288$ and $504$) are consistent with it.'' Its Open question 6 lists ``whether Step 1 of the proof of Theorem thm:quant is an
equality when $\ell\in S_f$''. \Cref{lem:congruence:quant-step1} answers both, and goes beyond the hypotheses of Theorem thm:quant: Step 1 is an equality at \emph{every} prime $\ell\notin S$,
with no hypothesis on $\ell$ at all: $\ell=2$ is allowed, $\ell$ may divide $N$, $E[\ell]$ may be reducible, and $\ell$ may lie in $S_f$. The reason is that the transport along $\iota$ works in both
directions (\Cref{ch:hecke} supplies the direction from $f$ to $g^S$), and the proof of the paper's direction above uses only $\ell\ne q$ for $q\in S$. This supersedes, at $\ell=2$, the
commentary of \Cref{ch:hecke}, which proves the equality under the hypotheses of Theorem thm:quant and records the equality at $\ell=2$ at level $539$ as observed, not proved: it is now
proved. At level $539$, \Cref{ex:congruence:level-539} gives $r(\curve{539c1})=288$ and $r(\curve{539d1})=504$, equal to the two depleted modules as integers; the $2$-parts and $3$-parts are
covered by the lemma, while the $7$-parts ($7\in S$) are not, and there the equality ($7^0$ and $7^1$) is observed, not proved.
\end{remark}

\begin{lemma}[Steps 2 to 4]\label{lem:congruence:quant-steps234}
In the situation of \Cref{lem:congruence:quant-setup}, with $\ell$ odd, $\ell\nmid N$ and $E[\ell]$ irreducible:
\begin{enumerate}[label=(\roman*)]
\item (Step 2) $v_\ell(\eta^{\mathrm{Hecke}}(g^S))\ge v_\ell(\eta^S_{g,\ell})$;
\item (Step 3) $v_\ell(\eta^S_{g,\ell})=v_\ell(\eta^\emptyset_{g,\ell})+\sum_{q\in S}e_\ell(E,q)$;
\item (Step 4) $v_\ell(\eta^\emptyset_{g,\ell})=v_\ell(\deg\varphi_{E_S'})=v_\ell(r_{E_S})$.
\end{enumerate}
\end{lemma}

\begin{proof}
The residual representation of $g$ is $\bar\rho\otimes\chi_S$, irreducible, and $\ell\nmid2N_S$, $N_S$ dividing $N$.

(i) Apply \Cref{lem:congruence:ddt417} with $\mathcal O=\Zl$, $\T$ replaced by $\T_{\mf m}$ ($\mf m$ the maximal ideal of $\pi$), $L=H^1(X_0(N),\Zl)_{\mf m}$, the pairing $\hat\delta_!$ (perfect on the $\mf m$-component, with
$\T_{\mf m}$ self-adjoint, \Cref{prop:congruence:hecke-local}(iv)), and $\pi\colon\T_{\mf m}\onto\Zl$. The congruence ideal of $\pi$ on $\T_{\mf m}$ is $\eta^{\mathrm{Hecke}}(g^S)\Zl$ (\Cref{prop:congruence:congruence-module}(iv)).
The module $L[\ker\pi]$ is $H^1(X_0(N),\Zl)[\ker\pi]$, which is $M_{g^S,\ell}$ of \Cref{def:congruence:eta-sigma} (level $N=N_S^S$; \Cref{prop:congruence:eta-sigma}(i), $\mf m$ being non-Eisenstein as its residual
representation is $\bar\rho\otimes\chi_S$): rank $d=2$, Gram determinant $(\eta^S_{g,\ell})^2$ up to a unit (\Cref{prop:congruence:eta-sigma}(ii)), nonzero by (ii) and (iii) below. So
$2v_\ell(\eta^{\mathrm{Hecke}}(g^S))\ge2v_\ell(\eta^S_{g,\ell})$. Freeness of $L$ over $\T_{\mf m}$ is not needed.

(ii) Apply \Cref{thm:congruence:dfg14}(c) to $g$ (level $N_S$, $\ell\nmid2N_S$, $\ell\notin S$, $\bar\rho_g$ irreducible) with $\Sigma=S$:
$v_\ell(\eta^S_{g,\ell})=v_\ell(\eta^\emptyset_{g,\ell})+\sum_{q\in S}v_\ell(L^{\nv}_q(A_g,1)^{-1})$. At $q\in S$, $\delta_q(g)\ge1$, so $L^{\nv}_q(A_g,1)=L_q(A_g,1)$; $A_g\simeq A_f$, the adjoint being unchanged by twisting
(\Cref{ch:lfunctions}); and $f$ has $\delta_q(f)=0$, so $e_\ell(E,q)=v_\ell(L_q(A_f,1)^{-1})$ by Definition def:terms. The values are those of Lemma lem:euler: $v_\ell(q^2-1)$ at an $\kod{I}_n^*$ prime and
$v_\ell(q-1)+v_\ell((q+1)^2-a_q(E_S)^2)$ at an $\kod{I}_0^*$ prime, $q=3$ included (\Cref{ch:lfunctions}; for $q=3$, $\ell\ne3$).

(iii) \Cref{thm:congruence:lem-hecke} applied to $E_S'$, which is optimal of conductor $N_S$, with $\ell\nmid2N_S$ and $E_S'[\ell]\simeq E_S[\ell]$ irreducible (the isogeny between $E_S$ and $E_S'$ has degree prime to $\ell$,
as $E_S[\ell]$ is irreducible). Its newform is $g$. No freeness is used.
\end{proof}

\begin{proof}[Proof of \Cref{thm:congruence:quant}]
\emph{The inequalities.} As $\ell^2\nmid N$, $v_\ell(\degphi)=v_\ell(r_E)$ \cite[Theorem 2.2]{ARS12}. By \Cref{lem:congruence:quant-step1} (only the inequality $\ge$ is needed here) and
\Cref{lem:congruence:quant-steps234},
\[
v_\ell(\degphi)=v_\ell(r_E)=v_\ell(\eta^{\mathrm{Hecke}}(g^S))\ge v_\ell(\eta^S_{g,\ell})=v_\ell(r_{E_S})+\sum_{q\in S}e_\ell(E,q),
\]
and the second inequality of \eqref{eq:congruence:quant} holds because $r_{E_S}$ is an integer.

\emph{The equality.} Let $\ell\notin S_f$. Then $\ell\notin S_{f_{E_S'}}$: $\ell\nmid2N_S$, and $\bar\rho_{E_S'}\simeq\bar\rho\otimes\chi_S$, whose restriction to $\Gal(\overline\Q/\Q(\sqrt{\ell^*}))$ is absolutely irreducible iff that of $\bar\rho$
is. Apply \Cref{thm:congruence:main} to $E$ (conductor $N$) and to $E_S'$ (conductor $N_S$). The $G_\Q$-modules $\ad T_\ell E$ and $\ad T_\ell E_S'$ are isomorphic (twisting does not change the adjoint, and
$T_\ell E_S\simeq T_\ell E_S'$ by an isogeny of degree prime to $\ell$), so the Selmer groups $H^1_f(\Q,W)$ coincide and, at every prime $q$, so do the Tamagawa factors $\Tam_q(\ad T)$, which are local
invariants of $\ad T$ (at a prime where $\ad T$ is unramified they are trivial, so the sums over $q\mid N$ and over $q\mid N_S$ agree). By \Cref{lem:congruence:quant-setup}(ii) the Euler terms agree at
$q\notin S$; at $q\in S$, $e_\ell(E_S',q)=0$, the reduction of $E_S'$ being good or multiplicative. Subtracting the two identities \eqref{eq:congruence:identity},
\[
v_\ell(\degphi)-v_\ell(\deg\varphi_{E_S'})=\sum_{q\in S}e_\ell(E,q),
\]
and $v_\ell(\deg\varphi_{E_S'})=v_\ell(r_{E_S})$ by \Cref{lem:congruence:quant-steps234}(iii).
\end{proof}

\begin{corollary}[every step is an equality when $\ell\notin S_f$]\label{cor:congruence:quant-equality}
In \Cref{thm:congruence:quant}, if $\ell\notin S_f$ then
\[
v_\ell(r_E)=v_\ell(\eta^{\mathrm{Hecke}}(g^S))=v_\ell(\eta^S_{g,\ell})=v_\ell(r_{E_S})+\sum_{q\in S}e_\ell(E,q).
\]
In particular the Hecke and cohomological congruence ideals of the depleted form $g^S$ at level $N$ have the same valuation, a freeness-type statement obtained without any multiplicity-one
theorem.
\end{corollary}

\begin{proof}
The chain in the proof of \Cref{thm:congruence:quant} begins and ends with $v_\ell(\degphi)=v_\ell(r_E)$ when the first inequality of \eqref{eq:congruence:quant} is an equality, so every inequality in it is an
equality.
\end{proof}

\begin{remark}[the outer inequality, and the word ``elementary'']\label{rem:congruence:outer-inequality}
The outer inequality, $v_\ell(\degphi)$ at least $\sum_{q\in S}e_\ell(E,q)$, needs only Step 1 (the paper's direction), Steps 2 and 3, and the integrality of $\eta^\emptyset_{g,\ell}$; neither Agashe-Ribet-Stein at level $N_S$ nor
any multiplicity-one theorem. It is still not elementary: Step 3 is the proposition of Diamond, Flach and Guo and Step 2 is Lemma lem:ddt417. The revised paper's abstract says ``by a twisting
argument and the level-raising formula of Diamond, Flach and Guo'', which is accurate; the version of 10 October 2026 said ``by an elementary twisting argument''.
\end{remark}

\begin{remark}[Remark rem:afterquant (1) and (3)]\label{rem:congruence:after-quant}
(1) Theorem thm:quant is level raising: from $N_S$ to $N_S^S=N$, at the price of the level-raising Euler factors of \Cref{prop:congruence:three-cases}. For $\ell=3$, a single prime of type $\kod{I}_n^*$ (necessarily
$q\ge5$, since $3\nmid N$) forces $3\mid\degphi$, because $3\mid q^2-1$; so does a prime of type $\kod{I}_0^*$ with $q\equiv1\pmod3$, because then $3\mid q-1$. Neither needs the hypothesis on $\Q(\sqrt{-3})$, and
neither is in Theorem thm:qual. (3) For the potentially good types $\kod{II},\kod{III},\kod{IV}$ and their duals the quadratic twist does not lower the level ($N_S=N$ in \Cref{prop:congruence:twist-inputs}(iv)), $f$ is its own
depletion at $S$, and \Cref{thm:congruence:dfg14} gives no factor; Step 1 has nothing to transport. The omitted factor at such primes is reached only through \Cref{thm:congruence:dfg37}, that is with
$\ell\notin S_f$, and the paper's Section sec:anomaly shows that for $\ell=3$ and image $\Ns$ the bound can fail through a type $\kod{III}$ term. The data of \Cref{ex:congruence:twist-cases}(c) (equality of
$v_\ell$ of the degrees in all $14{,}880$ pairs) are consistent with this: twisting at potentially good primes moves neither side.
\end{remark}

\section[Beyond the hypothesis: rem:ns, Open questions 1, 2]{Beyond the hypothesis: Remark rem:ns and Open questions 1 and 2}\label{sec:congruence:beyond}

This section records what the assembly of Section \ref{sec:congruence:assembly} gives, and what it would need, at the two boundaries of the hypothesis of Theorem thm:main: the
curves with $\ell=3$ and image $\Ns$, where $3\in S_f$ (the paper's Remark rem:ns, Section sec:anomaly and Open question 2), and the primes $\ell\mid N$ (Open question 1).

\begin{remark}[Remark rem:ns: the conditional identity for image $\Ns$]\label{rem:congruence:rem-ns}
Let $\ell=3$, $3\nmid N$, and let the image of $\bar\rho_{E,3}$ be $N_s(\F_3)$. Then $3\in S_f$ and neither \Cref{thm:congruence:dfg37} nor \Cref{thm:congruence:dfg27} applies. Steps 1 and 2 of
\Cref{thm:congruence:main} still hold ($E[3]$ is irreducible). Step 4 holds in a corrected form \emph{if} $H^1_f(\Q,V^D)=0$ (equivalently $H^1_f(\Q,\ad V_3E)=0$, \Cref{ch:blochkato}); the condition
$H^0(\Q,V^D)=0$ is automatic. Then Lemma lem:dfg21 gives, with $h_3(E)=1$ (Lemma lem:h0; \Cref{ch:galrep}),
\[
\len H^1_\Sigma(\Q,W)=\len H^1_f(\Q,W)+\sum_{p\in\Sigma}\bigl(v_3(L_p(A_f,1)^{-1})+v_3(\Tam_p(\ad T))\bigr)-h_3(E)
\]
(\Cref{ch:blochkato} proves this, conditionally). If in addition $v_3(\eta^\Sigma_{f,3})=\len H^1_\Sigma(\Q,W)$ held, as in \Cref{thm:congruence:dfg37}, Steps 1, 2 and 5 would give (eq:identity) with $h_3(E)$
subtracted on the right, which is the shape of Conjecture conj:euler. The paper now states the assumption $H^1_f(\Q,V^D)=0$ explicitly (``we assume it here''); Diamond, Flach and Guo's Theorem 2.7 does
not supply it for these curves, because it is proved only for $\lambda\notin S_f$.
\end{remark}

\begin{remark}[Open question 2]\label{rem:congruence:open-question-2}
For the curves with image $\Ns$ at $\ell=3$ two inputs are missing. (a) On the Hecke side, $v_3(\eta^\Sigma_{f,3})=\len H^1_\Sigma(\Q,W)$ would need the Taylor-Wiles method for a residual representation induced
from $\Q(\sqrt{-3})=\Q(\mu_3)$, where the auxiliary primes cannot kill the dual Selmer group (\Cref{rem:congruence:where-hypothesis}(c)). (b) On the Selmer side, $H^1_f(\Q,\ad V_3E)=0$. The paper's data
(failures of (eq:thm) by exactly one, never more, for these curves) and the equality in all $8$ pairs with image $\Ns$ in \Cref{ex:congruence:twist-cases}, which lie outside
\Cref{thm:congruence:quant}'s equality clause, are consistent with the corrected identity. We do not survey later work on residually dihedral representations, which we have not checked
against these questions.
\end{remark}

\begin{remark}[Open question 1: primes $\ell$ dividing $N$]\label{rem:congruence:open-question-1}
At $\ell\mid N$ three steps of \Cref{thm:congruence:main} fail. Step 1: Agashe-Ribet-Stein can fail when $\ell^2\mid N$, and $r_E\ne\degphi$ occurs (\curve{99a1}: $r=12$, $\degphi=4$ at $\ell=3$; \curve{539d1}: $r=504$,
$\degphi=72$ at $\ell=7$; \Cref{ch:moddeg}). Step 3: $\ell\mid N$ puts $\ell$ in $S_f$; the integral structures and the Fontaine-Laffaille condition at $\ell$ need $\ell\nmid N$. Step 5: the Euler factor at $q=\ell$ is not
the relevant local quantity (the paper's Remark rem:wild). A replacement would be a deformation condition at $\ell$ (finite flat, ordinary or semistable) whose local index plays the role of
$e_\ell(E,\ell)$; the paper's empirical weights $w_3$ at $q=\ell=3$ are the target. We record the question without an answer.
\end{remark}

\section{The imported statements}\label{sec:congruence:imported}

This section lists, for reference, every statement of \cite{DFG04} that the paper cites, what it says, where the paper uses it, and what this book has checked. It is the place that
\Cref{ch:blochkato,ch:lfunctions} point to for the status of these citations.

\begin{remark}[the citations of Diamond, Flach and Guo]\label{rem:congruence:imported-table}
Every entry was compared with the published text \cite{DFG04} for this chapter; the column ``where'' gives the page of the journal, and the content column what the check found.
\begin{center}\small
\begin{tabular}{p{0.19\textwidth}p{0.36\textwidth}p{0.2\textwidth}p{0.13\textwidth}}
\toprule
citation & content (book's notation) & paper's use & where\\
\midrule
Theorem 0.1 (= Theorem 2.15) & $\lambda$-part of Bloch-Kato for $A_f$, $A_f(1)$, $\lambda\notin S_f$ & Section 1.2 & pp.~666, 709\\
Theorem 0.2 (= Theorem 3.7) & $\len H^1_\Sigma=v_\lambda(\eta^\Sigma_f)$, $\Sigma\supseteq\set{p\mid N}$, $\lambda\notin S_f$, $\lambda\nmid\Sigma$ & thm:dfg37 & p.~666\\
definition of $S_f$ (Section 0.2 and display (31)) & $\lambda\mid Nk!$ or residually not absolutely irreducible over $\Q(\sqrt{\ell^*})$ & Section 4.2 & pp.~666, 695\\
Theorem 3.7 & as Theorem 0.2, with ``minimally ramified outside $\Sigma$'', $\ell\notin\Sigma$ & thm:dfg37 & p.~723\\
Proposition 1.4(c) & $\eta^\Sigma=\eta^\emptyset\prod L^{\nv}_p(A_f,1)^{-1}$ if $\bar{\mathcal M}_{f,\lambda}$ is irreducible; proved by homothety and $\gamma_{\mathrm{dR}}(f)=f^\Sigma$ & thm:dfg14 & p.~687\\
Theorem 2.7 & $H^0=H^1_f=0$ for $A_f$ and $A_f(1)$ & thm:dfg27 & p.~699\\
Lemma 2.1 & the exact sequence of Lemma lem:dfg21, for $\Sigma$ nonempty & lem:dfg21 & p.~696\\
Lemma 2.12 & $(f,f)_{\Gamma_1(N)}$ as a naive adjoint $L$-value; $N/(8\pi^3)$ on $\Gamma_0(N)$ for $k=2$ & Section 1 & p.~706\\
(57), p.~708 & $\Tam_p(M)=\Tam_p(M^D)$; definition of $\Tam_p$ & lem:fpr, Section 1.2 & p.~708\\
Sections 1.5.3, 1.6.2, 1.7.2, 1.7.3; (13), (14) & $\hat\delta_!$ through $[UwU]_\omega$, $M_f=M(N,\psi)_![I_f]$, $\delta_p$, $L^{\nv}_p$, $\eta^\Sigma_f$ & Section 4.1 & pp.~682 to 687\\
Sections 1.6.1, 1.8.2 & Hecke action compatible with $\hat\delta_!$ (Proposition 1.3); self-adjoint for the pairing (25) & Section 4.1 & pp.~682, 689\\
Section 2.1 & $H^1_f$ and $H^1_\Sigma$ & Section 4.2 & p.~695\\
Lemma 1.8, Corollary 1.9 & $H^1_c\to H^1$ an isomorphism after localising at a non-Eisenstein $\mf m$ & Section 4.1 & pp.~690, 691\\
\bottomrule
\end{tabular}
\end{center}
Two differences of wording matter. Lemma 2.1 of \cite{DFG04} asks that $\Sigma$ contain every prime at which $H^1_f(\Q_p,W)\ne H^1_{\ur}(\Q_p,W)$; the paper's ``every prime $p\ne\ell$ at which $W$ is ramified''
implies it, since the two groups agree at unramified $p\ne\ell$. And $H^1_\Sigma$ of \cite{DFG04} uses the unramified condition outside $\Sigma\cup\set\ell$, which is the Bloch-Kato condition when
$\Sigma\supseteq\set{p\mid N}$ (Section \ref{sec:congruence:dfg-theorems}). A longer version of \cite{DFG04} is now available \cite{DFG04long}, dated ``August 2001, with minor update December 2025''; it
numbers differently: according to a reading of arXiv:2512.02348 not re-checked for this chapter, its Theorem 0.1 is labelled ``= Theorem 8.10'', its Theorem 0.2 ``= Theorem 7.15'', and its
exceptional set is called $T_f$. The book, like the paper, cites the published numbering.
\end{remark}

\section{Exercises}\label{sec:congruence:exercises}

\begin{exercise}\label{ex:congruence:ex-three-forms}
Let $1\le r\le s$ and $\T=\set{(a,b,c)\in\mathcal O^3:a\equiv b\pmod{\varpi^r},\ a\equiv c\pmod{\varpi^s}}$, with $\pi(a,b,c)=a$. Compute $\wp$, $I$, $\eta(\pi)$ and $\mathrm{Cot}(\pi)$, and show that $\len\mathrm{Cot}(\pi)>v(\eta)$, so that
$\T$ is not a complete intersection. Do the same for the projections to $b$ and to $c$. (Answer: $\wp=0\times\varpi^r\mathcal O\times\varpi^s\mathcal O$, $I=\varpi^s\mathcal O\times0\times0$, $\eta=\varpi^s\mathcal O$,
$\mathrm{Cot}\simeq\mathcal O/\varpi^r\oplus\mathcal O/\varpi^s$, of length $r+s>s$. For the projection to $b$: $\wp$ has basis $(\varpi^r,0,\varpi^r)$, $(0,0,\varpi^s)$, $I=0\times\varpi^r\mathcal O\times0$ and
$\eta=\varpi^r\mathcal O$; the products of the basis elements span $\varpi^r\wp$, so $\mathrm{Cot}=\wp/\varpi^r\wp\simeq(\mathcal O/\varpi^r)^2$, of length $2r>r$. For the projection to $c$: $\wp=\varpi^s\mathcal O\times\varpi^r\mathcal O\times0$,
$I=0\times0\times\varpi^s\mathcal O$, $\eta=\varpi^s\mathcal O$ and $\mathrm{Cot}\simeq\mathcal O/\varpi^s\oplus\mathcal O/\varpi^r$, of length $r+s>s$. The inequality is strict for every projection, as it must be by
\Cref{thm:congruence:numerical-criterion}(ii): $\T$ is local and not a complete intersection, so no surjection to $\mathcal O$ gives equality.)
\end{exercise}

\begin{exercise}\label{ex:congruence:ex-wp-cap-I}
Let $\T=\mathcal O[x]/(x^2(x-\varpi))$ and $\pi(x)=0$. Show that $\Ann_\T(\wp)=x(x-\varpi)\T\subseteq\wp$, so $\eta=0$ and $\wp\cap I\ne0$, and that the factor of $\T_K$ at $\pi$ is $K[x]/(x^2)$. Then let $\pi'(x)=\varpi$ and
show $\eta(\pi')=\varpi^2\mathcal O$. Why does Lemma lem:ddt417 survive the first case? (Hint: \Cref{rem:congruence:ddt417-hypotheses}(a). For $\pi'$: $I=x^2\T$ and $\pi'(x^2)=\varpi^2$.)
\end{exercise}

\begin{exercise}\label{ex:congruence:ex-ddt-diag}
Let $\T_3^\flat$ act diagonally on $L=\mathcal O^3$ with the standard pairing. Run the proof of \Cref{lem:congruence:ddt417}: find $L[\wp]$, $L[I]$, $Q$, and compare $v(\Delta)$ with $v(\eta)$. Is $L$ free over $\T_3^\flat$? (Answer: $L[\wp]=\mathcal Oe_1$, $L[I]=\mathcal Oe_2\oplus\mathcal Oe_3$, $Q=0$,
$v(\Delta)=0<1=v(\eta)$; $L$ is not free, since by Step 6 a free module would give $v(\Delta)=d\,v(\eta)$.)
\end{exercise}

\begin{exercise}\label{ex:congruence:ex-alternating}
(a) Prove that for an alternating form on a free module of rank two over $\mathcal O$ the image of $\wedge^2$ is the ideal of values and the Gram determinant is its square (\Cref{prop:congruence:eta-sigma}(ii)).
(b) Show that in rank four the ideal of values does not determine the Gram determinant: compare the block-diagonal forms with blocks $\sm0a{-a}0$, $\sm0b{-b}0$ for $(a,b)=(1,1)$ and $(1,\varpi)$.
(Answer to (b): both have ideal of values $\mathcal O$; the determinants are $1$ and $\varpi^2$. So for a lattice of rank $d>2$ the Gram determinant,
not the ideal of values, is the right invariant; Diamond, Flach and Guo only need rank two.)
\end{exercise}

\begin{exercise}\label{ex:congruence:ex-idempotent}
Let $\T$ be a commutative $\Zl$-algebra of operators on a lattice $L$ with a perfect pairing for which every $t\in\T$ is self-adjoint. Show that every idempotent $e\in\T$ is self-adjoint, that $eL\perp(1-e)L$,
and that the pairing restricted to $eL$ is perfect. Deduce that for $\mf m=\mf m_{f,\ell}$ the lattice $M_{f,\lambda}$ lies in $e_{\mf m}H^1(X_0(N),\Zl)$ and that $\hat\delta_!$ is perfect there. (Hint: $\langle ex,(1-e)y\rangle=\langle x,e(1-e)y\rangle=0$; a
perfect pairing on $eL\oplus(1-e)L$ with the summands orthogonal is perfect on each.)
\end{exercise}

\begin{exercise}\label{ex:congruence:ex-level-raising}
Let $f=f_{\text{11a}}$ and $\ell=3$. For the primes $p<40$, $p\ne3,11$, compute $v_3$ of Wiles's factor $\frac{(p-1)((p+1)^2-a_p^2)}{p^3}$ and say which part, $p-1$ or $(p+1)^2-a_p^2$, contributes. Which primes satisfy Ribet's
level-raising condition $a_p\equiv\pm(p+1)\pmod3$? (Answer, with $a_p=-2,1,-2,4,-2,0,-1,0,7,3$ for $p=2,5,7,13,17,19,23,29,31,37$: the valuation is positive for $p=7$ ($1+1$), $13$ ($1+2$), $19$ ($2+0$),
$29$ ($0+2$), $31$ ($1+1$), $37$ ($2+0$); the level-raising condition holds for $p=7,13,29,31$; at $p=19,37$ only $p\equiv1\pmod 3$ contributes, which corresponds to congruences at level $11p^2$.)
\end{exercise}

\begin{exercise}\label{ex:congruence:ex-factor-11a}
For $f=f_{\text{11a}}$ compute $L^{\nv}_2(A_f,1)^{-1}$ and $L^{\nv}_7(A_f,1)^{-1}$, and deduce $v_3(\eta^{\{2\}}_{f,3})$ and $v_3(\eta^{\{7\}}_{f,3})$ from \Cref{thm:congruence:dfg14}. Compare the second with the Euler factor
of the symmetric square of \curve{539d1} at $7$, which PARI gives as $(1+3X-21X^2-343X^3)^{-1}$, and explain the coincidence. (Answer: $\frac12\cdot\frac{9-4}{4}=\frac58$, a $3$-adic unit, so
$v_3(\eta^{\{2\}})=0$; $\frac67\cdot\frac{64-4}{49}=\frac{360}{343}$, so $v_3(\eta^{\{7\}})=0+2=2$. And $1+\frac3{49}-\frac{21}{49^2}-\frac{343}{49^3}=\frac{360}{343}$: \curve{539d1} is the twist of \curve{11a1} by $\chi_{-7}$, the true adjoint
factor is twist-invariant, and \curve{11a1} is good at $7$, so its true factor is Wiles's factor. This is Step 3 of the proof of Theorem thm:quant, and $v_3(\eta^{\{7\}})=2=v_3(504)$.)
\end{exercise}

\begin{exercise}\label{ex:congruence:ex-sf}
Let $E$ have good reduction at $\ell\ge5$ and $E[\ell]$ irreducible. Using the description of $\bar\rho|_{I_\ell}$ by fundamental characters (ordinary: $\omega\oplus1$ up to semisimplification; supersingular:
$\omega_2\oplus\omega_2^\ell$), show that $\bar\rho\not\simeq\bar\rho\otimes\chi_{\ell^*}$, so $\ell\notin S_f$. Where does the argument fail for $\ell=3$? (Hint: $\chi_{\ell^*}|_{I_\ell}=\omega^{(\ell-1)/2}$; in the ordinary case
$\set{\omega,1}=\set{\omega^{(\ell+1)/2},\omega^{(\ell-1)/2}}$ forces $(\ell-1)\mid\frac{\ell+1}2$. For $\ell=3$ it holds: $\omega^2=1$ and $\omega^{1}=\omega$.)
\end{exercise}

\begin{exercise}\label{ex:congruence:ex-hur}
(a) Let $p\ne\ell$ and $W=T\otimes\Qlmod$ with $T$ unramified at $p$. Show that $H^1_{\ur}(\Q_p,W)=W/(\Frob_p-1)W$ is divisible and equals $H^1_f(\Q_p,W)$. (b) Let $T=\Zl^2$ with inertia acting through the
tame quotient by $t\mapsto\sm1{\ell^a}01$ and a Frobenius acting trivially on $T^{I_p}$ (possible when $p\equiv1\pmod{\ell^a}$). Show that $W^{I_p}\simeq\Qlmod\oplus\Z/\ell^a$ and that $H^1_{\ur}(\Q_p,W)$ is not divisible,
while $H^1_f(\Q_p,W)$, its maximal divisible subgroup (\Cref{ch:blochkato}), is $\Qlmod$. This is why Diamond, Flach and Guo write the unramified condition only outside a set $\Sigma$ that contains the
ramified primes. (Hint for (b): $(x,y)$ is fixed by $t$ iff $\ell^ay=0$ in $\Qlmod$.)
\end{exercise}

\begin{exercise}\label{ex:congruence:ex-77c}
Write out \Cref{thm:congruence:main} for \curve{77c1} at $\ell=3$: $N=77$, type $\kod{I}_3$ at $7$ and $\kod{I}_2$ at $11$, $\degphi=6$, $E[3]$ irreducible with image not $\Ns$, and Tate parameter at $7$ equal to
$5\cdot7^3+O(7^4)$. Compute every term of \Cref{cor:congruence:explicit} and deduce $\len H^1_f(\Q,\ad T_3E\otimes\Q_3/\Z_3)$. (Answer: $e_3=0$ at both primes (semistable); at $7$, $t_3=1$ and $\bar u=5$ is not a cube in $\F_7^\times$,
so the Tamagawa term is $1$; at $11$, $t_3=v_3(2)=0$. So $1=v_3(6)=\len H^1_f+1$ and the Selmer group is $0$. Compare \Cref{ex:congruence:main-examples}(c): \curve{539c1} is the twist of \curve{77c1} by
$\chi_{-7}$, with the same Selmer group and Tamagawa factors and one more Euler term.)
\end{exercise}

\begin{exercise}\label{ex:congruence:ex-ko}
In the situation of \Cref{rem:congruence:rem-ko-1} with $\ell=3$, $n=9$ and $q\equiv1\pmod9$, show that the excess $\min(v_3(9),v_3(q-1),v_3(c(u)))$ is $2$ if $\bar u$ is a ninth power in $\F_q^\times$, $1$ if it is a cube but
not a ninth power, and $0$ otherwise; so the bound of (eq:thm) can be exceeded by $2$ at one prime. (Hint: $c(u)$ is the class of $\bar u$ in $\F_q^\times/(\F_q^\times)^{9}$, a cyclic group of order $9$, and $v_3$ of a class is
read as in Appendix B, B-C10.)
\end{exercise}

\begin{exercise}\label{ex:congruence:ex-121}
Show that $\iota_{\{11\}}(f_{\text{121c}})=f_{\text{121a}}$ (both sides lie in $S_2(\Gamma_0(121))$, so it suffices to compare $a_n$ for $n\le22$, Sturm's bound; for instance $a_2,a_3,a_5,a_7,a_{13}$ are $-1,2,1,2,-1$ for \curve{121a1} and $1,2,1,-2,1$ for \curve{121c1}), and deduce $r(\curve{121a1})=r(\curve{121c1})$ from \Cref{thm:congruence:twist}. Why does the same twist not lower the level? (Answer: $\chi_{-11}(p)$ is $-1,+1,+1,-1,-1$ for $p=2,3,5,7,13$; the twist by $\chi_{-11}$ exchanges the types $\kod{II}$
and $\kod{IV}^*$ at $11$ (the twist of \curve{121a1}, type $\kod{II}$, is \curve{121c2}, type $\kod{IV}^*$; \curve{121c1}, type $\kod{IV}$, is only isogenous to it, being the twist of \curve{121a2}, type $\kod{II}^*$), all of
conductor exponent $2$, so $N_S=N=121$ and both hypotheses of the theorem hold. The congruence number depends only on the newform, hence only on the isogeny class: $r=6$ for both classes.)
\end{exercise}

\begin{exercise}\label{ex:congruence:ex-q3}
(Theorem thm:quant at $q=3$.) The curves \curve{333d1} and \curve{477a1} have type $\kod{I}_0^*$ at $3$ and $\degphi=28$. PARI gives the Euler factor of the symmetric square at $3$ as $(1-6X+18X^2-27X^3)^{-1}$ for both.
Compute $L_3(A_f,1)^{-1}$, compare with the formula $\frac23\cdot\frac{16-a_3(E')^2}9$ for a twist $E'$ with good reduction at $3$, find $a_3(E')^2$, and check the inequality of \Cref{thm:congruence:quant} at $\ell=7$.
(Answer: $1-\frac69+\frac{18}{81}-\frac{27}{729}=\frac{14}{27}$, so $a_3(E')^2=9$; $e_7(E,3)=v_7(\frac{14}{27})=1\le v_7(28)=1$. As $7\notin S_f$, the equality clause gives $v_7(\deg\varphi_{E_S'})=0$.)
\end{exercise}

\begin{exercise}\label{ex:congruence:ex-chain}
For \curve{539d1} and $\ell=3$ compute each term of the chain of \Cref{cor:congruence:quant-equality}: $v_3(r_E)$, $v_3(\eta^{\mathrm{Hecke}}(g^S))$ with $g=f_{\text{11a}}$, $v_3(\eta^{\{7\}}_{g,3})$, $v_3(r_{E_S})$ and $e_3(E,7)$, using
\Cref{ex:congruence:level-539} and \Cref{ex:congruence:ex-factor-11a}. Then do the same at $\ell=2$ for Step 1 only, which \Cref{lem:congruence:quant-step1} covers although $\ell=2$ is outside the hypotheses of Theorem thm:quant.
(Answer: $2=2=2=0+2$. At $\ell=2$: $v_2(504)=3=v_2(504)$.)
\end{exercise}

\begin{exercise}\label{ex:congruence:ex-step1-prime}
Explain where the proof of the paper's inequality (the second half of the proof of \Cref{lem:congruence:quant-step1}) breaks down for $\ell=q\in S$, when $E$ has type $\kod{I}_0^*$ at $q$ and $E_S$ is ordinary at $q$
($q\nmid a_q(g)$). (Answer: the roots $\alpha,\beta$ of $X^2-a_q(g)X+q$ satisfy $\alpha\beta=q$ with $\alpha$ a $q$-adic unit, so $\beta$ is not a unit; the stabilisation $g(z)-\alpha g(qz)$, with $U_q$-eigenvalue $\beta\equiv0$, lies in the same
maximal ideal as $g^S$ at $\ell=q$. It belongs to the old space of $g$ itself, and $\iota$ sends it to a form with the system of $f$, which is not in $f^\perp$: the transport fails. Congruences between
$g^S$ and this stabilisation contribute to $\eta^{\mathrm{Hecke}}(g^S)$ but not to $r_E$.)
\end{exercise}

\section{Notes and references}\label{sec:congruence:notes}

\emph{Congruence modules.} Congruence primes go back to Hida \cite{Hida81}, who related them to the Petersson norm and the special value of the symmetric square; \cite{HidaBook} is an
exposition, and \cite{Ribet83} treats weight $2$ through modular curves. The algebraic formalism of \Cref{def:congruence:congruence-module} is that of \cite[Section 4]{DDT}; Lenstra's
paper \cite{Lenstra95} proves the numerical criterion in the form of \Cref{thm:congruence:numerical-criterion} and its relation to the Gorenstein property. Lemma lem:ddt417 is the inequality half of
\cite[Lemma 4.17]{DDT}, as the paper cites it; we have not compared the numbering with the published text of \cite{DDT}, and the book proves only the part the paper states. Wiles's
numerical criterion is in \cite[Appendix]{Wiles95} and \cite[Section 5]{DDT}; the theorem number sometimes quoted for it in \cite{DDT} is not verified here.

\emph{Diamond, Flach and Guo.} The paper cites the published version \cite{DFG04}; so does the book. Every item of it that the paper or the book cites (the abstract, the definition of $S_f$ in
Section 0.2 and as display (31), Theorems 0.1, 0.2, 2.7, 2.15, 3.7, Propositions 1.4, 2.2 and 2.16, Lemmas 1.8, 2.1, 2.5 and 2.12, Corollary 1.9, the displays (13), (14), (15), (57), the definition
of minimal ramification and the sections listed in \Cref{rem:congruence:imported-table}) was checked against the published text for this chapter, with the page numbers given there. The longer version \cite{DFG04long} (``August 2001, with minor update December
2025'') contains the foundational details omitted from the published paper, with different numbering. The Taylor-Wiles method is \cite{Wiles95,TaylorWiles}; \cite{Diamond96} gives the freeness
of the cohomology and the version without the Gorenstein hypothesis; \cite{DDT} (Chapters 2 to 5) is the standard account, and we followed its organisation in
Section \ref{sec:congruence:taylor-wiles}. Level raising and Ihara's lemma are \cite{Ribet84,Ribet90} and \cite{DiamondTaylor94}; Fontaine-Laffaille theory is \cite{FontaineLaffaille}, and the
local formula of Fontaine and Perrin-Riou is \cite[I.4.2.2]{FPR94}, as recalled in \cite{DFG04}.

\emph{The modular degree.} Agashe, Ribet and Stein \cite[Theorem 2.2]{ARS12} (theorem number as cited by the paper and checked in \Cref{ch:moddeg}); the multiplicity-one and Gorenstein results
behind Remark rem:free are \cite{Mazur77}, \cite[Theorem 3.5]{RibetStein} (as cited by the paper) and \cite{Tilouine97}; the Shimura subgroup is \cite{LingOesterle,Ribet88}. The formula of Kim and Ota
is \cite[Theorem 1.1]{KimOta}, and Pollack and Weston's earlier result is \cite[Theorem 6.11]{PollackWeston11} (numbering of the arXiv version; hypotheses in \Cref{ch:modgalrep}); Ribet and Takahashi's comparison with Shimura curves is \cite[Theorems 1 and 2]{RibetTakahashi97}, sharpened in \cite[Theorem 3.2, Corollary 3.5]{Takahashi01}
(checked against the sources in \Cref{ch:moddeg}).

\emph{What is new here.} Relative to the paper: \Cref{lem:congruence:quant-step1} (Step 1 of the proof of Theorem thm:quant is an equality at every prime outside $S$, which answers the question
left open in Remark rem:afterquant (2) and Open question 6, the divisibility being due to \Cref{ch:hecke}); the observation that Step 1 uses no hypothesis on $\ell$ beyond $\ell\notin S$;
\Cref{cor:congruence:quant-equality}; the length form of Lemma lem:ddt417 over any discrete valuation ring; and the recomputation of the data of the paragraph after Theorem thm:twist
(\Cref{ex:congruence:twist-cases}). All numbers in this chapter were computed with PARI/GP \cite{PARI} (version 2.17, through the Python interface \texttt{cypari2}) and Cremona's files
\cite{ecdata} for conductors below $10{,}000$; how to reproduce such computations is explained in \Cref{ch:computation}, and the notation is collected in \Cref{ch:notation}. The paper's own
account of these results, with the book's commentary, is in \Cref{ch:paper}.
